\documentclass[10pt]{article}
\usepackage[T1,T2A]{fontenc}
\usepackage[utf8]{inputenc}
\usepackage[english,ukrainian]{babel}
\usepackage{geometry}
\usepackage{longtable}
\usepackage{array}
\usepackage{multirow}
\usepackage{hyperref}
\usepackage{mathptmx}
\renewcommand{\familydefault}{\rmdefault}

\geometry{a4paper,left=2cm,right=2cm,top=2cm,bottom=2cm}

\begin{document}

\begin{flushright}
ЗАТВЕРДЖЕНО\\
Наказом ТОВ "Криптографічні Телесистеми"\\
від \_\_ \_\_\_\_\_\_\_\_\_ 2026 р. № \_\_\_\\
(в редакції наказу ТОВ "Криптографічні Телесистеми"\\
від \_\_ \_\_\_\_\_\_\_\_\_ 2026 р. № \_\_)
\end{flushright}

\vspace{2cm}

\begin{center}
\textbf{\Large Система електронного урядування}\\
\textbf{\Large «ERP/1»}\\
\vspace{1cm}
\textbf{\Large ЦІЛЬОВИЙ ПРОФІЛЬ БЕЗПЕКИ}\\
\vspace{0.5cm}
\large UA.46207985.СЗІ.ЦПБ-67
\end{center}

\newpage

\footnotesize
\begin{longtable}{|>{\centering\arraybackslash}p{0.5cm}|>{\raggedright\arraybackslash}p{4cm}|>{\raggedright\arraybackslash}p{3cm}|>{\raggedright\arraybackslash}p{2cm}|>{\raggedright\arraybackslash}p{5cm}|}
\hline
\textbf{№} & \textbf{Вимога з безпеки інформації} & \textbf{Вимога ГПБ} & \multicolumn{2}{c|}{\textbf{ЦПБ}} \\
\cline{4-5}
& & & \textbf{Захід захисту} & \textbf{Налаштований зміст заходу захисту}\\
\hline
\endfirsthead

\hline
\textbf{№} & \textbf{Вимога з безпеки інформації} & \textbf{Вимога ГПБ} & \multicolumn{2}{c|}{\textbf{ЦПБ}} \\
\cline{4-5}
& & & \textbf{Захід захисту} & \textbf{Налаштований зміст заходу захисту}\\
\hline
\endhead

\hline
\endfoot

\hline
\endlastfoot

\multicolumn{5}{|l|}{\textbf{1. УПРАВЛІННЯ ДОСТУПОМ (AC) (AC)}} \\ \hline
1 & ПОЛІТИКА ТА ПРОЦЕДУРИ УПРАВЛІННЯ ДОСТУПОМ (AC-1) & Параметр: Визначено персонал або ролі, на які поширюється політика контролю доступу\newline Тип: list\newline Значення: \textbf{admin, security\_\allowbreak{}officer}\newline\newline Параметр: Визначено персонал або ролі, на які поширюються процедури контролю доступу\newline Тип: list\newline Значення: \textbf{admin, security\_\allowbreak{}officer}\newline\newline Параметр: Визначено посадову особу, яка керуватиме політикою та процедурами контролю доступу\newline Тип: list\newline Значення: \textbf{admin, security\_\allowbreak{}officer}\newline\newline Параметр: Визначено частоту, з якою переглядається та оновлюється поточна політика контролю доступу\newline Тип: list\newline Значення: \textbf{default\_\allowbreak{}deny\_\allowbreak{}rule, abac\_\allowbreak{}rule\_\allowbreak{}1}\newline\newline Параметр: Визначено події, які вимагають перегляду та оновлення поточної політики контролю доступу\newline Тип: list\newline Значення: \textbf{default\_\allowbreak{}deny\_\allowbreak{}rule, abac\_\allowbreak{}rule\_\allowbreak{}1}\newline\newline Параметр: Визначено частоту, з якою переглядаються та оновлюються поточні процедури контролю доступу\newline Тип: integer\newline Значення: \textbf{30} & AC-1 & a. Розробити, задокументувати та поширити [Призначення: серед визначеного організацією персоналу або ролей]:\\ 
1. 2. [Вибір (один або декілька): Рівень організації; Рівень місії/бізнес-процесу; рівень системи] політики контролю доступу, яка:\\ 
(a) містить мету, сферу застосування, ролі, відповідальність, зобов'язання керівництва, координацію між організаційними підрозділами та систему контролю відповідності (compliances);\\ 
(b) відповідає чинному законодавству, нормативним документам, директивам, нормам, політикам, стандартам і керівним документам. Процедури, що сприяють реалізації політики управління доступом і відповідних заходів управління доступом.\\ 
b. Призначити на посаду [Призначення: визначену організацією посадову особу] для управління, документування і розповсюдження політики та процедур контролю доступом.\\ 
c. Переглянути та оновити:\\ 
1. поточну політику управління доступом [Призначення: з визначеною організацією частотою] та [Призначення: події, визначені організацією];\\ 
2. поточні процедури управління доступом [Призначення: з визначеною організацією частотою] та [Завдання: події, визначені організацією]. \\ \hline
2 & УПРАВЛІННЯ ОБЛІКОВИМИ ЗАПИСАМИ (AC-2) & Параметр: Визначено передумови та критерії членства в групах і ролях; визначено атрибути (за необхідності) для кожного облікового запису; визначено персонал або ролі, необхідні для затвердження запитів на створення облікових записів; визначено політику, процедури, передумови та критерії створення, активації, зміни, деактивації та видалення облікових записів; визначено персонал або ролі, які мають бути повідомлені; визначено період часу, протягом якого адміністратори облікових записів повинні бути повідомлені про те, що облікові записи більше не потрібні; визначено термін, протягом якого необхідно повідомляти адміністраторів облікових записів про звільнення або переведення користувачів; визначено період часу, протягом якого необхідно повідомляти адміністраторів облікових записів про зміни у використанні системи або необхідність знати про зміни для окремої особи; визначено атрибути, необхідні для авторизації доступу до системи (за потреби); AC-02\_\allowbreak{}ODP[10] AC-02a.[01] AC-02a.[02] AC-02b AC-02с AC-02d.01 AC-02d.02 AC-02d.03[01] AC-02d.03[02] AC-02e AC-02f.[01] AC-02f.[02] AC-02f.[03] AC-02f.[04] AC-02f.[05] AC-02g AC-02h.01 AC-02h.02 AC-02h.03 AC-02i.01 AC-02i.02 AC-02i.03 визначено періодичність перегляду облікових записів; визначено та задокументовано типи облікових записів, дозволених для використання в системі; визначено та задокументовано типи облікових записів, які заборонено використовувати в системі; призначені менеджери облікових записів; необхідні умови та критерії для членства в групах та ролях; визначено авторизованих користувачів системи; вказано приналежність до групи або ролі; для кожного облікового запису вказуються повноваження доступу (тобто привілеї); атрибути (за необхідності) вказуються для кожного облікового запису; для запитів на створення облікових записів потрібні схвалення від персоналу або ролей; облікові записи створюються відповідно до політики, процедур, передумов та критеріїв; облікові записи активуються відповідно до політики, процедур, передумов та критеріїв; облікові записи змінюються відповідно до політики, процедур, передумов та критеріїв; облікові записи деактивуються відповідно до політики, процедур, передумов та критеріїв; облікові записи видаляються відповідно до політики, процедур, передумов та критеріїв; контролюється використання облікових записів; адміністратори облікових записів та персонал або ролі отримують повідомлення протягом періоду часу, коли облікові записи більше не потрібні; адміністратори облікових записів та персонал або ролі отримують повідомлення протягом періоду часу, коли користувачі звільнені чи переведені; адміністратори облікових записів та персонал або ролі отримують повідомлення протягом періоду часу, коли використовуються індивідуальні системи або наявні зміни, які потребують нових знань. доступ до системи здійснюється на підставі дійсної авторизації доступу; доступ до системи авторизується на основі передбачуваного використання системи; доступ до системи авторизовано на основі атрибутів (за необхідності); AC-02j AC-02k.[01] AC-02k.[02] AC-02l.[01] AC-02l.[02] облікові записи переглядаються на відповідність вимогам управління обліковими записами частота; створено процес повторного випуску облікових даних спільного доступу або групових облікових записів (якщо вони розгорнуті), коли користувачів вилучено з групи; впроваджено процес повторного випуску облікових даних спільного доступу або групових облікових записів (якщо вони розгорнуті), коли користувачів вилучено з групи; процеси управління обліковими записами узгоджуються з процесами звільнення персоналу; процеси управління обліковими записами узгоджуються з процесами переводу персоналу\newline Тип: integer\newline Значення: \textbf{24} & AC-2 & a. Визначити та задокументувати типи облікових записів системи, дозволених для використання в ІС для підтримки цілей, завдань, функцій і процесів організації.\\ 
b. Призначити менеджерів облікових записів для управління системними обліковими записами.\\ 
c. Створити умови для групового та рольового членства.\\ 
d. Визначити авторизованих користувачів інформаційної системи, членство в групі та ролі, а також дозволи доступу (наприклад, привілеї) та інші атрибути (за потреби) для кожного облікового запису.\\ 
e. Вимагати схвалення [Призначення: визначеною організацією відповідальною особою або роллю] запитів на створення облікових записів системи.\\ 
f. Створювати, активувати, змінювати, деактивувати та видаляти системні облікові записи відповідно до [Призначення: визначених організацією політики, процедур та умов].\\ 
g. Впровадити моніторинг використання облікових записів системи.\\ 
h. Повідомляти адміністраторів облікових записів у межах [Призначення: визначеного організацією часового періоду для кожної ситуації]:\\ 
1. коли облікові записи більше не потрібні;\\ 
2. коли користувачі звільнені чи переведені;\\ 
3. коли використовуються індивідуальні системи або наявні зміни, які потребують нових знань.\\ 
i. Авторизувати доступ до системи на основі:\\ 
1. Дійсної авторизації доступу.\\ 
2. Передбачуваного використання системи.\\ 
3. Інших атрибутів, що вимагаються організацією.\\ 
j. Проводити перегляд облікових записів на відповідність вимогам управління обліковими записами з [Призначення: визначеною організацією частотою].\\ 
k. Впровадити процес повторного випуску облікових даних спільного/групового облікового запису (якщо він буде розгорнутий), коли особи виходять з групи.\\ 
l. Узгодити процеси управління обліковими записами з процесами звільнення та переводу (передачі повноважень) персоналу. \\ \hline
3 & УПРАВЛІННЯ ОБЛІКОВИМИ ЗАПИСАМИ - ДЕАКТИВАЦІЯ ОБЛІКОВИХ ЗАПИСІВ (AC-2(3)) & Параметр: Облікові записи деактивуються протягом часового періоду, коли облікові записи більше не пов'язані з користувачем або фізичною особою\newline Тип: list\newline Значення: \textbf{admin, security\_\allowbreak{}officer}\newline\newline Параметр: Облікові записи відключено протягом часового періоду, коли облікові записи порушують політику організації; облікові записи деактивуються протягом , якщо вони були неактивними протягом \newline Тип: list\newline Значення: \textbf{default\_\allowbreak{}deny\_\allowbreak{}rule, abac\_\allowbreak{}rule\_\allowbreak{}1}\newline\newline Параметр: Визначено період часу, протягом якого необхідно деактивувати облікові записи\newline Тип: integer\newline Значення: \textbf{30}\newline\newline Параметр: Визначено період часу неактивності, після закінчення якого облікові записи будуть деактивовані; AC-02(03)(a) облікові записи деактивуються протягом часового періоду, коли термін дії облікових записів минув\newline Тип: list\newline Значення: \textbf{login, logout, failed\_\allowbreak{}attempt} & AC-2(3) & Автоматично деактивувати облікові записи коли: a) їх строк дії минув; b) вони більше не пов'язані з користувачем; c) вони порушують організаційну політику; d) вони були неактивними впродовж [Призначення: визначеного організацією періоду часу]. \\ \hline
4 & УПРАВЛІННЯ ОБЛІКОВИМИ ЗАПИСАМИ - ВИХІД ІЗ СИСТЕМИ ЗА ВІДСУТНОСТІ АКТИВНОСТІ (AC-2(5)) & Параметр: Користувачі повинні виходити з системи, коли період очікуваної бездіяльності або опис часу, коли потрібно вийти з системи\newline Тип: integer\newline Значення: \textbf{30}\newline\newline Параметр: Визначено часовий період очікуваної бездіяльності або опис, коли потрібно вийти з системи\newline Тип: integer\newline Значення: \textbf{30} & AC-2(5) & Вимагати від користувачів виходити із системи, коли [Призначення: вичерпано визначений організацією періоду часу очікування або опис того, коли необхідно вийти із системи]. \\ \hline
5 & УПРАВЛІННЯ ОБЛІКОВИМИ ЗАПИСАМИ - ДЕАКТИВАЦІЯ ОБЛІКОВИХ ЗАПИСІВ ОСІБ З ВИСОКИМ РІВНЕМ РИЗИКУ (AC-2(13)) & Параметр: Облікові записи користувачів деактивуються протягом періоду часу з моменту виявлення значних ризиків\newline Тип: integer\newline Значення: \textbf{30}\newline\newline Параметр: Визначено період часу, протягом якого необхідно деактивувати облікові записи фізичних осіб, які становлять значний ризик\newline Тип: integer\newline Значення: \textbf{30} & AC-2(13) & Деактивувати облікові записи користувачів, які становлять значний ризик, у межах [Призначення: визначеного організацією періоду часу] після виявлення ризику. \\ \hline
6 & ЗАБЕЗПЕЧЕННЯ ДОСТУПУ (AC-3) & Параметр: Затверджені повноваження на логічний доступ до інформації та ресурсів системи виконуються відповідно до чинних політик(правил) управління доступом\newline Тип: list\newline Значення: \textbf{default\_\allowbreak{}deny\_\allowbreak{}rule, abac\_\allowbreak{}rule\_\allowbreak{}1} & AC-3 & Застосовувати затверджені повноваження для логічного доступу до інформації та ресурсів системи відповідно до чинної політики (правил) управління доступом. \\ \hline
7 & УПРАВЛІННЯ ІНФОРМАЦІЙНИМИ ПОТОКАМИ (AC-4) & Параметр: Затверджені повноваження застосовуються для контролю потоку інформації всередині системи та між підключеними системами на основі політики управління інформаційними потоками\newline Тип: list\newline Значення: \textbf{default\_\allowbreak{}deny\_\allowbreak{}rule, abac\_\allowbreak{}rule\_\allowbreak{}1}\newline\newline Параметр: Визначено політики управління інформаційними потоками всередині системи та між підключеними системами\newline Тип: list\newline Значення: \textbf{default\_\allowbreak{}deny\_\allowbreak{}rule, abac\_\allowbreak{}rule\_\allowbreak{}1} & AC-4 & Застосувати затверджені повноваження для управління потоком інформації всередині системи та між пов'язаними системами на основі [Призначення: визначеними організацією політиками управління інформаційним потоком]. \\ \hline
8 & РОЗМЕЖУВАННЯ ОБОВ'ЯЗКІВ (AC-5) &  & AC-5 & a. Розмежувати і документувати [Призначення: визначені організацією обов'язки окремих осіб].\\ 
b. Установити правила авторизації доступу для підтримки розмежування обов'язків. \\ \hline
9 & МІНІМІЗАЦІЯ ПОВНОВАЖЕНЬ (AC-6) &  & AC-6 & Впровадити принцип мінімізації повноважень, який дозволяє користувачам (або процесам, що діють від імені користувачів) здійснювати лише такі авторизовані звернення, які необхідні для виконання визначених завдань відповідно до цілей (призначення, місії) організації та функцій. \\ \hline
10 & МІНІМІЗАЦІЯ ПОВНОВАЖЕНЬ - АВТОРИЗОВАНИЙ ДОСТУП ДО ФУНКЦІЙ БЕЗПЕКИ (AC-6(1)) & Параметр: Визначені особи або ролі з авторизованим доступом до функцій безпеки та інформації, що має відношення до безпеки\newline Тип: list\newline Значення: \textbf{admin, security\_\allowbreak{}officer} & AC-6(1) &  \\ \hline
11 & МІНІМІЗАЦІЯ ПОВНОВАЖЕНЬ - НЕПРИВІЛЕЙОВАНИЙ ДОСТУП ДО НЕЗАХИЩЕНИХ ФУНКЦІЙ (AC-6(2)) & Параметр: Користувачі облікових записів (або ролей) системи з доступом до функцій безпеки або інформації, що стосується безпеки, повинні використовувати непривілейовані облікові записи або ролі під час доступу до незахищених функцій\newline Тип: list\newline Значення: \textbf{admin, security\_\allowbreak{}officer} & AC-6(2) & Вимагати від користувачів облікових записів системи або ролей, які мають доступ до [Призначення: визначених організацією функцій безпеки або інформації, що стосується безпеки], використовувати непривілейовані облікові записи чи ролі під час доступу до незахищених функцій. \\ \hline
12 & МІНІМІЗАЦІЯ ПОВНОВАЖЕНЬ - ПРИВІЛЕЙОВАНІ ОБЛІКОВІ ЗАПИСИ (AC-6(5)) & Параметр: Привілейовані облікові записи в системі обмежено персонал або ролі\newline Тип: list\newline Значення: \textbf{admin, security\_\allowbreak{}officer}\newline\newline Параметр: Визначено персонал або ролі, яким мають бути обмежені привілейовані облікові записи в системі\newline Тип: list\newline Значення: \textbf{admin, security\_\allowbreak{}officer} & AC-6(5) & Обмежити привілейовані облікові записи в системі згідно з [Призначення: визначеним організацією персоналом або ролями]. \\ \hline
13 & МІНІМІЗАЦІЯ ПОВНОВАЖЕНЬ - ПЕРЕГЛЯД ПОВНОВАЖЕНЬ КОРИСТУВАЧА (AC-6(7)) & Параметр: Повноваження, призначені ролям і класам, переглядаються з частотою для перевірки необхідності таких повноважень\newline Тип: integer\newline Значення: \textbf{30}\newline\newline Параметр: Визначено частоту перегляду повноважень, призначених ролям або класам користувачів\newline Тип: integer\newline Значення: \textbf{30}\newline\newline Параметр: Визначено ролі або класи користувачів, яким призначено повноваження \newline Тип: list\newline Значення: \textbf{admin, security\_\allowbreak{}officer} & AC-6(7) & a) Переглядати [Призначення: з визначеною організацією частотою] повноваження призначених для [Призначення: визначених організацією посад або класів користувачів] для перевірки необхідності таких повноважень; b) За необхідності перепризначити або зняти повноваження, правильного відображення цілей (місії) організації та потреб організації. для \\ \hline
14 & МІНІМІЗАЦІЯ ПОВНОВАЖЕНЬ - АУДИТ ВИКОРИСТАННЯ ПРИВІЛЕЙОВАНИХ ФУНКЦІЙ (AC-6(9)) &  & AC-6(9) & Реєструвати виконання привілейованих функцій. \\ \hline
15 & МІНІМІЗАЦІЯ ПОВНОВАЖЕНЬ - ЗАБОРОНА НЕПРИВІЛЕЙОВАНИМ КОРИСТУВАЧАМ ВИКОНУВАТИ ПРИВІЛЕЙОВАНІ ФУНКЦІЇ (AC-6(10)) &  & AC-6(10) & Вжити заходи для запобігання можливості виконувати привілейовані функції непривілейованими користувачами. \\ \hline
16 & НЕВДАЛІ СПРОБИ ВХОДУ В СИСТЕМУ (AC-7) & Параметр: Застосовано обмеження на кількість послідовних неуспішних спроб входу користувача протягом періоду часу\newline Тип: integer\newline Значення: \textbf{30}\newline\newline Параметр: Визначається кількість послідовних неуспішних спроб входу користувача, дозволених протягом певного періоду часу\newline Тип: integer\newline Значення: \textbf{30}\newline\newline Параметр: Визначено період часу, яким обмежується кількість послідовних неуспішних спроб входу користувача\newline Тип: integer\newline Значення: \textbf{30}\newline\newline Параметр: Вибрано одне або декілька з наступних ЗНАЧЕНЬ ПАРАМЕТРІВ: \{заблокувати обліковий запис або вузол на період часу; заблокувати обліковий запис або вузол до зняття адміністратором; затримати наступний запит на вхід за алгоритмом затримки; повідомити системного адміністратора; виконати іншу дію\}\newline Тип: string\newline Значення: \textbf{AES-256-GCM}\newline\newline Параметр: Період часу, на який буде заблоковано обліковий запис або вузол (якщо вибрано)\newline Тип: integer\newline Значення: \textbf{30}\newline\newline Параметр: Визначено алгоритм затримки наступного запиту на вхід (якщо вибрано)\newline Тип: string\newline Значення: \textbf{AES-256-GCM}\newline\newline Параметр: Інша дія, яка буде виконана після перевищення максимальної кількості невдалих спроб (якщо вибрано)\newline Тип: integer\newline Значення: \textbf{3} & AC-7 & a. Встановити обмеження на [Призначення: визначену організацією кількість] послідовних неуспішних спроб входу користувача в систему впродовж [Призначення: визначеного організацією часового періоду].\\ 
b. Автоматично виконати [Вибір (один або декілька): блокування облікового запису/вузла на [Призначення: визначений організацією часовий період]; блокування облікового запису/вузла, доки він не буде розблокований адміністратором; затримання наступної команди входу в систему за [Надання: визначеним організацією алгоритмом затримки]; виконати [Призначення: визначені організацією дії]], коли перевищено максимальну кількість невдалих спроб входу в систему. \\ \hline
17 & НЕВДАЛІ СПРОБИ ВХОДУ В СИСТЕМУ - ОЧИЩЕННЯ АБО СТИРАННЯ МОБІЛЬНОГО ПРИСТРОЮ (AC-7(2)) & Параметр: Інформація очищується або стирається з мобільних пристроїв на основі вимог або методів очищення або стирання після <AC-07(02)\_\allowbreak{}ODP[03 кількість> послідовних, невдалих спроб входу на пристрій\newline Тип: integer\newline Значення: \textbf{3}\newline\newline Параметр: Визначається кількість послідовних невдалих спроб входу в систему до того, як інформація буде очищена або стерта з мобільних пристроїв\newline Тип: integer\newline Значення: \textbf{3} & AC-7(2) & Очистити або стерти інформацію з [Призначення: визначених організацією мобільних пристроїв] на основі [Призначення: визначених організацією вимог та методик очищення чи стирання] після [Призначення: визначеної організацією кількості] послідовних невдалих спроб входу в систему з пристрою. \\ \hline
18 & ПОПЕРЕДЖЕННЯ ПРО ВИКОРИСТАННЯ СИСТЕМИ (AC-8) & Параметр: Визначається сповіщення або банер про використання системи, який система буде показувати користувачам перед наданням доступу до системи; визначено умови використання системи, які система відображатиме перед наданням подальшого доступу; повідомлення про використання системи відображається користувачам перед наданням доступу до системи, яке містить повідомлення про конфіденційність і безпеку, що відповідають чинним законам, нормативним документам, наказам, директивам, політикам, правилам, стандартам і керівним принципам; у повідомленні про використання системи зазначено, що користувачі отримують доступ до урядової системи; у повідомленні про використання системи зазначено, що використання системи може контролюватися, реєструватися та підлягати аудиту; у повідомленні про використання системи зазначено, що несанкціоноване використання системи заборонено і тягне за собою кримінальну та цивільну відповідальність; у повідомленні про використання системи зазначено, що використання системи означає згоду на моніторинг і запис дій; сповіщення або банер залишається на екрані до тих пір, поки користувачі не визнають умови використання та не приймуть явні дії для входу в систему або подальшого доступу до неї; для загальнодоступних систем інформація про використання системи умови відображається перед наданням подальшого доступу до загальнодоступної системи; для загальнодоступних систем відображаються будь-які посилання на моніторинг, запис або аудит, які узгоджуються з положеннями про конфіденційність таких систем, що зазвичай забороняють ці види діяльності; для загальнодоступних систем додається опис авторизованого використання системи\newline Тип: list\newline Значення: \textbf{default\_\allowbreak{}deny\_\allowbreak{}rule, abac\_\allowbreak{}rule\_\allowbreak{}1} & AC-8 & a.\\ 
b. Демонструвати користувачам [Призначення: визначене організацією сповіщення або банер про використання системи] перед тим, як надавати доступ до системи, що забезпечує безпеку та приватність відповідно до чинних законів, нормативних документів, наказів, директив, політик, правил, стандартів і керівних принципів, які зазначають, що:\\ 
1. користувачі здійснюють доступ до урядової системи;\\ 
2. використання системи може контролюватися, реєструватися та підлягати аудиту;\\ 
3. несанкціоноване використання системи забороняється та приводить до кримінальної та цивільної відповідальності;\\ 
4. використання системи означає згоду на моніторинг і запис дій користувача. Зберігати сповіщення або банер на екрані, доки користувачі не визнають умови використання та не приймуть явних дій для входу в систему або подальшого доступу до системи.\\ 
c. Для загальнодоступних систем:\\ 
1. демонструвати інформацію про умови використання системи [Призначення: визначені організацією умови], перш ніж надавати подальший доступ до загальнодоступної системи;\\ 
2. демонструвати посилання, якщо такі є, на моніторинг, запис або аудит, які узгоджуються з акомодацією приватності для таких систем, які зазвичай забороняють такі дії;\\ 
3. мати опис авторизованого використання системи. \\ \hline
19 & БЛОКУВАННЯ ПРИСТРОЮ (AC-11) & Параметр: Вибрано одне або декілька з наступних ЗНАЧЕНЬ ПАРАМЕТРІВ: \{ініціювання блокування пристрою після періоду неактивності; вимога до користувача ініціювати блокування пристрою перед тим, як залишити систему без нагляду\}\newline Тип: integer\newline Значення: \textbf{30}\newline\newline Параметр: Часовий проміжок бездіяльності, після якого ініціюється блокування пристрою (якщо вибрано)\newline Тип: integer\newline Значення: \textbf{30} & AC-11 & a. Заборонити подальший доступ до системи шляхом ініціювання блокування пристрою після [Призначення: визначеного організацією періоду] бездіяльності або після отримання запиту від користувача.\\ 
b. Зберігати блокування пристрою, поки користувач не відновить доступ, використовуючи встановлені процедури ідентифікації та автентифікації. \\ \hline
20 & БЛОКУВАННЯ ПРИСТРОЮ - ПРИХОВАНІ ДИСПЛЕЇ (AC-11(1)) &  & AC-11(1) &  \\ \hline
21 & Припинення сеансу (AC-12) & Параметр: Визначено умови або події, що вимагають припинення сеансу\newline Тип: list\newline Значення: \textbf{login, logout, failed\_\allowbreak{}attempt} & AC-12 & Сеанс користувача має завершуватися автоматично після [Призначення: визначених організацією умов або тригерних подій, що вимагають припинення сеансу]. \\ \hline
22 & АВТОМАТИЗОВАНЕ МАРКУВАННЯ (AC-15) &  & AC-15 &  \\ \hline
23 & АТРИБУТИ БЕЗПЕКИ ТА ПРИВАТНОСТІ (AC-16) & Параметр: Визначено типи атрибутів безпеки, які мають бути пов'язані зі значеннями атрибутів безпеки для інформації, що зберігається, обробляється та/або передається; визначено значення атрибутів конфіденційності для типів атрибутів конфіденційності; визначено системи, для яких мають бути встановлені дозволені атрибути безпеки; визначено системи, для яких мають бути встановлені дозволені атрибути конфіденційності; визначено атрибути безпеки, визначені як частина AC-16(a), які дозволені для систем; визначено атрибути конфіденційності, визначені як частина AC-16(a), які дозволені для систем; визначено значення атрибутів або діапазони для встановлених атрибутів; визначено частоту, з якою слід переглядати атрибути безпеки на предмет відповідності; визначено частоту, з якою слід переглядати атрибути конфіденційності на предмет відповідності; надано засоби для асоціювання типів атрибутів безпеки зі значеннями атрибутів безпеки для інформації, що зберігається, обробляється та/або передається надано засоби для асоціювання типів атрибутів конфіденційності зі значеннями атрибутів конфіденційності для інформації, що зберігається, обробляється та/або передається виникають зв'язки з атрибутами; зв'язки атрибутів зберігаються разом з інформацією; на основі атрибутів, визначених у AC-16\_\allowbreak{}ODP[01] для систем, встановлюються наступні дозволені атрибути безпеки: атрибути безпеки; на основі атрибутів, визначених у AC-16\_\allowbreak{}ODP[02] для систем, встановлюються наступні дозволені атрибути приватності: атрибути конфіденційності ; AC-16(d)\newline Тип: list\newline Значення: \textbf{default\_\allowbreak{}deny\_\allowbreak{}rule, abac\_\allowbreak{}rule\_\allowbreak{}1}\newline\newline Параметр: Визначено наступні допустимі значення або діапазони атрибутів для кожного з встановлених атрибутів: значення або діапазони атрибутів; проводиться аудит змін до атрибутів; атрибути безпеки перевіряються на відповідність частота; атрибути конфіденційності перевіряються на відповідність частота\newline Тип: list\newline Значення: \textbf{default\_\allowbreak{}deny\_\allowbreak{}rule, abac\_\allowbreak{}rule\_\allowbreak{}1} & AC-16 & a. Визначити засоби для асоціювання (пов'язання) [Призначення: визначених організацією типів атрибутів безпеки та приватності], що приймають [Призначення: визначені організацією значення атрибутів безпеки та приватності] з інформацією, яка зберігається, обробляється та/або передається.\\ 
b. Пов'язані атрибути безпеки та приватності мають створюватися і зберігатися разом з інформацією.\\ 
c. Встановити дозволені [Призначення: визначені організацією атрибути безпеки] для [Призначення: систем, визначених організацією].\\ 
d. Визначити дозволені [Призначення: визначені організацією значення або діапазони] для кожного з встановлених атрибутів безпеки та приватності. \\ \hline
24 & ВІДДАЛЕНИЙ ДОСТУП (AC-17) &  & AC-17 & a. Встановити та задокументувати обмеження на використання, вимоги до конфігурації/підключення та рекомендації щодо здійснення кожного типу віддаленого доступу.\\ 
b. Авторизувати віддалений доступ до системи, перш ніж будуть дозволені такі підключення. \\ \hline
25 & ВІДДАЛЕНИЙ ДОСТУП - КЕРОВАНІ ТОЧКИ КОНТРОЛЮ ДОСТУПУ (AC-17(3)) &  & AC-17(3) &  \\ \hline
26 & ВІДДАЛЕНИЙ ДОСТУП - ПРИВІЛЕЙОВАНІ КОМАНДИ ТА ДОСТУП (AC-17(4)) &  & AC-17(4) &  \\ \hline
27 & БЕЗДРОТОВИЙ ДОСТУП (AC-18) &  & AC-18 & a. Установити обмеження на використання, вимоги до конфігурації/підключення та рекомендації щодо здійснення бездротового доступу.\\ 
b. Авторизувати бездротовий доступ до системи, перш ніж будуть дозволені такі підключення. \\ \hline
28 & БЕЗДРОТОВИЙ ДОСТУП - АВТЕНТИФІКАЦІЯ ТА ШИФРУВАННЯ (AC-18(1)) & Параметр: Бездротовий доступ до системи захищений за допомогою шифрування\newline Тип: string\newline Значення: \textbf{AES-256-GCM} & AC-18(1) &  \\ \hline
29 & БЕЗДРОТОВИЙ ДОСТУП - ВІДКЛЮЧЕННЯ БЕЗДРОТОВОЇ МЕРЕЖІ (AC-18(3)) &  & AC-18(3) &  \\ \hline
30 & КОНТРОЛЬ ДОСТУПУ ДЛЯ МОБІЛЬНИХ ПРИСТРОЇВ (AC-19) &  & AC-19 & a. Встановити обмеження на використання, вимоги до конфігурації, вимоги до підключення і рекомендації щодо впровадження мобільних пристроїв, контрольованих організацією.\\ 
b. Авторизувати підключення мобільних пристроїв до систем, які експлуатуються організацією. \\ \hline
31 & КОНТРОЛЬ ДОСТУПУ ДЛЯ МОБІЛЬНИХ ПРИСТРОЇВ - ВИКОРИСТАННЯ ПИСЬМОВИХ ТА ПОРТАТИВНИЙ ПРИСТРОЇВ ДЛЯ ЗБЕРІГАННЯ ДАНИХ (AC-19(1)) &  & AC-19(1) &  \\ \hline
32 & КОНТРОЛЬ ДОСТУПУ ДЛЯ МОБІЛЬНИХ ПРИСТРОЇВ - ВИКОРИСТАННЯ ПЕРСОНАЛЬНИХ ПОРТАТИВНИХ ПРИСТРОЇВ ЗБЕРІГАННЯ ДАНИХ (AC-19(2)) & Параметр: КОНТРОЛЬ ДОСТУПУ ДЛЯ МОБІЛЬНИХ ПРИСТРОЇВ - ВИКОРИСТАННЯ ПЕРСОНАЛЬНИХ ПОРТАТИВНИХ ПРИСТРОЇВ ЗБЕРІГАННЯ ДАНИХ [Вилучено: Включено в MP-07]. AC-19(03) КОНТРОЛЬ ДОСТУПУ ДЛЯ МОБІЛЬНИХ ПРИСТРОЇВ - ВИКОРИСТАННЯ ПОРТАТИВНИХ ПРИСТРОЇВ ЗБЕРІГАННЯ ДАНИХ З НЕІДЕНТИФІКОВАНИМ ВЛАСНИКОМ [Вилучено: Включено в MP-07]\newline Тип: list\newline Значення: \textbf{admin, security\_\allowbreak{}officer} & AC-19(2) &  \\ \hline
33 & КОНТРОЛЬ ДОСТУПУ ДЛЯ МОБІЛЬНИХ ПРИСТРОЇВ - ОБМЕЖЕННЯ ДЛЯ ЗАСЕКРЕЧЕНОЇ ІНФОРМАЦІЇ (AC-19(4)) & Параметр: Використання незахищених мобільних пристроїв на об'єктах, що містять системи, які обробляють, зберігають або передають секретну інформацію, заборонено, за винятком випадків, коли на це є спеціальний дозвіл уповноваженої посадової особи; AC-19(04)(b)(01) заборона підключення незахищених мобільних пристроїв до систем з обмеженим доступом застосовується до осіб, яким уповноважена посадова особа дозволила використовувати незахищені мобільні пристрої на об'єктах, що містять системи, які обробляють, зберігають або передають інформацію з обмеженим доступом; AC-19(04)(b)(02) дозвіл уповноваженої посадової особи на підключення незахищених мобільних пристроїв до незахищених систем вимагається від осіб, яким дозволено використовувати незахищені мобільні пристрої на об'єктах, що містять системи, які обробляють, зберігають або передають інформацію з обмеженим доступом; AC-19(04)(b)(03) заборона використання внутрішніх або зовнішніх модемів чи бездротових інтерфейсів у складі незахищених мобільних пристроїв поширюється на осіб, яким уповноваженою посадовою особою дозволено використовувати незахищені мобільні пристрої під час виконання службових обов'язків, а також на осіб, які не мають права на використання таких пристроїв AC-19(04)(b)(04)[01] вибірковий огляд та перевірка незахищених мобільних пристроїв та інформації, що зберігається на них, посадовими особами є обов'язковими; AC-19(04)(b)(04)[02] дотримання політики обробки інцидентів застосовується у разі виявлення секретної інформації під час випадкового огляду та перевірки незахищених мобільних пристроїв\newline Тип: list\newline Значення: \textbf{admin, security\_\allowbreak{}officer}\newline\newline Параметр: Підключення захищенихмобільних пристроїв до засекречених систем обмежено відповідно до політики безпеки\newline Тип: list\newline Значення: \textbf{default\_\allowbreak{}deny\_\allowbreak{}rule, abac\_\allowbreak{}rule\_\allowbreak{}1}\newline\newline Параметр: Визначено посадових осіб із захисту інформації, відповідальних за огляд та перевірку захищених мобільних пристроїв та інформації, що зберігається на цих пристроях\newline Тип: list\newline Значення: \textbf{admin, security\_\allowbreak{}officer}\newline\newline Параметр: Визначено політики безпеки, що обмежують підключення захищених мобільних пристроїв до засекречених систем\newline Тип: list\newline Значення: \textbf{default\_\allowbreak{}deny\_\allowbreak{}rule, abac\_\allowbreak{}rule\_\allowbreak{}1} & AC-19(4) &  \\ \hline
34 & КОНТРОЛЬ ДОСТУПУ ДЛЯ МОБІЛЬНИХ ПРИСТРОЇВ - ПОВНЕ ШИФРУВАННЯ ПРИСТРОЇВ ТА СХОВИЩ ІНФОРМАЦІЇ (AC-19(5)) &  & AC-19(5) &  \\ \hline
35 & ВИКОРИСТАННЯ ЗОВНІШНІХ СИСТЕМ (AC-20) & Параметр: Вибрано одне або більше з наступних ЗНАЧЕНЬ ПАРАМЕТРІВ: \{встановити умови та положення; визначити заходи захисту\}\newline Тип: list\newline Значення: \textbf{}\newline\newline Параметр: Визначено умови та положення, що відповідають довірчим відносинам, встановленим з іншими організаціями, які володіють, експлуатують та/або обслуговують зовнішні системи (якщо вибрано)\newline Тип: list\newline Значення: \textbf{} & AC-20 & a. [Вибір (один або кілька): Встановіть [Призначення: умови, визначені організацією]; Визначте [Призначення: визначені організацією засоби контролю, які, як стверджується, будуть реалізовані на зовнішніх системах]], узгоджені з довірчими відносинами, встановленими з іншими організаціями, які володіють, експлуатують та/або обслуговують зовнішні системи, дозволяючи уповноваженим особам:\\ 
1. доступ до системи із зовнішніх систем;\\ 
2. обробляти, зберігати або передавати керовану організацією інформацію за допомогою зовнішніх систем;\\ 
b. Заборонити використання [Призначення: організаційно-визначені типи зовнішніх систем]. \\ \hline
36 & ВИКОРИСТАННЯ ЗОВНІШНІХ СИСТЕМ - ОБМЕЖЕННЯ НА АВТОРИЗОВАНЕ ВИКОРИСТАННЯ (AC-20(1)) & Параметр: Авторизовані особи мають право використовувати зовнішню систему для доступу до системи або для обробки, зберігання чи передачі інформації, що контролюється організацією, лише після перевірки виконання заходів безпеки та конфіденційності, зазначених у політиці безпеки та конфіденційності організації, а також планах безпеки та конфіденційності (якщо такі застосовуються)\newline Тип: list\newline Значення: \textbf{admin, security\_\allowbreak{}officer}\newline\newline Параметр: Авторизовані особи мають право використовувати зовнішню систему для доступу до системи або для обробки, зберігання чи передачі інформації, що контролюється організацією, лише після збереження погоджених угод про підключення або обробку системи з структурою організації, на якій розміщена зовнішня система\newline Тип: list\newline Значення: \textbf{admin, security\_\allowbreak{}officer} & AC-20(1) &  \\ \hline
37 & ВИКОРИСТАННЯ ЗОВНІШНІХ СИСТЕМ - ПЕРЕНОСНІ ПРИСТРОЇ ЗБЕРІГАННЯ ДАНИХ (AC-20(2)) & Параметр: Використання портативних пристроїв носіїв інформації уповноваженими особами обмежено у зовнішніх системах за допомогою обмеження \newline Тип: list\newline Значення: \textbf{admin, security\_\allowbreak{}officer}\newline\newline Параметр: Визначено обмеження на використання авторизованими особами портативних носіїв інформації у зовнішніх системах\newline Тип: list\newline Значення: \textbf{admin, security\_\allowbreak{}officer} & AC-20(2) &  \\ \hline
38 & ПУБЛІЧНО ДОСТУПНИЙ КОНТЕНТ (AC-22) & Параметр: Визначені особи уповноважені на розміщення інформації в загальнодоступній системі\newline Тип: list\newline Значення: \textbf{admin, security\_\allowbreak{}officer}\newline\newline Параметр: Уповноважені особи проходять навчання, щоб гарантувати, що загальнодоступна інформація не містить інформацію з обмеженим доступом\newline Тип: list\newline Значення: \textbf{admin, security\_\allowbreak{}officer}\newline\newline Параметр: Зміст у загальнодоступній системі перевіряється на наявність інформації з обмеженим доступом з частотою\newline Тип: integer\newline Значення: \textbf{30}\newline\newline Параметр: Визначено частоту, з якою слід переглядати вміст загальнодоступної системи на предмет наявності там інформації з обмеженим доступом\newline Тип: integer\newline Значення: \textbf{30} & AC-22 & a. Призначити осіб, що уповноважені на розміщення інформації в загальнодоступній системі.\\ 
b. Навчати уповноважених осіб тому, щоб загальнодоступна інформація не містила інформацію з обмеженим доступом.\\ 
c. Переглядати запропонований зміст інформації до публікації в загальнодоступній системі, щоб гарантувати, що там не міститься інформація з обмеженим доступом.\\ 
d. Переглядати вміст загальнодоступної системи на предмет наявності там інформації з обмеженим доступом з [Призначення: визначеною організацією частотою]; така інформація має бути видалена в разі її виявлення. \\ \hline
\multicolumn{5}{|l|}{\textbf{2. ОБІЗНАНІСТЬ ТА НАВЧАННЯ (AT) (AT)}} \\ \hline
39 & ПОЛІТИКА ТА ПРОЦЕДУРИ ПІДВИЩЕННЯ ОБІЗНАНОСТІ ТА НАВЧАННЯ (AT-1) &  & AT-1 & a. Розробити, задокументувати та поширити [Призначення: серед визначеного організацією персоналу або ролей]:\\ 
1. [Вибір (один або декілька): Рівень організації; Рівень місії/бізнес-процесу; рівень системи] політики обізнаності та навчання у сфері забезпечення безпеки та приватності, яка:\\ 
(a) містить мету, сферу застосування, ролі, обов'язки, відповідальність керівництва, координацію між організаційними підрозділами та систему контролю відповідності (complaince);\\ 
(b) відповідає чинним законам, нормативним документам, директивам, нормам, політикам, стандартам та керівним документам.\\ 
2. Процедури, що сприяють реалізації політики підвищення обізнаності та професійної підготовки в галузі безпеки, приватності, а також пов'язаних з ними заходів захисту інформації та персональних даних.\\ 
b. Призначити [Призначення: визначену організацією посадову особу] для управління політикою та процедурами підвищення обізнаності та навчання у сфері забезпечення безпеки та приватності.\\ 
c. Переглядати та оновлювати:\\ 
1. Поточну політика [Призначення: частота, визначена організацією] і наступне [Призначення: події, визначені організацією];\\ 
2. Процедури [Призначення: частота, визначена організацією] та наступні [Призначення: події, визначені організацією]. \\ \hline
40 & НАВЧАННЯ З ПІДВИЩЕННЯ ОБІЗНАНОСТІ (AT-2) & Параметр: Оновлюється зміст навчання грамотності та підвищення обізнаності з частотою\newline Тип: integer\newline Значення: \textbf{30}\newline\newline Параметр: Визначено періодичність проведення навчання грамотності з питань безпеки для користувачів системи (в тому числі менеджерів, вищого керівництва та підрядників) після початкового тренінгу\newline Тип: string\newline Значення: \textbf{щорічно}\newline\newline Параметр: Визначено періодичність проведення навчання грамотності з питань конфіденційності для користувачів системи (в тому числі менеджерів, вищого керівництва та підрядників) після початкового тренінгу\newline Тип: string\newline Значення: \textbf{щорічно}\newline\newline Параметр: Визначено події, які потребують навчання користувачів системи грамотності з питань безпеки\newline Тип: list\newline Значення: \textbf{login, logout, failed\_\allowbreak{}attempt}\newline\newline Параметр: Визначено частоту оновлення навчання грамотності та змісту обізнаності; AT-02\_\allowbreak{}ODP[07] визначено події після яких необхідне оновлення навчання грамотності та змісту обізнаності; AT-02(a)[01][01] навчання з грамотності з питань безпеки надається користувачам системи (включаючи менеджерів, керівників вищої ланки та підрядників) як частина початкового навчання для нових користувачів; AT-02(a)[01][02] навчання з грамотності з питань конфіденційності надається користувачам системи (включаючи менеджерів, керівників вищої ланки та підрядників) як частина початкового навчання для нових користувачів; AT-02(a)[01][03] для користувачів системи (включно з менеджерами, вищим керівництвом та підрядниками) проводиться навчання з безпеки з періодичністю після цього; AT-02(a)[01][04] для користувачів системи (включно з менеджерами, вищим керівництвом та підрядниками) проводиться навчання з конфіденційності з періодичністю після цього; AT-02(a)[02][01] тренінги з грамотності з питань безпеки проводяться для користувачів системи (включаючи менеджерів, керівників вищої ланки та підрядників), коли цього вимагають зміни в системі або після подій; AT-02(a)[02][02] тренінги з грамотності з питань конфіденціності проводяться для користувачів системи (включаючи менеджерів, керівників вищої ланки та підрядників), коли цього вимагають зміни в системі або після подій\newline Тип: list\newline Значення: \textbf{login, logout, failed\_\allowbreak{}attempt}\newline\newline Параметр: Визначено події, які потребують навчання користувачів системи грамотності з питань конфіденційності\newline Тип: list\newline Значення: \textbf{login, logout, failed\_\allowbreak{}attempt} & AT-2 & Впровадити базові тренінги з підвищення обізнаності у сфері безпеки та приватності для користувачів системи (включно з менеджерами, керівниками компаній і підрядниками):\\ 
a. Забезпечити навчання грамотності з питань безпеки та конфіденційності для користувачів системи (включаючи менеджерів, керівників вищої ланки та підрядників):\\ 
1. як частину початкового навчання для нових користувачів і [Призначення: частота, визначена організацією] після цього;\\ 
2. якщо цього потребують системні зміни або наступні [Призначення: події, визначені організацією].\\ 
b. Використовувати наведені нижче методи, щоб підвищити рівень безпеки та конфіденційності користувачів системи [Завдання: визначені організацією методи поінформованості];\\ 
c. Оновлювати навчання грамотності та зміст обізнаності [Завдання: частота, визначена організацією] і наступні [Завдання: події, визначені організацією];\\ 
d. Включити уроки, отримані з внутрішніх або зовнішніх інцидентів безпеки або порушень, у навчання грамотності та методи підвищення обізнаності. \\ \hline
41 & НАВЧАННЯ З ПІДВИЩЕННЯ ОБІЗНАНОСТІ | НАВЧАННЯ ЩОДО РОЗПІЗНАВАННЯ ВНУТРІШНЬОЇ ЗАГРОЗИ &  & AT-2(2) & Ввести до програми навчання вправи з розпізнавання та виявлення потенційних індикаторів внутрішніх загроз. \\ \hline
42 & НАВЧАННЯ З ПІДВИЩЕННЯ ОБІЗНАНОСТІ | НАВЧАННЯ З ПИТАНЬ СОЦІАЛЬНОЇ ІНЖЕНЕРІЇ &  & AT-2(3) & Ввести до програми навчання вправи з підвищення обізнаності щодо розпізнавання та повідомлення про потенційні та фактичні атаки, з використанням методів соціальної інженерії та інтелектуального аналізу соціальних даних. \\ \hline
43 & РОЛЬОВЕ НАВЧАННЯ (AT-3) & Параметр: Оновлюється вміст навчання на основі ролей з частотою\newline Тип: integer\newline Значення: \textbf{30}\newline\newline Параметр: Визначено ролі та обов'язки для тренінгів з безпеки на основі ролей\newline Тип: list\newline Значення: \textbf{admin, security\_\allowbreak{}officer}\newline\newline Параметр: Визначено ролі та обов'язки для тренінгів з конфіденційності на основі ролей\newline Тип: list\newline Значення: \textbf{admin, security\_\allowbreak{}officer}\newline\newline Параметр: Визначено частоту проведення тренінгів на основі ролей з безпеки та конфіденційності для призначеного персоналу після початкової підготовки\newline Тип: list\newline Значення: \textbf{admin, security\_\allowbreak{}officer}\newline\newline Параметр: Визначено частоту оновлення змісту навчання на основі ролей\newline Тип: integer\newline Значення: \textbf{30}\newline\newline Параметр: Визначено події, які потребують оновлення змісту навчання на основі ролей; AT-03(a)[01][01] навчання з безпеки на основі ролей проводиться для ролей та обов'язків перед авторизацією доступу до системи, інформації або виконанням призначених обов'язків; AT-03(a)[01][02] навчання з конфіденційності на основі ролей проводиться для ролей та обов'язків перед авторизацією доступу до системи, інформації або виконанням призначених обов'язків; AT-03(a)[01][03] навчання з безпеки на основі ролей проводиться для ролей та обов'язків з частотою після цього; AT-03(a)[01][04] навчання з конфіденційності на основі ролей проводиться для ролей та обов'язків з частотою після цього; AT-03(a)[02][01] навчання з питань безпеки на основі ролей проводиться для персоналу, який виконує певні ролі та обов'язки у сфері безпеки, коли цього вимагають зміни в системі; AT-03(a)[02][02] навчання з питань конфіденційності на основі ролей проводиться для персоналу, який виконує певні ролі та обов'язки у сфері безпеки, коли цього вимагають зміни в системі\newline Тип: list\newline Значення: \textbf{admin, security\_\allowbreak{}officer} & AT-3 & a. Забезпечити проведення навчання з питань безпеки та приватності на основі ролей для працівників з ролями та обов'язками: [Призначення: визначені організацією ролі та обов'язки]:\\ 
1. перед авторизацією доступу до системи, інформації або виконанням призначених обов'язків і [Призначення: частота, визначена організацією] після цього;\\ 
2. коли цього потребують системні зміни.\\ 
b. Оновити навчальний контент на основі ролей [Призначення: частота, визначена організацією] і наступні [Призначення: події, визначені організацією];\\ 
c. Включіть у рольове навчання, інформацію, отриману з внутрішніх або зовнішніх інцидентів та порушень безпеки. \\ \hline
\multicolumn{5}{|l|}{\textbf{3. АУДИТ ТА ПІДЗВІТНІСТЬ (AU) (AU)}} \\ \hline
44 & Політика та процедури аудиту та підзвітності (AU-1) & Параметр: Розроблено та задокументовано політику аудиту та підзвітності\newline Тип: list\newline Значення: \textbf{default\_\allowbreak{}deny\_\allowbreak{}rule, abac\_\allowbreak{}rule\_\allowbreak{}1}\newline\newline Параметр: Політика аудиту та підзвітності доведена до персонал або ролі\newline Тип: list\newline Значення: \textbf{admin, security\_\allowbreak{}officer}\newline\newline Параметр: Розроблені та задокументовані процедури аудиту та підзвітності, що сприяють впровадженню політики аудиту та підзвітності, а також відповідні заходи контролю аудиту та підзвітності\newline Тип: list\newline Значення: \textbf{default\_\allowbreak{}deny\_\allowbreak{}rule, abac\_\allowbreak{}rule\_\allowbreak{}1}\newline\newline Параметр: Посадова особа призначається для управління політикою та процедурами аудиту та підзвітності AU-01(c)[01][01] переглядається та оновлюється поточна політика аудиту та підзвітності з частота; AU-01(c)[01][02] переглядається та оновлюється поточна політика аудиту та підзвітності після подій; AU-01(c)[02][01] переглядається та оновлюється поточні процедури аудиту та підзвітності з частота; AU-01(c)[02][02] переглядається та оновлюється поточні процедури аудиту та підзвітності після подій\newline Тип: list\newline Значення: \textbf{admin, security\_\allowbreak{}officer}\newline\newline Параметр: Визначено персонал або ролі, до яких має бути доведена політика аудиту та підзвітності\newline Тип: list\newline Значення: \textbf{admin, security\_\allowbreak{}officer}\newline\newline Параметр: Визначено персонал або ролі, на які поширюються процедури аудиту та підзвітності\newline Тип: list\newline Значення: \textbf{admin, security\_\allowbreak{}officer}\newline\newline Параметр: Визначено посадову особу, яка управлятиме політикою та процедурами аудиту та підзвітності\newline Тип: list\newline Значення: \textbf{admin, security\_\allowbreak{}officer}\newline\newline Параметр: Визначено частоту, з якою переглядається та оновлюється поточна політика аудиту та підзвітності\newline Тип: list\newline Значення: \textbf{default\_\allowbreak{}deny\_\allowbreak{}rule, abac\_\allowbreak{}rule\_\allowbreak{}1}\newline\newline Параметр: Визначено події, які потребують перегляду та оновлення поточної політики аудиту та підзвітності\newline Тип: list\newline Значення: \textbf{default\_\allowbreak{}deny\_\allowbreak{}rule, abac\_\allowbreak{}rule\_\allowbreak{}1}\newline\newline Параметр: Визначено частоту, з якою переглядаються та оновлюються поточні процедури аудиту та підзвітності\newline Тип: integer\newline Значення: \textbf{30}\newline\newline Параметр: Визначено події, які потребують перегляду та оновлення поточної процедури аудиту та підзвітності\newline Тип: list\newline Значення: \textbf{login, logout, failed\_\allowbreak{}attempt} & AU-1 & a. Розробити, задокументувати та поширити [Призначення: серед персоналу або ролей, що їх визначила організація]:\\ 
1. [Вибір (один або декілька): Рівень організації; Рівень місії/бізнес-процесу; рівень системи] політика аудиту та підзвітності, яка:\\ 
(a) містить мету, сферу застосування, ролі, обов'язки, відповідальність керівництва, координацію між організаційними підрозділами та систему контролю відповідності (complaince);\\ 
(b) відповідає чинним законам, нормативним документам, директивам, нормам, політикам, стандартам та керівним документам.\\ 
2. Процедури, що сприяють здійсненню політики аудиту та підзвітності, а також пов'язані з ними заходи аудиту та підзвітності.\\ 
b. Призначити [Призначення: визначену організацією старшу посадову особу] для управління політикою та процедурами аудиту та підзвітності.\\ 
c. Переглядати та оновлювати поточний аудит та підзвітність:\\ 
1. політику [Призначення: частота, визначена організацією] та наступне [Призначення: події, визначені організацією];\\ 
2. процедури аудиту [Призначення: визначеною організацією частотою] та [Завдання: події, визначені організацією]. \\ \hline
45 & ПОДІЇ АУДИТУ (AU-2) & Параметр: Зазначені типи подій реєструються 02\_\allowbreak{}ODP[03] частота або ситуація>; <AU-\newline Тип: string\newline Значення: \textbf{щорічно}\newline\newline Параметр: Переглядаються та оновлюються типи подій, вибрані для реєстрації, частота. у системі\newline Тип: string\newline Значення: \textbf{щорічно}\newline\newline Параметр: Визначено частоту або ситуацію, що вимагає проведення аудиту для кожної ідентифікованої події\newline Тип: list\newline Значення: \textbf{login, logout, failed\_\allowbreak{}attempt}\newline\newline Параметр: Частота перегляду та оновлення типів подій, обраних для журналювання\newline Тип: string\newline Значення: \textbf{щорічно} & AU-2 & a. Визначити типи подій, які система може реєструвати для підтримки функції аудиту: [Призначення: типи подій, визначені організацією, які система здатна реєструвати];\\ 
b. Координувати функції аудиту безпеки з іншими організаційними підрозділами, які вимагають інформації, пов'язаної з аудитом, для посилення взаємної підтримки та допомоги у виборі типів подій, що перевіряються;\\ 
c. Визначити, які типи подій підлягають аудиту: [Призначення: визначені організацією події, що підлягають аудиту (підмножина подій, що підлягають аудиту, визначених в AU-2 a.), а також частота (або ситуація, що вимагає) проведення аудиту для кожної ідентифікованої події]\\ 
d. Обґрунтувати, чому типи подій, що перевіряються, вважаються достатніми для підтримки розслідувань інцидентів (постфактум), пов'язаних з безпекою та приватністю;\\ 
e. Перегляньте й оновіть типи подій, вибрані для журналювання [Призначення: частота, визначена організацією]. \\ \hline
46 & ЗМІСТ ЗАПИСІВ АУДИТУ (AU-3) & Параметр: Записи аудиту містять інформацію, яка встановлює який тип події стався\newline Тип: list\newline Значення: \textbf{login, logout, failed\_\allowbreak{}attempt}\newline\newline Параметр: Записи аудиту містять інформацію, яка встановлює джерело події\newline Тип: list\newline Значення: \textbf{login, logout, failed\_\allowbreak{}attempt}\newline\newline Параметр: Записи аудиту містять інформацію, яка встановлює наслідки події\newline Тип: list\newline Значення: \textbf{login, logout, failed\_\allowbreak{}attempt}\newline\newline Параметр: Записи аудиту містять інформацію, яка встановлює результат події та ідентифікатор будь-яких осіб або суб'єктів, пов'язаних з подією\newline Тип: list\newline Значення: \textbf{login, logout, failed\_\allowbreak{}attempt} & AU-3 & Переконатися, що записи аудиту містять інформацію, яка встановлює наступне:\\ 
a. який тип події стався;\\ 
b. коли відбулася подія;\\ 
c. де відбулася подія;\\ 
d. джерело події;\\ 
e. наслідки події;\\ 
f. результат події та ідентифікатор будь-яких осіб або суб'єктів, пов'язаних з подією. \\ \hline
47 & ЗМІСТ ЗАПИСІВ АУДИТУ - ДОДАТКОВА ІНФОРМАЦІЯ ПРО АУДИТ (AU-3(1)) &  & AU-3(1) &  \\ \hline
48 & РЕАГУВАННЯ НА ВІДМОВИ ОБРОБКИ ДАНИХ АУДИТУ (AU-5) & Параметр: Персонал або ролі отримують сповіщення у разі збою процесу обробки даних аудиту періоду часу\newline Тип: list\newline Значення: \textbf{admin, security\_\allowbreak{}officer}\newline\newline Параметр: Додаткові дії виконуються у разі збою процесу обробки даних аудиту\newline Тип: list\newline Значення: \textbf{login, logout, failed\_\allowbreak{}attempt}\newline\newline Параметр: Визначено персонал або ролі, які отримують сповіщення про збої в процесі обробки даних аудиту \newline Тип: list\newline Значення: \textbf{admin, security\_\allowbreak{}officer}\newline\newline Параметр: Визначено період часу, протягом якого персонал або ролі отримують сповіщення про збої в процесі обробки даних аудиту\newline Тип: list\newline Значення: \textbf{admin, security\_\allowbreak{}officer}\newline\newline Параметр: Визначено додаткові дії, яких слід вжити у випадку збою в процесі обробки даних аудиту \newline Тип: list\newline Значення: \textbf{login, logout, failed\_\allowbreak{}attempt} & AU-5 & a. Сповіщати [Призначення: визначені організацією персонал або посади] у разі збою обробки даних аудиту в [Призначення: визначений організацією період часу].\\ 
b. Виконати наступні додаткові дії: [Призначення: визначені організацією дії, які необхідно зробити]. \\ \hline
49 & ОГЛЯД, АНАЛІЗ І ЗВІТНІСТЬ АУДИТУ (AU-6) & Параметр: Записи аудиту системи переглядаються та аналізуються частота для виявлення ознак неналежної або незвичної діяльності та потенційного впливу неналежної або незвичної діяльності\newline Тип: string\newline Значення: \textbf{щотижня}\newline\newline Параметр: Звіт аудиту відправляється персоналу або ролям\newline Тип: list\newline Значення: \textbf{admin, security\_\allowbreak{}officer}\newline\newline Параметр: Визначено частоту, з якою переглядаються та аналізуються записи аудиту системи\newline Тип: integer\newline Значення: \textbf{30}\newline\newline Параметр: Визначено персонал або ролі які отримують результати оглядів та аналізів системних записів\newline Тип: list\newline Значення: \textbf{admin, security\_\allowbreak{}officer} & AU-6 & a. Переглядати та аналізувати записи системного аудиту [Призначення: з визначеною організацією частотою] для виявлення [Призначення: визначеної організацією неналежної або незвичайної діяльності].\\ 
b. Відправляти звіт про аудит [Призначення: визначеним організацією персоналу або посадам].\\ 
c. Налаштувати рівні огляду аудиту, аналізу та звітності в рамках системи, коли змінюється рівень ризику на основі інформації від правоохоронних органів, розвідувальної інформації або від інших достовірних джерел інформації. \\ \hline
50 & ОГЛЯД, АНАЛІЗ І ЗВІТНІСТЬ АУДИТУ - ЗІСТАВЛЯННЯ СХОВИЩ АУДИТУ (AU-6(3)) &  & AU-6(3) & Аналізувати та зіставляти записи аудиту в різних сховищах задля забезпечення ситуативної обізнаності в масштабах організації. \\ \hline
51 & СКОРОЧЕННЯ ЗАПИСІВ АУДИТУ ТА ФОРМУВАННЯ ЗВІТУ (AU-7) & Параметр: Забезпечено можливість скорочення записів перевірок аудитом та звітів, які не змінюють оригінальний зміст або час упорядкування записів аудиту\newline Тип: integer\newline Значення: \textbf{30}\newline\newline Параметр: Реалізовано можливість скорочення записів перевірок аудитом та звітів, які не змінюють оригінальний зміст або час упорядкування записів аудиту\newline Тип: integer\newline Значення: \textbf{30} & AU-7 & Забезпечити та реалізувати можливості скорочення записів перевірок аудитом і звітів, до рівня, який:\\ 
a. підтримує перевірку, аналіз і звітність аудиту на вимогу та розслідування (постфактум) інцидентів безпеки;\\ 
b. не змінює оригінальний вміст або час упорядкування записів аудиту. \\ \hline
52 & ПОЗНАЧКА ЧАСУ (AU-8) & Параметр: Внутрішній системний годинник використовується для створення позначок часу для записів аудиту\newline Тип: integer\newline Значення: \textbf{30}\newline\newline Параметр: Позначки часу застосовуються для записів аудиту, які відповідають деталізація вимірювання часу і які використовують всесвітній координований час, мають фіксоване місцеве зміщення місцевого часу від всесвітнього координованого часу або включають зміщення місцевого часу як частину позначки часу\newline Тип: integer\newline Значення: \textbf{30}\newline\newline Параметр: Визначено деталізацію вимірювання часу для часових позначок записів аудиту\newline Тип: integer\newline Значення: \textbf{30} & AU-8 & a. Використовувати внутрішньосистемний годинник для створення позначок часу для записів аудиту.\\ 
b. Застосовувати позначки часу, які відповідають [Призначення: деталізація вимірювання часу, визначена організацією] і використовують всесвітній координований час, мають фіксоване зміщення місцевого часу відносно всесвітнього координованого часу або включають зміщення місцевого часу як частину позначки часу. \\ \hline
53 & ЗАХИСТ ІНФОРМАЦІЇ АУДИТУ (AU-9) & Параметр: Персонал або ролі отримують сповіщення при виявленні несанкціонованого доступу, зміни або видалення інформації аудиту\newline Тип: list\newline Значення: \textbf{admin, security\_\allowbreak{}officer}\newline\newline Параметр: Визначено персонал або ролі, які мають бути сповіщені при виявленні несанкціонованого доступу, зміни або видалення інформації аудиту\newline Тип: list\newline Значення: \textbf{admin, security\_\allowbreak{}officer} & AU-9 & a. Захист інформації аудиту та інструментів несанкціонованого доступу, зміни та видалення; журналювання аудиту від\\ 
b. Сповіщення [Призначення: персонал або ролі, визначені організацією] у разі виявлення несанкціонованого доступу, зміни або видалення інформації аудиту. \\ \hline
54 & ЗАХИСТ ІНФОРМАЦІЇ АУДИТУ - ДОСТУП, ЯКИЙ НАДАЄТЬСЯ ЧЕРЕЗ ЧЛЕНСТВО В ПІДМНОЖИНИ ПРИВІЛЕЙОВАНИХ КОРИСТУВАЧІВ (AU-9(4)) &  & AU-9(4) & Авторизувати доступ до управління функціональністю аудиту тільки для [Призначення: визначеної організацією підмножини привілейованих користувачів]. \\ \hline
55 & НЕСПРОСТОВНІСТЬ (AU-10) & Параметр: Надаються неспростовні докази того, що особа (або процес, що діє від імені особи) виконала дії\newline Тип: list\newline Значення: \textbf{admin, security\_\allowbreak{}officer}\newline\newline Параметр: Визначено дії, на які поширюється принцип неспростовності\newline Тип: list\newline Значення: \textbf{login, logout, failed\_\allowbreak{}attempt} & AU-10 & Надавайте неспростовні докази того, що особа (або процес, який діє від імені особи) виконала [Призначення: дії, визначені організацією, на які поширюється принцип неспростовності]. \\ \hline
56 & НЕСПРОСТОВНІСТЬ - ЦИФРОВІ ПІДПИСИ (AU-10(5)) &  & AU-10(5) &  \\ \hline
57 & ЗБЕРЕЖЕННЯ ЗАПИСІВ АУДИТУ (AU-11) & Параметр: Записи аудиту зберігаються впродовж період часу, щоб забезпечити підтримку розслідування (постфактум) інцидентів безпеки та конфіденційності, а також відповідати нормативним та вимогам організації щодо збереження даних аудиту\newline Тип: integer\newline Значення: \textbf{30}\newline\newline Параметр: Визначено період часу для зберігання записів аудиту, який узгоджується з політикою зберігання записів\newline Тип: list\newline Значення: \textbf{default\_\allowbreak{}deny\_\allowbreak{}rule, abac\_\allowbreak{}rule\_\allowbreak{}1} & AU-11 & Зберігати записи аудиту впродовж [Призначення: визначеного організацією періоду часу, відповідно політиці зберігання записів], щоб забезпечити підтримку розслідувань (постфактум) інцидентів безпеки та приватності, а також для задоволення вимог нормативних і документів організації щодо збереження даних аудиту. \\ \hline
58 & ГЕНЕРАЦІЯ ДАНИХ АУДИТУ (AU-12) & Параметр: Персонал або ролі може/можуть вибирати типи подій, які будуть реєструватися певними компонентами системи; AU-12(c) згенеровано записи аудиту для типів подій, визначених у AU02\_\allowbreak{}ODP[02], які включають вміст записів аудиту, визначений у AU03\newline Тип: list\newline Значення: \textbf{admin, security\_\allowbreak{}officer}\newline\newline Параметр: Визначено персонал або ролі, яким дозволено обирати типи подій, що мають реєструватися певними компонентами системи\newline Тип: list\newline Значення: \textbf{admin, security\_\allowbreak{}officer} & AU-12 & a. Забезпечити генерацію даних аудиту для типів подій, що перевіряються в AU-2а в [Призначення: визначених організацією компонентах системи].\\ 
b. Дозволити [Призначення: визначеному організацією персоналу або посадам] вибирати, які типи подій, що перевіряються, повинні перевірятися окремими компонентами системи;\\ 
c. Генерувати записи аудиту для типів подій, визначених в AU-2с. з вмістом згідно з AU-3. \\ \hline
\multicolumn{5}{|l|}{\textbf{4. ОЦІНЮВАННЯ, АВТОРИЗАЦІЯ ТА МОНІТОРИНГ (CA) (CA)}} \\ \hline
59 & ПОЛІТИКА І ПРОЦЕДУРИ ОЦІНЮВАННЯ, АКРЕДИТАЦІЯ ТА МОНІТОРИНГ БЕЗПЕКИ (CA-1) & Параметр: Розроблено та задокументовано політику оцінювання, авторизації та моніторингу\newline Тип: list\newline Значення: \textbf{default\_\allowbreak{}deny\_\allowbreak{}rule, abac\_\allowbreak{}rule\_\allowbreak{}1}\newline\newline Параметр: Політика оцінювання, авторизації та моніторингу поширюється серед персоналу або ролей\newline Тип: list\newline Значення: \textbf{admin, security\_\allowbreak{}officer}\newline\newline Параметр: Розроблені та задокументовані процедури оцінювання, авторизації та моніторингу, що сприяють впровадженню політики оцінювання, авторизації та моніторингу, а також пов'язані з ними засоби контролю оцінювання, авторизації та моніторингу\newline Тип: list\newline Значення: \textbf{default\_\allowbreak{}deny\_\allowbreak{}rule, abac\_\allowbreak{}rule\_\allowbreak{}1}\newline\newline Параметр: Посадова особа призначається для управління розробкою, документуванням та розповсюдженням політики та процедур оцінювання, авторизації та моніторингу; CA-01(c)[01][01] переглядається та оновлюється поточна політика оцінки, авторизації та моніторингу з частотою; CA-01(c)[01][02] переглядається та оновлюється поточна політика оцінки, авторизації та моніторингу після подій; CA-01(c)[02][01] переглядаються та оновлюються поточні процедури оцінки, авторизації та моніторингу з частотою; CA-01(c)[02][02] переглядаються та оновлюються поточні процедури оцінки, авторизації та моніторингу після подій; ПОТЕНЦІЙНІ МЕТОДИ ТА ОБ'ЄКТИ ОЦІНКИ:\newline Тип: list\newline Значення: \textbf{admin, security\_\allowbreak{}officer}\newline\newline Параметр: Визначено персонал або ролі, серед яких має бути поширена політика оцінювання, авторизації та моніторингу\newline Тип: list\newline Значення: \textbf{admin, security\_\allowbreak{}officer}\newline\newline Параметр: Визначено персонал або ролі, серед яких мають бути поширені процедури оцінювання, авторизації та моніторингу\newline Тип: list\newline Значення: \textbf{admin, security\_\allowbreak{}officer}\newline\newline Параметр: Визначено посадову особу, яка управлятиме розробкою, документуванням і розповсюдженням політики та процедур оцінювання, авторизації та моніторингу\newline Тип: list\newline Значення: \textbf{admin, security\_\allowbreak{}officer}\newline\newline Параметр: Визначено частоту, з якою переглядається та оновлюється поточна політика оцінювання, авторизації та моніторингу\newline Тип: list\newline Значення: \textbf{default\_\allowbreak{}deny\_\allowbreak{}rule, abac\_\allowbreak{}rule\_\allowbreak{}1}\newline\newline Параметр: Визначено події, які потребують перегляду та оновлення поточної політики оцінки, авторизації та моніторингу\newline Тип: list\newline Значення: \textbf{default\_\allowbreak{}deny\_\allowbreak{}rule, abac\_\allowbreak{}rule\_\allowbreak{}1}\newline\newline Параметр: Визначено частоту, з якою переглядається та оновлюється поточні процедури оцінювання, авторизації та моніторингу\newline Тип: integer\newline Значення: \textbf{30}\newline\newline Параметр: Визначено події, які потребують перегляду та оновлення поточні процедури оцінки, авторизації та моніторингу\newline Тип: list\newline Значення: \textbf{login, logout, failed\_\allowbreak{}attempt} & CA-1 & a. Розробити, задокументувати та поширити серед [Призначення: визначеного організацією персоналу або посад]:\\ 
1. [Вибір (один або декілька): Рівень організації; Рівень місії/бізнес-процесу; рівень системи] політика оцінювання, авторизації та моніторингу, яка:\\ 
(a) містить мету, сферу застосування, ролі, обов'язки, відповідальність керівництва, координацію між організаційними підрозділами та систему контролю відповідності (complaince);\\ 
(b) відповідає чинним законам, нормативним документам, наказам, положенням, політиці, стандартам і керівним принципам.\\ 
2. Процедури, що сприяють реалізації політики оцінювання, авторизації та моніторингу безпеки та приватності, а також пов'язаних з ними заходів оцінювання, авторизації та моніторингу безпеки та приватності.\\ 
b. Призначити [Призначення: посадова особа, визначена організацією] для управління розробкою, документуванням і розповсюдженням політики та процедур оцінювання, авторизації та моніторингу;\\ 
c. Переглядати та оновлювати поточне оцінювання, авторизацію та моніторинг:\\ 
1. Політику [Призначення: частота, визначена організацією] та наступне [Призначення: події, визначені організацією];\\ 
2. Процедури [Призначення: частота, визначена організацією] та наступні [Призначення: події, визначені організацією]. \\ \hline
60 & ОЦІНЮВАННЯ (CA-2) & Параметр: Розроблено план контрольної оцінки, який описує обсяг оцінки, включаючи заходи захисту та посилені заходи, що підлягають оцінюванню. CA-02(b)[02] розроблено план контрольної оцінки, який описує обсяг оцінки, включаючи процедури оцінювання, які використовуватимуться для визначення ефективності заходів. CA-02(b)[03][01] розроблено план контрольної оцінки, який описує обсяг оцінки, включаючи середовище оцінювання. CA-02(b)[03][02] розроблено план контрольної оцінки, який описує обсяг оцінки, включаючи групу оцінювання. CA-02(b)[03][03] розроблено план контрольної оцінки, який описує обсяг оцінки, включаючи ролі й обов'язки з оцінювання\newline Тип: list\newline Значення: \textbf{admin, security\_\allowbreak{}officer}\newline\newline Параметр: Розроблено план контрольної оцінки, який описує обсяг оцінки, включаючи процедури оцінювання, які використовуватимуться для визначення ефективності заходів. CA-02(b)[03][01] розроблено план контрольної оцінки, який описує обсяг оцінки, включаючи середовище оцінювання. CA-02(b)[03][02] розроблено план контрольної оцінки, який описує обсяг оцінки, включаючи групу оцінювання. CA-02(b)[03][03] розроблено план контрольної оцінки, який описує обсяг оцінки, включаючи ролі й обов'язки з оцінювання\newline Тип: list\newline Значення: \textbf{admin, security\_\allowbreak{}officer}\newline\newline Параметр: План оцінки заходів захисту розглядається та затверджується уповноваженою посадовою особою або призначеним представником перед проведенням оцінки\newline Тип: list\newline Значення: \textbf{admin, security\_\allowbreak{}officer}\newline\newline Параметр: Заходи захисту оцінюються в системі та середовищі її функціонування частота оцінки, щоб визначити, наскільки коректно реалізовані заходи захисту, чи працюють вони за призначенням і чи дають бажаний результат щодо дотримання встановлених вимог до безпеки\newline Тип: string\newline Значення: \textbf{щорічно}\newline\newline Параметр: Заходи захисту оцінюються в системі та середовищі її функціонування частота оцінки, щоб визначити, наскільки коректно реалізовані заходи захисту, чи працюють вони за призначенням і чи дають бажаний результат щодо дотримання встановлених вимог до конфіденційності; CA-02(e) готується звіт оцінювання\newline Тип: string\newline Значення: \textbf{щорічно}\newline\newline Параметр: Результати оцінювання з безпеки надаються особам або ролям. оцінювання , який документує результати\newline Тип: list\newline Значення: \textbf{admin, security\_\allowbreak{}officer}\newline\newline Параметр: Визначено частоту, з якою слід оцінювати засоби контролю в системі та середовищі її функціонування\newline Тип: integer\newline Значення: \textbf{30}\newline\newline Параметр: Визначені особи або ролі, яким мають бути надані результати оцінювання з безпеки\newline Тип: list\newline Значення: \textbf{admin, security\_\allowbreak{}officer} & CA-2 & a. Виберіть відповідного оцінювача або команду з оцінки для типу оцінювання, яке буде проводитися;\\ 
b. Розробіть план контрольної оцінки, який описує обсяг оцінки, в тому числі:\\ 
1. заходи захисту та посилені заходи, що підлягають оцінюванню;\\ 
2. процедури оцінювання, ефективності заходів; які використовуватимуться для визначення\\ 
3. середовище оцінювання, групу оцінювання, ролі й обов'язки з оцінювання.\\ 
c. Забезпечити розгляд і затвердження плану оцінювання уповноваженою офіційною особою або призначеним для проведення оцінювання представником;\\ 
d. Оцінити заходи захисту в системі та в її середовищі функціонування з [Призначення: визначеною організацією частотою] для визначення, наскільки коректно реалізовані заходи безпеки і чи працюють вони за призначенням і дають бажаний результат щодо дотримання встановлених вимог безпеки та приватності;\\ 
e. Підготовити звіт оцінювання безпеки, який документує результати оцінювання;\\ 
f. Надати результати оцінювання з безпеки [Призначення: особам або ролям, визначеним організацією]. \\ \hline
61 & ВЗАЄМОДІЯ СИСТЕМ (CA-3) & Параметр: Характеристики інтерфейсу задокументовані як частина кожної угоди про обмін\newline Тип: integer\newline Значення: \textbf{30}\newline\newline Параметр: Вимоги до безпеки задокументовані як частина кожної угоди про обмін\newline Тип: integer\newline Значення: \textbf{30}\newline\newline Параметр: Вимоги щодо конфіденційності задокументовані як частина кожної угоди про обмін\newline Тип: integer\newline Значення: \textbf{30}\newline\newline Параметр: Заходи захисту задокументовані як частина кожної угоди про обмін\newline Тип: integer\newline Значення: \textbf{30}\newline\newline Параметр: Відповідальність за кожну систему задокументована як частина кожної угоди про обмін\newline Тип: integer\newline Значення: \textbf{30}\newline\newline Параметр: Характер переданої інформації документується як частина кожної угоди про обмін\newline Тип: integer\newline Значення: \textbf{30}\newline\newline Параметр: Угоди переглядаються та оновлюються частота\newline Тип: string\newline Значення: \textbf{щорічно}\newline\newline Параметр: Визначено частоту, з якою необхідно переглядати та оновлювати угоди\newline Тип: integer\newline Значення: \textbf{30} & CA-3 & a. схвалити та керувати обміном інформацією між системою та іншими системами за допомогою [Вибір (один або кілька): угоди безпеки взаємозв'язку; договори безпеки обміну інформацією; меморандуми про взаєморозуміння; угоди про рівень обслуговування; угоди користувача; угоди про нерозголошення; [Доручення: тип договору, визначений організацією]];.\\ 
b. документувати, як частину угоди про обмін, характеристики інтерфейсу, вимоги до безпеки та приватності, засоби контролю та відповідальність для кожної системи, а також характер переданої інформації;\\ 
c. здійснювати перегляд та оновлення угод з [Призначення: визначеною організацією частотою]. \\ \hline
62 & ПЛАН УСУНЕННЯ НЕДОЛІКІВ ТА КОНТРОЛЬНІ ПОКАЗНИКИ (CA-5) & Параметр: Розроблено план усунення недоліків та контрольні показники для системи, та задокументувано заплановані дії організації з коригування, спрямовані на усунення недоліків та зауважень, виявлених під час оцінки заходів захисту, а також на зменшення або усунення відомих вразливостей в системі\newline Тип: list\newline Значення: \textbf{login, logout, failed\_\allowbreak{}attempt}\newline\newline Параметр: Існуючий план оновлюються частота на основі результатів оцінювання заходів, незалежних аудитів та постійного моніторингу\newline Тип: string\newline Значення: \textbf{щорічно}\newline\newline Параметр: Визначено частоту оновлення чинного плану усунення недоліків та контрольних показників на основі результатів оцінювання заходів захисту, незалежних аудитів та постійного моніторингу\newline Тип: integer\newline Значення: \textbf{30} & CA-5 & a. Розробити для системи план усунення недоліків та контрольні показники з метою документування запланованих коригувальних дій організації для усунення недоліків і зауважень, які виявлені в ході оцінювання заходів захисту, а також для зменшення або усунення відомих вразливостей у системі.\\ 
b. Оновлювати чинний план усунення недоліків та контрольні показники з [Призначення: визначеною організацією частотою] на основі результатів оцінювання заходів, незалежних аудитів та постійного моніторингу. \\ \hline
63 & БЕЗПЕРЕРВНИЙ МОНІТОРИНГ (CA-7) & Параметр: Безперервний моніторинг на рівні системи включає встановлені частоти для моніторингу ефективності заходів захисту\newline Тип: integer\newline Значення: \textbf{30}\newline\newline Параметр: Безперервний моніторинг на рівні системи включає встановлені частоти для оцінки ефективності заходів захисту\newline Тип: integer\newline Значення: \textbf{30}\newline\newline Параметр: Безперервний моніторинг на рівні системи включає в себе дії з реагування на результати аналізу інформації, пов'язаної з безпекою та приватністю\newline Тип: list\newline Значення: \textbf{login, logout, failed\_\allowbreak{}attempt}\newline\newline Параметр: Безперервний моніторинг на рівні системи включає повідомлення про статус безпеки системи для персоналу або ролей частота\newline Тип: list\newline Значення: \textbf{admin, security\_\allowbreak{}officer}\newline\newline Параметр: Безперервний моніторинг на рівні системи включає повідомлення про стан конфіденційності системи персоналу або ролям частота\newline Тип: list\newline Значення: \textbf{admin, security\_\allowbreak{}officer}\newline\newline Параметр: Визначено частоту, з якою слід моніторити ефективність заходів захисту\newline Тип: integer\newline Значення: \textbf{30}\newline\newline Параметр: Визначено частоту, з якою слід оцінювати ефективність заходів захисту\newline Тип: integer\newline Значення: \textbf{30}\newline\newline Параметр: Визначено персонал або ролі, яким повідомляється про стан безпеки системи\newline Тип: list\newline Значення: \textbf{admin, security\_\allowbreak{}officer}\newline\newline Параметр: Визначено частоту, з якою повідомляється про стан безпеки системи\newline Тип: integer\newline Значення: \textbf{30}\newline\newline Параметр: Визначено персонал або ролі, яким повідомляється про стан конфіденційності системи\newline Тип: list\newline Значення: \textbf{admin, security\_\allowbreak{}officer}\newline\newline Параметр: Визначено частоту, з якою повідомляється про стан конфіденційності системи\newline Тип: integer\newline Значення: \textbf{30} & CA-7 & Розробити стратегію безперервного моніторингу безпеки та приватності й упровадити програму безперервного моніторингу безпеки та приватності, яка охоплює:\\ 
a. встановлення показників безпеки та приватності, які необхідно відстежувати: [Призначення: визначені організацією метрики];\\ 
b. встановлення [Призначення: визначена організацією частота] для моніторингу та [Призначення: визначена організацією частота] для безперервного оцінювання ефективності заходів захисту;\\ 
c. поточні оцінювання заходів захисту відповідно до стратегії безперервного моніторингу організації;\\ 
d. постійний моніторинг стану безпеки та приватності відповідно до встановлених організацією метрик і відповідно до стратегії безперервного моніторингу організації;\\ 
e. зіставлення та аналіз інформації, отриманої в результаті оцінювання та моніторингу безпеки та приватності;\\ 
f. дії реагування за результатами аналізу інформації, пов'язаної з безпекою та приватністю;\\ 
g. повідомлення про статус безпеки та приватності системи [Призначення: визначені організацією персонал або ролі] з [Призначення: визначеною організацією частотою]. \\ \hline
\multicolumn{5}{|l|}{\textbf{5. УПРАВЛІННЯ КОНФІГУРАЦІЄЮ (CM) (CM)}} \\ \hline
64 & Політика та процедури управління конфігурацією (CM-1) & Параметр: Розроблено та задокументовано політику управління конфігурацією\newline Тип: list\newline Значення: \textbf{default\_\allowbreak{}deny\_\allowbreak{}rule, abac\_\allowbreak{}rule\_\allowbreak{}1}\newline\newline Параметр: Політика управління конфігурацією поширюється на персонал або ролі\newline Тип: list\newline Значення: \textbf{admin, security\_\allowbreak{}officer}\newline\newline Параметр: Розроблені та задокументовані процедури управління , що сприяють реалізації політики управління конфігурацією та пов'язаних з нею заходів управління конфігурацією\newline Тип: list\newline Значення: \textbf{default\_\allowbreak{}deny\_\allowbreak{}rule, abac\_\allowbreak{}rule\_\allowbreak{}1}\newline\newline Параметр: Посадова особа призначається для управління розробкою, документуванням і розповсюдженням політики та процедур керування конфігурацією; CM-01(c)[01][01] переглядається та оновлюється поточна політика управління конфігурацією частота; CM-01(c)[01][02] переглядається та оновлюється поточна політика управління конфігурацією після події; CM-01(c)[02][01] переглядаються та оновлюються поточні процедури управління конфігурацією частота; CM-01(c)[02][02] переглядаються та оновлюються поточні процедури управління конфігурацією після події; ЗНАЧЕННЯ конфігурацією\newline Тип: list\newline Значення: \textbf{admin, security\_\allowbreak{}officer}\newline\newline Параметр: Визначено персонал або ролі, на яких поширюється політика управління конфігурацією\newline Тип: list\newline Значення: \textbf{admin, security\_\allowbreak{}officer}\newline\newline Параметр: Визначено персонал або ролі, на яких поширюється процедури управління конфігурацією\newline Тип: list\newline Значення: \textbf{admin, security\_\allowbreak{}officer}\newline\newline Параметр: Визначено посадову особу, яка управляє розробкою, документуванням і розповсюдженням політики та процедур керування конфігурацією\newline Тип: list\newline Значення: \textbf{admin, security\_\allowbreak{}officer}\newline\newline Параметр: Визначено частоту, з якою переглядається та оновлюється поточна політика управління конфігурацією\newline Тип: list\newline Значення: \textbf{default\_\allowbreak{}deny\_\allowbreak{}rule, abac\_\allowbreak{}rule\_\allowbreak{}1}\newline\newline Параметр: Визначено події, після яких переглядається та оновлюється поточна політика управління конфігурацією\newline Тип: list\newline Значення: \textbf{default\_\allowbreak{}deny\_\allowbreak{}rule, abac\_\allowbreak{}rule\_\allowbreak{}1}\newline\newline Параметр: Визначено частоту, з якою переглядаються та оновлюються поточні процедури управління конфігурацією\newline Тип: integer\newline Значення: \textbf{30}\newline\newline Параметр: Визначено події, після яких переглядаються та оновлюються поточні процедури управління конфігурацією\newline Тип: list\newline Значення: \textbf{login, logout, failed\_\allowbreak{}attempt} & CM-1 & a. Розробити, задокументувати та поширити серед [Призначення: визначених організацією персоналу або ролей]:\\ 
1. [Вибір (один або декілька): Рівень організації; Рівень місії/бізнес-процесу; рівень системи] політики управління конфігурацією, яка:\\ 
2.\\ 
(a) містить мету, сферу застосування, ролі, обов'язки, відповідальність керівництва, координацію між організаційними підрозділами та систему контролю відповідності (complaince);\\ 
(b) відповідає чинним законам, нормативним документам, наказам, положенням, політикам, стандартам і керівним принципам; процедури, що сприяють реалізації політики управління конфігурацією та пов'язаних з нею заходів управління конфігурацією.\\ 
b. Призначити [Призначення: посадова особа, визначена організацією] для управління розробкою, документуванням і розповсюдженням політики та процедур керування конфігурацією.\\ 
c. Переглядати та оновлювати поточну політику управління конфігурацією:\\ 
1. Політика [Призначення: частота, визначена організацією] та наступні [Призначення: події, визначені організацією];\\ 
2. Процедури [Призначення: частота, визначена організацією] та наступні [Призначення: події, визначені організацією]. \\ \hline
65 & БАЗОВА КОНФІГУРАЦІЯ (CM-2) & Параметр: Переглядаються та оновлюються базові налаштування системи частота\newline Тип: string\newline Значення: \textbf{щорічно}\newline\newline Параметр: Визначено частоту налаштувань; перегляду та оновлення базових\newline Тип: integer\newline Значення: \textbf{30} & CM-2 & a. Розробити, задокументувати та підтримувати за допомогою заходів конфігурації поточні базові налаштування системи.\\ 
b. Переглядати та оновлювати базові налаштування системи:\\ 
1. з [Призначення: визначеною організацією частотою];\\ 
2. за потреби внаслідок [Призначення: визначених організацією обставин];\\ 
3. коли встановлені нові або оновлені компоненти системи. \\ \hline
66 & БАЗОВА КОНФІГУРАЦІЯ - КОНФІГУРАЦІЯ СИСТЕМ ТА КОМПОНЕНТІВ ДЛЯ СФЕР З ВИСОКИМ РИЗИКОМ (CM-2(7)) & Параметр: Системи або компоненти системи з конфігураціями видаються особам, що перебувають у місцях, які організація вважає, становлять значний ризик\newline Тип: list\newline Значення: \textbf{admin, security\_\allowbreak{}officer}\newline\newline Параметр: Заходи безпеки застосовуються до систем або компонентів системи, коли особи повертаються з поїздки\newline Тип: list\newline Значення: \textbf{admin, security\_\allowbreak{}officer}\newline\newline Параметр: Визначено системи або компоненти систем, які мають видаватися особам, що перебувають у місцях зі значним ризиком\newline Тип: list\newline Значення: \textbf{admin, security\_\allowbreak{}officer} & CM-2(7) & (a) Видавати [Призначення: визначених організацією систем або компонентів систем] з [Призначенням: визначеними організацією конфігураціями] особам, що перебувають у місцях, які організація вважає місцями зі значним ризиком;\\ 
(b) Застосувати [Призначення: визначені організацією запобіжні заходи безпеки] до компонентів, коли особи повертаються з поїздки. \\ \hline
67 & УПРАВЛІННЯ ЗМІНАМИ КОНФІГУРАЦІЇ (CM-3) & Параметр: Записи про зміни конфігурації у системі зберігаються протягом <період часу CM-03\_\allowbreak{}ODP[01]>\newline Тип: integer\newline Значення: \textbf{30}\newline\newline Параметр: Визначено період часу, протягом якого зберігатимуться записи про зміни конфігурації\newline Тип: integer\newline Значення: \textbf{30}\newline\newline Параметр: Вибрано одне або декілька з наступних ЗНАЧЕНЬ ПАРАМЕТРІВ: \{частота; коли умови\}\newline Тип: list\newline Значення: \textbf{}\newline\newline Параметр: Визначено частоту, з якою викликаються елементи управління змінами конфігурації (якщо вибрано)\newline Тип: integer\newline Значення: \textbf{30}\newline\newline Параметр: Визначено умови, за яких викликаються елементи управління змінами конфігурації (якщо вибрано); CM-03(a) визначено та задокументовано типи змін до системи, які контролюються конфігурацією\newline Тип: list\newline Значення: \textbf{} & CM-3 & a. Визначити типи змін у системі, які контролюються конфігурацією.\\ 
b. Переглядати запропоновані зміни в конфігурації, контрольовані системою, і схвалити або відхилити ці зміни з явним урахуванням аналізу наслідків безпеки.\\ 
c. Документувати рішення зі зміни конфігурації системи.\\ 
d. Впровадити схвалені зміни конфігурації в систему.\\ 
e. Зберігати записи змін конфігурації системі впродовж [Призначення: певного періоду часу, визначеного організацією].\\ 
f. Здійснювати моніторинг і аналіз дій, пов'язаних зі змінами конфігурації системи.\\ 
g. Координувати та впроваджувати нагляд за діяльністю з управління змінами конфігурації за допомогою [Призначення: елементу управління змінами конфігурації, визначеного організацією], який викликається [Вибір (один або кілька): [Призначення: з визначеною організацією частотою]; [Призначення: визначені організацією умови зміни конфігурації]]. \\ \hline
68 & АНАЛІЗ ВПЛИВУ НА БЕЗПЕКУ ТА ПРИВАТНІСТЬ (CM-4) &  & CM-4 & Аналізувати зміни в системі, щоб визначити потенційну загрозу безпеці та приватності перед реалізацією змін. \\ \hline
69 & АНАЛІЗ ВПЛИВУ НА БЕЗПЕКУ ТА ПРИВАТНІСТЬ - ВЕРИФІКАЦІЯ ФУНКЦІЙ БЕЗПЕКИ ТА ПРИВАТНОСТІ (CM-4(2)) &  & CM-4(2) & Після змін у системі переконайтеся, що відповідні заходи захисту реалізовано правильно і вони функціонують належним чином та дають бажаний результат щодо дотримання вимог безпеки та приватності для системи. \\ \hline
70 & ОБМЕЖЕННЯ ДОСТУПУ ДО ЗМІНИ (CM-5) &  & CM-5 & Визначити, задокументувати, затвердити та забезпечити застосування фізичних і логічних обмежень доступу, пов'язаних зі змінами в системі. \\ \hline
71 & НАЛАШТУВАННЯ КОНФІГУРАЦІЇ (CM-6) & Параметр: Зміни в налаштуваннях конфігурації відстежуються відповідно до політики та процедур організації\newline Тип: list\newline Значення: \textbf{default\_\allowbreak{}deny\_\allowbreak{}rule, abac\_\allowbreak{}rule\_\allowbreak{}1}\newline\newline Параметр: Зміни налаштувань конфігурації керуються відповідно до політики та процедур організації\newline Тип: list\newline Значення: \textbf{default\_\allowbreak{}deny\_\allowbreak{}rule, abac\_\allowbreak{}rule\_\allowbreak{}1} & CM-6 & a. Встановити та задокументувати параметри конфігурації компонентів, які застосовуються в системі, які відображають найбільш обмежений режим, що відповідає експлуатаційним вимогам, використовуючи [Призначення: визначені організацією загальні безпечні конфігурації].\\ 
b. Реалізувати конфігураційні установки.\\ 
c. Визначити, задокументувати та затвердити будь-які відхилення від встановлених конфігураційних параметрів конфігурації для [Призначення: визначених організацією компонентів системи] на основі [Призначення: визначених організацією експлуатаційних вимог].\\ 
d. Відстежувати та керувати змінами конфігураційних параметрів відповідно до організаційної політики та процедур. \\ \hline
72 & МІНІМАЛЬНО НЕОБХІДНА ФУНКЦІОНАЛЬНІСТЬ (CM-7) & Параметр: Використання протоколи заборонено або обмежено\newline Тип: string\newline Значення: \textbf{TLS 1.3}\newline\newline Параметр: Визначено протоколи, які необхідно заборонити або обмежити; CM-07\_\allowbreak{}ODP[05] визначено програмне забезпечення, яке необхідно заборонити або обмежити\newline Тип: string\newline Значення: \textbf{TLS 1.3} & CM-7 & a. Налаштуйте систему для забезпечення лише [Призначення: основні функції, визначені організацією для місії];\\ 
b. Заборонити або обмежити використання таких функцій, портів, протоколів, програмного забезпечення та/або служб: [Призначення: визначені організацією заборонені або обмежені функції, системні порти, протоколи, програмне забезпечення та/або служби]. \\ \hline
73 & МІНІМАЛЬНО НЕОБХІДНА ФУНКЦІОНАЛЬНІСТЬ - ПЕРІОДИЧНИЙ ПЕРЕГЛЯД (CM-7(1)) & Параметр: Система переглядається частота для виявлення непотрібних та/або незахищених функцій, портів, протоколів і послуг\newline Тип: string\newline Значення: \textbf{TLS 1.3}\newline\newline Параметр: Протоколи, які вважаються непотрібними та/або незахищеними, вимкнено\newline Тип: string\newline Значення: \textbf{TLS 1.3}\newline\newline Параметр: Визначено частоту, з якою слід переглядати систему для виявлення непотрібних та/або незахищених функцій, портів, протоколів і послуг\newline Тип: string\newline Значення: \textbf{TLS 1.3}\newline\newline Параметр: Визначено протоколи, які слід вимкнути, якщо вони вважаються непотрібними або незахищеними\newline Тип: string\newline Значення: \textbf{TLS 1.3} & CM-7(1) &  \\ \hline
74 & МІНІМАЛЬНО НЕОБХІДНА ФУНКЦІОНАЛЬНІСТЬ - АВТОРИЗОВАНЕ ПРОГРАМНЕ ЗАБЕЗПЕЧЕННЯ -- БІЛИЙ СПИСОК (CM-7(5)) & Параметр: Політика \textquotedbl{}заборонити все, дозволити за винятком\textquotedbl{} застосовуэться, щоб дозволити виконання авторизованих програм у системі\newline Тип: list\newline Значення: \textbf{default\_\allowbreak{}deny\_\allowbreak{}rule, abac\_\allowbreak{}rule\_\allowbreak{}1}\newline\newline Параметр: Переглядається та оновлюється список авторизованих програм частота\newline Тип: list\newline Значення: \textbf{}\newline\newline Параметр: Визначено частоту перегляду та оновлення списку авторизованих програм\newline Тип: integer\newline Значення: \textbf{30} & CM-7(5) & (a) Визначити [Призначення: визначені організацією програмне забезпечення, яке авторизовано виконується в системі]\\ 
(b) Впровадити політику «заборони всього, за винятком деяких», щоб дозволити виконання авторизованих програм у системі.\\ 
(c) Переглядати та оновлювати список авторизованих програм [Призначення: з визначеною організацією частотою]. \\ \hline
75 & ІНВЕНТАРИЗАЦІЯ КОМПОНЕНТІВ СИСТЕМИ (CM-8) & Параметр: Переглядається та оновлюється опис компонентів системи частота\newline Тип: string\newline Значення: \textbf{щорічно}\newline\newline Параметр: Визначено частоту перегляду та оновлення опису компонентів системи\newline Тип: integer\newline Значення: \textbf{30} & CM-8 & a.\\ 
b. Розробити та задокументувати процес інвентаризації компонентів системи, який:\\ 
1. точно описує поточну систему;\\ 
2. охоплює всі компоненти в межах акредитації системи;\\ 
3. не включає повторний облік компонентів або компонентів, будь-якої іншої системи;\\ 
4. визначає рівень деталізації, який є необхідним для відстеження та звітування;\\ 
5. включає інформацію для досягнення підзвітності компонентів системи: [Призначення: визначена організацією інформація, необхідна для досягнення ефективної підзвітності компонентів системи]. Переглядати та оновлювати опис компонентів системи з [Призначення: визначеною організацією частотою]. \\ \hline
76 & ІНВЕНТАРИЗАЦІЯ КОМПОНЕНТІВ СИСТЕМИ - ОНОВЛЕННЯ ПІД ЧАС ВСТАНОВЛЕННЯ ТА ВИДАЛЕННЯ (CM-8(1)) &  & CM-8(1) &  \\ \hline
77 & РОЗТАШУВАННЯ ІНФОРМАЦІЇ (CM-12) &  & CM-12 &  \\ \hline
\multicolumn{5}{|l|}{\textbf{6. ПЛАНУВАННЯ БЕЗПЕРЕРВНОЇ РОБОТИ (CP) (CP)}} \\ \hline
78 & РЕЗЕРВНЕ КОПІЮВАННЯ (CP-9) & Параметр: Резервне копіювання інформації користувача, що міститься в компонентах системи, здійснюється частота\newline Тип: string\newline Значення: \textbf{щорічно}\newline\newline Параметр: Виконується резервне копіювання інформації системи, що міститься в системі частота; CP-09(с) створюються резервні копії документації системи, включаючи документацію, пов'язану з безпекою та конфіденційністю частота\newline Тип: string\newline Значення: \textbf{щорічно}\newline\newline Параметр: Визначено частоту, з якою слід проводити резервне копіювання інформації користувача відповідно до часу відновлення та цілей відновлення\newline Тип: integer\newline Значення: \textbf{30}\newline\newline Параметр: Визначено частоту проведення резервного копіювання інформації системи, що відповідає завдань відновлення і встановлених цілей відновлення\newline Тип: integer\newline Значення: \textbf{30}\newline\newline Параметр: Визначено частоту, з якою слід проводити резервне копіювання документації системи відповідно до часу відновлення та цілей точки відновлення\newline Тип: integer\newline Значення: \textbf{30} & CP-9 &  \\ \hline
79 & РЕЗЕРВНЕ КОПІЮВАННЯ - КРИПТОГРАФІЧНИЙ ЗАХИСТ (CP-9(8)) & Параметр: Реалізовано криптографічні механізми для запобігання несанкціонованому розкриттю та зміні резервної інформації\newline Тип: string\newline Значення: \textbf{AES-256-GCM} & CP-9(8) &  \\ \hline
\multicolumn{5}{|l|}{\textbf{7. ІДЕНТИФІКАЦІЯ ТА АВТЕНТИФІКАЦІЯ (IA) (IA)}} \\ \hline
80 & Політика та процедури ідентифікації та автентифікації (IA-1) & Параметр: Розроблено та задокументовано політику ідентифікації та автентифікації\newline Тип: list\newline Значення: \textbf{default\_\allowbreak{}deny\_\allowbreak{}rule, abac\_\allowbreak{}rule\_\allowbreak{}1}\newline\newline Параметр: Політика ідентифікації та автентифікації поширюється на персонал або ролі\newline Тип: list\newline Значення: \textbf{admin, security\_\allowbreak{}officer}\newline\newline Параметр: Розроблені та задокументовані процедури ідентифікації та автентифікації, що сприяють впровадженню політики ідентифікації та автентифікації, а також відповідні заходи ідентифікації та перевірки автентичності\newline Тип: list\newline Значення: \textbf{default\_\allowbreak{}deny\_\allowbreak{}rule, abac\_\allowbreak{}rule\_\allowbreak{}1}\newline\newline Параметр: Посадова особа призначається для управління політикою та процедурами ідентифікації та автентифікації; IA-01(c)[01][01] переглядається та оновлюється поточна політика ідентифікації та автентифікації частота; IA-01(c)[01][02] переглядається та оновлюється поточна політика ідентифікації та автентифікації після подій; IA-01(c)[02][01] переглядаються та оновлюються поточні процедури ідентифікації та автентифікації частота; IA-01(c)[02][02] переглядаються та оновлюються поточні процедури ідентифікації та автентифікації після подій\newline Тип: list\newline Значення: \textbf{admin, security\_\allowbreak{}officer}\newline\newline Параметр: Визначено персонал або ролі, на які поширюється політика ідентифікації та автентифікації\newline Тип: list\newline Значення: \textbf{admin, security\_\allowbreak{}officer}\newline\newline Параметр: Визначено персонал або ролі, на які поширюються процедури ідентифікації та автентифікації\newline Тип: list\newline Значення: \textbf{admin, security\_\allowbreak{}officer}\newline\newline Параметр: Визначено посадову особу, яка управляє політикою та процедурами ідентифікації та автентифікації\newline Тип: list\newline Значення: \textbf{admin, security\_\allowbreak{}officer}\newline\newline Параметр: Визначено частоту, з якою переглядається та оновлюється поточна політика ідентифікації та автентифікації\newline Тип: list\newline Значення: \textbf{default\_\allowbreak{}deny\_\allowbreak{}rule, abac\_\allowbreak{}rule\_\allowbreak{}1}\newline\newline Параметр: Визначено події, які потребують перегляду та оновлення поточної політики ідентифікації та автентифікації\newline Тип: list\newline Значення: \textbf{default\_\allowbreak{}deny\_\allowbreak{}rule, abac\_\allowbreak{}rule\_\allowbreak{}1}\newline\newline Параметр: Визначено частоту, з якою переглядаються та оновлюються поточні процедури ідентифікації та автентифікації\newline Тип: integer\newline Значення: \textbf{30}\newline\newline Параметр: Визначено події, які потребують перегляду та оновлення процедур ідентифікації та автентифікації\newline Тип: list\newline Значення: \textbf{login, logout, failed\_\allowbreak{}attempt} & IA-1 &  \\ \hline
81 & ІДЕНТИФІКАЦІЯ ТА АВТЕНТИФІКАЦІЯ КОРИСТУВАЧІВ (IA-2) &  & IA-2 &  \\ \hline
82 & ІДЕНТИФІКАЦІЯ ТА АВТЕНТИФІКАЦІЯ (КОРИСТУВАЧІВ ОРГАНІЗАЦІЇ) - БАГАТОФАКТОРНА АВТЕНТИФІКАЦІЯ ПРИВІЛЕЙОВАНИХ ОБЛІКОВИХ ЗАПИСІВ (IA-2(1)) &  & IA-2(1) &  \\ \hline
83 & ІДЕНТИФІКАЦІЯ ТА АВТЕНТИФІКАЦІЯ (КОРИСТУВАЧІВ ОРГАНІЗАЦІЇ) - БАГАТОФАКТОРНА АВТЕНТИФІКАЦІЯ НЕПРИВІЛЕЙОВАНИХ (IA-2(2)) &  & IA-2(2) &  \\ \hline
84 & Ідентифікація та автентифікація (користувачів організації) - Доступ до облікових записів --- стійкість до відтворення (IA-2(8)) &  & IA-2(8) &  \\ \hline
85 & ІДЕНТИФІКАЦІЯ ТА АВТЕНТИФІКАЦІЯ (КОРИСТУВАЧІВ ОРГАНІЗАЦІЇ) - ПРИЙНЯТТЯ ПОВНОВАЖЕНЬ ДЛЯ ВЕРИФІКАЦІЇ ОСОБИСТОЇ ІНФОРМАЦІЇ (PIV CREDENTIALS) (IA-2(12)) & Параметр: Приймаються та електронним шляхом підтверджуються повноваження облікових даних особистої ідентифікації\newline Тип: list\newline Значення: \textbf{admin, security\_\allowbreak{}officer} & IA-2(12) &  \\ \hline
86 & ІДЕНТИФІКАЦІЯ ТА АВТЕНТИФІКАЦІЯ ПРИСТРОЇВ (IA-3) &  & IA-3 &  \\ \hline
87 & УПРАВЛІННЯ ІДЕНТИФІКАЦІЄЮ (IA-4) & Параметр: Управління ідентифікаторами здійснюється шляхом отримання дозволу від персоналу або ролей на призначення ідентифікатора особі, групі, ролі або пристрою\newline Тип: list\newline Значення: \textbf{admin, security\_\allowbreak{}officer}\newline\newline Параметр: Управління ідентифікаторами здійснюється шляхом вибору ідентифікатора, який ідентифікує окрему особу, групу, ролі або пристрій\newline Тип: list\newline Значення: \textbf{admin, security\_\allowbreak{}officer}\newline\newline Параметр: Управління ідентифікаторами здійснюється шляхом призначення ідентифікатора особі, групі, ролі або пристрою; IA-04(d) ідентифікатори управляються шляхом запобігання повторному використанню ідентифікаторів впродовж період часу\newline Тип: list\newline Значення: \textbf{admin, security\_\allowbreak{}officer}\newline\newline Параметр: Визначено персонал або ролі, від яких необхідно отримати дозвіл на призначення ідентифікатора\newline Тип: list\newline Значення: \textbf{admin, security\_\allowbreak{}officer}\newline\newline Параметр: Визначено період часу для запобігання повторному використанню ідентифікаторів\newline Тип: integer\newline Значення: \textbf{30} & IA-4 &  \\ \hline
88 & УПРАВЛІННЯ ІДЕНТИФІКАЦІЄЮ - ІДЕНТИФІКАЦІЯ СТАТУСУ КОРИСТУВАЧА (IA-4(4)) &  & IA-4(4) &  \\ \hline
89 & АВТЕНТИФІКАТОР УПРАВЛІННЯ & Параметр: Управління системними автентифікаторами здійснюється шляхом перевірки, як частини початкового розподілу автентифікатора, особи, групи, ролі або пристрою, який отримує автентифікатор\newline Тип: list\newline Значення: \textbf{admin, security\_\allowbreak{}officer}\newline\newline Параметр: Управління системними автентифікаторами здійснюється шляхом зміни/оновлення автентифікаторів у встановлений період часу або коли відбуваються події\newline Тип: list\newline Значення: \textbf{login, logout, failed\_\allowbreak{}attempt}\newline\newline Параметр: Визначено період часу для зміни або оновлення автентифікаторів за типом автентифікатора\newline Тип: integer\newline Значення: \textbf{30} & IA-5 &  \\ \hline
90 & УПРАВЛІННЯ АВТЕНТИФІКАТОРОМ - АВТЕНТИФІКАЦІЯ НА ОСНОВІ ПАРОЛЯ (IA-5(1)) & Параметр: Для автентифікації на основі паролів підтримується та оновлюється список часто використовуваних, очікуваних або скомпрометованих паролів частота, а також коли є підозра, що паролі організації були скомпрометовані прямо чи опосередковано\newline Тип: list\newline Значення: \textbf{admin, security\_\allowbreak{}officer}\newline\newline Параметр: Для автентифікації на основі паролів, коли паролі створюються або оновлюються користувачами, паролі перевіряються на відсутність у списку загальновживаних, очікуваних або скомпрометованих паролів в IA-05(01)(a)\newline Тип: list\newline Значення: \textbf{admin, security\_\allowbreak{}officer}\newline\newline Параметр: Для автентифікації на основі паролів, паролі передаються лише криптографічно захищеними каналами\newline Тип: list\newline Значення: \textbf{admin, security\_\allowbreak{}officer}\newline\newline Параметр: Для автентифікації на основі паролів паролі зберігаються за допомогою затвердженого алгоритму гешування, переважно використовуючи ключову геш-функцію\newline Тип: list\newline Значення: \textbf{admin, security\_\allowbreak{}officer}\newline\newline Параметр: Для автентифікації на основі пароля дозволяється вибір користувачем довгих паролів і фраз, що включають пробіли та всі друковані символи\newline Тип: list\newline Значення: \textbf{admin, security\_\allowbreak{}officer}\newline\newline Параметр: Для автентифікації на основі пароля використовуються автоматизовані інструменти, які допомагають користувачеві у виборі надійних автентифікаторів паролів\newline Тип: list\newline Значення: \textbf{admin, security\_\allowbreak{}officer}\newline\newline Параметр: Для автентифікації на основі пароля застосовуються склад та правила складності\newline Тип: list\newline Значення: \textbf{default\_\allowbreak{}deny\_\allowbreak{}rule, abac\_\allowbreak{}rule\_\allowbreak{}1}\newline\newline Параметр: Визначено частоту оновлення списку часто використовуваних, очікуваних або скомпрометованих паролів\newline Тип: list\newline Значення: \textbf{admin, security\_\allowbreak{}officer}\newline\newline Параметр: Визначено склад та правила складності автентифікатора\newline Тип: list\newline Значення: \textbf{default\_\allowbreak{}deny\_\allowbreak{}rule, abac\_\allowbreak{}rule\_\allowbreak{}1} & IA-5(1) &  \\ \hline
91 & УПРАВЛІННЯ АВТЕНТИФІКАТОРОМ - АВТЕНТИФІКАЦІЯ НА ОСНОВІ ВІДКРИТОГО КЛЮЧА (IA-5(2)) &  & IA-5(2) &  \\ \hline
92 & УПРАВЛІННЯ АВТЕНТИФІКАТОРОМ - АВТЕНТИФІКАЦІЯ НА ОСНОВІ АПАРАТНИХ ТОКЕНІВ (IA-5(11)) &  & IA-5(11) &  \\ \hline
93 & УПРАВЛІННЯ АВТЕНТИФІКАТОРОМ - ЕФЕКТИВНІСТЬ БІОМЕТРИЧНОЇ АВТЕНТИФІКАЦІЇ (IA-5(12)) & Параметр: Визначено вимоги до якості біометрії; для біометричної автентифікації використовувати механізми, які задовольняють вимоги\newline Тип: string\newline Значення: \textbf{автоматизований засіб моніторингу} & IA-5(12) &  \\ \hline
94 & УПРАВЛІННЯ АВТЕНТИФІКАТОРОМ - ПРОДУКТИ ТА ПОСЛУГИ, ЗАТВЕРДЖЕНІ УПОВНОВАЖЕНИМ ОРГАНОМ (IA-5(15)) &  & IA-5(15) &  \\ \hline
95 & УПРАВЛІННЯ АВТЕНТИФІКАТОРОМ - АВТОМАТИЗОВАНІ ЗАСОБИ ВИЯВЛЕННЯ АТАК ІЗ ВИКОРИСТАННЯМ БІОМЕТРИЧНИХ АВТЕНТИФІКАТОРІВ (IA-5(17)) & Параметр: Використовуються механізми виявлення атак із використанням штучно виготовлених артефактів для біометричних автентифікаторів\newline Тип: string\newline Значення: \textbf{автоматизований засіб моніторингу} & IA-5(17) &  \\ \hline
96 & ПРИХОВУВАННЯ ЗВОРОТНОГО ЗВ'ЯЗКУ АВТЕНТИФІКАТОРА & Параметр: Забезпечино приховану зворотну передачу інформації автентифікації в про- цесі автентифікації для забезпечення захисту інформації від можливої експлуатації та використання неавторизованими особами\newline Тип: list\newline Значення: \textbf{admin, security\_\allowbreak{}officer} & IA-6 &  \\ \hline
97 & АВТЕНТИФІКАЦІЯ КРИПТОГРАФІЧНОГО МОДУЛЯ (IA-7) & Параметр: Впроваджено механізми автентифікації в криптографічний модуль, який відповідає вимогам чинних законів, виконавчих розпоряджень, директив, політик, правил, стандартів та рекомендацій для такої автентифікації\newline Тип: list\newline Значення: \textbf{default\_\allowbreak{}deny\_\allowbreak{}rule, abac\_\allowbreak{}rule\_\allowbreak{}1} & IA-7 &  \\ \hline
98 & ПОВТОРНА АВТЕНТИФІКАЦІЯ (IA-11) &  & IA-11 &  \\ \hline
\multicolumn{5}{|l|}{\textbf{8. РЕАГУВАННЯ НА ІНЦИДЕНТИ (IR) (IR)}} \\ \hline
99 & Політика та процедури реагування на інциденти (IR-1) & Параметр: Розроблено та задокументовано політику реагування на інциденти\newline Тип: list\newline Значення: \textbf{default\_\allowbreak{}deny\_\allowbreak{}rule, abac\_\allowbreak{}rule\_\allowbreak{}1}\newline\newline Параметр: Політика реагування на інциденти поширюється серед персоналу або ролей\newline Тип: list\newline Значення: \textbf{admin, security\_\allowbreak{}officer}\newline\newline Параметр: Розроблені та задокументовані процедури реагування на інциденти, що сприяють впровадженню політики реагування на інциденти та пов'язаних з нею заходів захисту з реагування на інциденти\newline Тип: list\newline Значення: \textbf{default\_\allowbreak{}deny\_\allowbreak{}rule, abac\_\allowbreak{}rule\_\allowbreak{}1}\newline\newline Параметр: Посадова особа призначається для управління розробкою, документуванням та розповсюдженням політики та процедур реагування на інциденти; IR-01(c)[01][01] переглядається та оновлюється поточна політика реагування на інциденти частота; IR-01(c)[01][02] поточна політика реагування на інциденти переглядається та оновлюється після подій; IR-01(c)[02][01] переглядаються та оновлюються поточні процедури реагування на інциденти частота; IR-01(c)[02][02] поточні процедури реагування на інциденти переглядаються та оновлюються після подій\newline Тип: list\newline Значення: \textbf{admin, security\_\allowbreak{}officer}\newline\newline Параметр: Визначено персонал або ролі, до яких має бути доведена політика реагування на інциденти\newline Тип: list\newline Значення: \textbf{admin, security\_\allowbreak{}officer}\newline\newline Параметр: Визначено персонал або ролі, до яких мають бути доведені процедури реагування на інциденти\newline Тип: list\newline Значення: \textbf{admin, security\_\allowbreak{}officer}\newline\newline Параметр: Визначено посадову особу, яка керуватиме політикою та процедурами реагування на інциденти\newline Тип: list\newline Значення: \textbf{admin, security\_\allowbreak{}officer}\newline\newline Параметр: Визначено частоту, з якою переглядається та оновлюється поточна політика реагування на інциденти\newline Тип: list\newline Значення: \textbf{default\_\allowbreak{}deny\_\allowbreak{}rule, abac\_\allowbreak{}rule\_\allowbreak{}1}\newline\newline Параметр: Визначаються події, які потребують перегляду та оновлення поточної політики реагування на інциденти\newline Тип: list\newline Значення: \textbf{default\_\allowbreak{}deny\_\allowbreak{}rule, abac\_\allowbreak{}rule\_\allowbreak{}1}\newline\newline Параметр: Визначено частоту, з якою переглядаються та оновлюються поточні процедури реагування на інциденти\newline Тип: integer\newline Значення: \textbf{30}\newline\newline Параметр: Визначено події, які потребують перегляду та оновлення процедур реагування на інциденти\newline Тип: list\newline Значення: \textbf{login, logout, failed\_\allowbreak{}attempt} & IR-1 &  \\ \hline
100 & НАВЧАННЯ З РЕАГУВАННЯ НА ІНЦИДЕНТИ (IR-2) & Параметр: Навчання з реагування на інциденти надається користувачам системи відповідно до призначених ролей та обов'язків протягом періоду часу з моменту прийняття на себе ролі або обов'язків з реагування на інциденти або отримання доступу до системи\newline Тип: list\newline Значення: \textbf{admin, security\_\allowbreak{}officer}\newline\newline Параметр: Користувачам системи надається навчання з реагування на інциденти відповідно до призначених ролей та обов'язків частота\newline Тип: string\newline Значення: \textbf{щорічно}\newline\newline Параметр: Зміст навчання з реагування на інциденти переглядається та оновлюється частота\newline Тип: string\newline Значення: \textbf{щорічно}\newline\newline Параметр: Визначено частоту, з якою користувачі повинні проходити навчання з реагування на інциденти\newline Тип: integer\newline Значення: \textbf{30}\newline\newline Параметр: Визначено частоту перегляду та оновлення змісту навчання з реагування на інциденти\newline Тип: integer\newline Значення: \textbf{30}\newline\newline Параметр: Визначено події, які ініціюють перегляд змісту навчання з реагування на інциденти\newline Тип: list\newline Значення: \textbf{login, logout, failed\_\allowbreak{}attempt} & IR-2 &  \\ \hline
101 & ПЕРЕВІРКА РЕАГУВАНЬ НА ІНЦИДЕНТИ (IR-3) & Параметр: Ефективність реагування системи на інциденти перевіряється частота за допомогою тестів\newline Тип: string\newline Значення: \textbf{щорічно}\newline\newline Параметр: Визначено частоту, з якою необхідно перевіряти ефективність реагування системи на інциденти\newline Тип: integer\newline Значення: \textbf{30} & IR-3 &  \\ \hline
102 & Обробка інциденту (IR-4) &  & IR-4 &  \\ \hline
103 & Моніторинг інциденту (IR-5) &  & IR-5 &  \\ \hline
104 & ЗВІТНІСТЬ ПРО ІНЦИДЕНТИ (IR-6) & Параметр: Персонал зобов'язаний повідомляти про підозрілі інциденти протягом періоду часу\newline Тип: list\newline Значення: \textbf{admin, security\_\allowbreak{}officer}\newline\newline Параметр: Визначено період часу, протягом якого персонал повинен повідомляти про підозрілі інциденти до уповноваженого органу\newline Тип: list\newline Значення: \textbf{admin, security\_\allowbreak{}officer} & IR-6 &  \\ \hline
105 & ПІДТРИМКА РЕАГУВАННЯ НА ІНЦИДЕНТИ (IR-7) &  & IR-7 &  \\ \hline
106 & План реагування на інциденти (IR-8) & Параметр: Розроблено план реагування на інцидент, який розглядається та затверджується персоналом або ролями частота\newline Тип: list\newline Значення: \textbf{admin, security\_\allowbreak{}officer}\newline\newline Параметр: Розроблено план реагування на інциденти, в якому чітко визначено відповідальність за реагування на інциденти для організацій, персоналу або ролей\newline Тип: list\newline Значення: \textbf{admin, security\_\allowbreak{}officer}\newline\newline Параметр: Копії плану реагування на інцидент розповсюджуються серед <IR- 08\_\allowbreak{}ODP[04] персоналу з реагування на інциденти>\newline Тип: list\newline Значення: \textbf{admin, security\_\allowbreak{}officer}\newline\newline Параметр: План реагування на інциденти оновлюється з урахуванням змін у системі та організації або проблем, що виникають під час впровадження, виконання або тестування плану\newline Тип: integer\newline Значення: \textbf{30}\newline\newline Параметр: Зміни в плані реагування на інцидент повідомляються персоналу з реагування на інциденти\newline Тип: list\newline Значення: \textbf{admin, security\_\allowbreak{}officer}\newline\newline Параметр: Визначено персонал або ролі, які переглядають та затверджують план реагування на інциденти\newline Тип: list\newline Значення: \textbf{admin, security\_\allowbreak{}officer}\newline\newline Параметр: Визначено періодичність перегляду та затвердження плану реагування на інциденти\newline Тип: string\newline Значення: \textbf{щорічно}\newline\newline Параметр: Визначені організації, персонал або ролі, які несуть відповідальність за реагування на інциденти; IR-08\_\allowbreak{}ODP[04] визначено персонал з реагування на інцидент (ідентифікований за іменами та/або за ролями), якому мають бути роздані копії плану реагування на інцидент\newline Тип: list\newline Значення: \textbf{admin, security\_\allowbreak{}officer}\newline\newline Параметр: Визначено персонал з реагування на інцидент (ідентифікований за іменами та/або ролями), якому повідомляються зміни до плану реагування на інцидент\newline Тип: list\newline Значення: \textbf{admin, security\_\allowbreak{}officer} & IR-8 &  \\ \hline
\multicolumn{5}{|l|}{\textbf{9. ТЕХНІЧНЕ ОБСЛУГОВУВАННЯ (MA) (MA)}} \\ \hline
107 & ПОЛІТИКА ТА ПРОЦЕДУРИ ТЕХНІЧНОГО ОБСЛУГОВУВАННЯ (MA-1) & Параметр: Розроблено та задокументовано політику технічного обслуговування\newline Тип: list\newline Значення: \textbf{default\_\allowbreak{}deny\_\allowbreak{}rule, abac\_\allowbreak{}rule\_\allowbreak{}1}\newline\newline Параметр: Політика технічного обслуговування поширюється на персонал або ролі\newline Тип: list\newline Значення: \textbf{admin, security\_\allowbreak{}officer}\newline\newline Параметр: Розроблені та задокументовані процедури технічного обслуговування, що сприяють впровадженню політики технічного обслуговування та пов'язаних з нею заходів технічного обслуговування\newline Тип: list\newline Значення: \textbf{default\_\allowbreak{}deny\_\allowbreak{}rule, abac\_\allowbreak{}rule\_\allowbreak{}1}\newline\newline Параметр: Посадова особа призначається для управління розробкою, документуванням та розповсюдженням політики та процедур технічного обслуговування; MA-01(c)[01][01] переглядається та оновлюється поточна політика обслуговування частота; MA-01(c)[01][02] переглядається та оновлюється поточна політика обслуговування після подій; MA-01(c)[02][01] переглядаються та оновлюються поточні процедури технічного обслуговування частота; MA-01(c)[02][02] переглядаються та оновлюються поточні процедури технічного обслуговування після подій\newline Тип: list\newline Значення: \textbf{admin, security\_\allowbreak{}officer}\newline\newline Параметр: Визначено персонал або ролі, на які поширюється політика технічного обслуговування\newline Тип: list\newline Значення: \textbf{admin, security\_\allowbreak{}officer}\newline\newline Параметр: Визначено персонал або ролі, на які поширюються процедури технічного обслуговування\newline Тип: list\newline Значення: \textbf{admin, security\_\allowbreak{}officer}\newline\newline Параметр: Визначено посадову особу, яка керуватиме політикою та процедурами технічного обслуговування\newline Тип: list\newline Значення: \textbf{admin, security\_\allowbreak{}officer}\newline\newline Параметр: Визначено частоту, з якою переглядається та оновлюється поточна політика технічного обслуговування\newline Тип: list\newline Значення: \textbf{default\_\allowbreak{}deny\_\allowbreak{}rule, abac\_\allowbreak{}rule\_\allowbreak{}1}\newline\newline Параметр: Визначено події, які потребують перегляду та оновлення поточної політики технічного обслуговування\newline Тип: list\newline Значення: \textbf{default\_\allowbreak{}deny\_\allowbreak{}rule, abac\_\allowbreak{}rule\_\allowbreak{}1}\newline\newline Параметр: Визначено частоту, з якою переглядаються та оновлюються поточні процедури технічного обслуговування\newline Тип: integer\newline Значення: \textbf{30}\newline\newline Параметр: Визначено події, які потребують перегляду та оновлення процедур технічного обслуговування\newline Тип: list\newline Значення: \textbf{login, logout, failed\_\allowbreak{}attempt} & MA-1 & a. Планувати, документувати та переглядати записи з технічного обслуговування, ремонту або заміни компонентів системи відповідно до вимог виробника та постачальників та/або вимог організації.\\ 
b. Затвердити та здійснювати моніторинг усіх заходів з технічного обслуговування, незалежно від того, виконуються вони на місці або віддалено, а також чи обслуговуються системи або системні компоненти на місці, чи переміщуються в інше місце.\\ 
c. Вимагати, щоб [Призначення: визначені організацією персонал чи ролі] явно схвалили видалення системи або компоненту системи з організаційного обладнання для технічного обслуговування, ремонту чи заміни поза об'єктами експлуатації.\\ 
d. Очищати обладнання з погляду видалення всієї інформації з носіїв до вилучення обладнання організації для технічного обслуговування, ремонту чи заміни поза об'єктами експлуатації.\\ 
e. Перевірити всі потенційно порушені заходи захисту, щоб переконатися, що вони, як і раніше, працюють належним чином після дій з обслуговування, ремонту або заміни.\\ 
f. Вносити [Призначення: визначену організацією інформацію, пов'язану з технічним обслуговуванням] до записів з технічного обслуговування. \\ \hline
108 & ІНСТРУМЕНТИ ДЛЯ ОБСЛУГОВУВАННЯ (MA-3) & Параметр: Переглядаються раніше затверджені інструменти технічного обслуговування частота\newline Тип: string\newline Значення: \textbf{щорічно}\newline\newline Параметр: Визначено частоту, з якою слід переглядати раніше затверджені інструменти технічного обслуговування\newline Тип: integer\newline Значення: \textbf{30} & MA-3 & a. Затвердити, контролювати та відстежувати використання засобів технічного обслуговування.\\ 
b. Переглядати раніше затверджені інструменти технічного [Призначення: з частотою, визначеною організацією]. обслуговування \\ \hline
109 & ІНСТРУМЕНТИ ДЛЯ ОБСЛУГОВУВАННЯ - ПЕРЕВІРКА ІНСТРУМЕНТІВ (MA-3(1)) & Параметр: Оглядаються інструменти для технічного обслуговування, які доставлені на об'єкт обслуговуючим персоналом, на предмет неправильних або несанкціонованих модифікацій\newline Тип: list\newline Значення: \textbf{admin, security\_\allowbreak{}officer} & MA-3(1) &  \\ \hline
110 & ІНСТРУМЕНТИ ДЛЯ ОБСЛУГОВУВАННЯ - ПЕРЕВІРКА НОСІЇВ ІНФОРМАЦІЇ (MA-3(2)) &  & MA-3(2) & Перед використанням носіїв у системі перевірити носії, що містять діагностичні та тестові програми на наявність шкідливого коду. \\ \hline
111 & ІНСТРУМЕНТИ ДЛЯ ОБСЛУГОВУВАННЯ - ЗАПОБІГАННЯ НЕСАНКЦІОНОВАНОМУ ПЕРЕМІЩЕННЮ (MA-3(3)) & Параметр: Переміщення обладнання для технічного обслуговування, що містить інформацію організації, запобігається шляхом отримання дозволу від персонал або ролі\newline Тип: list\newline Значення: \textbf{admin, security\_\allowbreak{}officer}\newline\newline Параметр: Визначено персонал або ролі, які можуть надавати дозвіл на переміщення обладнання з об'єкту\newline Тип: list\newline Значення: \textbf{admin, security\_\allowbreak{}officer} & MA-3(3) & Запобігти переміщенню обладнання для технічного обслуговування, що містить організаційну інформацію, шляхом:\\ 
(a) перевірки відсутності організаційної інформації, розміщеної на обладнанні;\\ 
(b) очищення або знищення обладнання;\\ 
(c) утримання обладнання на об'єкті;\\ 
(d) отримання дозволу від [Призначення: визначених організацією персоналу чи ролей], які явно дозволяють переміщення обладнання з об'єкта. \\ \hline
112 & ВІДДАЛЕНЕ ОБСЛУГОВУВАННЯ (MA-4) & Параметр: Впроваджено віддалені дії з обслуговування та діагностики\newline Тип: list\newline Значення: \textbf{login, logout, failed\_\allowbreak{}attempt}\newline\newline Параметр: Відстежуються віддалені дії з обслуговування та діагностики\newline Тип: list\newline Значення: \textbf{login, logout, failed\_\allowbreak{}attempt}\newline\newline Параметр: Використання віддалених засобів технічного обслуговування та діагностики дозволено лише відповідно до політики організації\newline Тип: list\newline Значення: \textbf{default\_\allowbreak{}deny\_\allowbreak{}rule, abac\_\allowbreak{}rule\_\allowbreak{}1} & MA-4 &  \\ \hline
113 & ТЕХНІЧНИЙ ПЕРСОНАЛ (MA-5) & Параметр: Запроваджено процес авторизації технічного персоналу\newline Тип: list\newline Значення: \textbf{admin, security\_\allowbreak{}officer}\newline\newline Параметр: Ведеться перелік авторизованих організацій або персоналу з технічного обслуговування\newline Тип: list\newline Значення: \textbf{admin, security\_\allowbreak{}officer}\newline\newline Параметр: Персонал без супроводу, який виконує технічне обслуговування системи, має необхідні дозволи на доступ\newline Тип: list\newline Значення: \textbf{admin, security\_\allowbreak{}officer}\newline\newline Параметр: Персонал організації з необхідними повноваженнями доступу та технічною компетентністю призначений/призначені для нагляду за діяльністю з технічного обслуговування персоналу, який не має необхідних дозволів на доступ\newline Тип: list\newline Значення: \textbf{admin, security\_\allowbreak{}officer} & MA-5 & a. Встановити процедуру авторизації технічного персоналу та вести перелік авторизованих організацій технічного обслуговування або персоналу.\\ 
b. Перевіряти, що персонал, який не супроводжується та виконує технічне обслуговування в системі, має необхідні дозволи на доступ.\\ 
c. Визначити персонал організації з необхідними повноваженнями щодо доступу та технічною компетенцією для нагляду за персоналом з технічного обслуговування, який не має необхідних дозволів на доступ. \\ \hline
\multicolumn{5}{|l|}{\textbf{10. ЗАХИСТ НОСІЇВ ІНФОРМАЦІЇ (MP) (MP)}} \\ \hline
114 & ПОЛІТИКА ТА ПРОЦЕДУРИ ЩОДО ЗАХИСТУ НОСІЇВ ІНФОРМАЦІЇ (MP-1) & Параметр: Розроблено та задокументовано політику захисту носіїв інформації\newline Тип: list\newline Значення: \textbf{default\_\allowbreak{}deny\_\allowbreak{}rule, abac\_\allowbreak{}rule\_\allowbreak{}1}\newline\newline Параметр: Політика захисту носіїв інформації поширюється на персонал або ролі\newline Тип: list\newline Значення: \textbf{admin, security\_\allowbreak{}officer}\newline\newline Параметр: Розроблено та задокументовано процедури захисту носіїв інформації, що сприятимуть реалізації політики захисту носіїв інформації та заходів захисту носіїв інформації\newline Тип: list\newline Значення: \textbf{default\_\allowbreak{}deny\_\allowbreak{}rule, abac\_\allowbreak{}rule\_\allowbreak{}1}\newline\newline Параметр: Посадова особа призначається для управління розробкою, документуванням та розповсюдженням політики та процедур захисту носіїв інформації. MP-01(c)[01][01] переглядається та оновлюється поточна політика захисту носіїв інформації частота; MP-01(c)[01][02] переглядається та оновлюється поточна політика захисту носіїв інформації після подій; MP-01(c)[02][01] переглядаються та оновлюються поточні процедури захисту носіїв інформації частота; MP-01(c)[02][02] переглядаються та оновлюються поточні процедури захисту носіїв інформації після подій\newline Тип: list\newline Значення: \textbf{admin, security\_\allowbreak{}officer}\newline\newline Параметр: Визначено персонал або ролі, серед яких має бути поширена політика захисту носіїв інформації\newline Тип: list\newline Значення: \textbf{admin, security\_\allowbreak{}officer}\newline\newline Параметр: Визначено персонал або ролі, серед яких мають бути поширені процедури захисту носіїв інформації\newline Тип: list\newline Значення: \textbf{admin, security\_\allowbreak{}officer}\newline\newline Параметр: Визначено посадову особу, яка керуватиме політикою та процедурами захисту носіїв інформації\newline Тип: list\newline Значення: \textbf{admin, security\_\allowbreak{}officer}\newline\newline Параметр: Визначено частоту, з якою переглядається та оновлюється поточна політика захисту носіїв інформації\newline Тип: list\newline Значення: \textbf{default\_\allowbreak{}deny\_\allowbreak{}rule, abac\_\allowbreak{}rule\_\allowbreak{}1}\newline\newline Параметр: Визначено події, які потребують перегляду та оновлення чинної політики захисту носіїв інформації\newline Тип: list\newline Значення: \textbf{default\_\allowbreak{}deny\_\allowbreak{}rule, abac\_\allowbreak{}rule\_\allowbreak{}1}\newline\newline Параметр: Визначено частоту, з якою переглядаються та оновлюються чинні процедури захисту носіїв інформації\newline Тип: integer\newline Значення: \textbf{30}\newline\newline Параметр: Визначено події, які потребують перегляду та оновлення процедур захисту носіїв інформації\newline Тип: list\newline Значення: \textbf{login, logout, failed\_\allowbreak{}attempt} & MP-1 & a. Розробити, задокументувати та поширити серед [Призначення: визначеного організацією персоналу або посад]:\\ 
1. 2. політику захисту носіїв інформації, яка:\\ 
(a) містить мету, сферу застосування, ролі, обов'язки, відповідальність керівництва, координацію між організаційними підрозділами та систему контролю відповідності (complaince);\\ 
(b) відповідає чинному законодавству, виконавчим наказам, директивам, нормам, політикам, стандартам та керівним принципам; процедури, які сприяють здійсненню політики та заходів захисту носіїв інформації.\\ 
b. Призначити [Призначення: визначену організацією посадову особу] для управління розробкою, документування, та розповсюдження політики та процедурами захисту носіїв інформації.\\ 
c. Переглядати та оновлювати чинну систему захисту носіїв інформації:\\ 
1. поточну політику захисту носіїв інформації [Призначення: з визначеною організацією частотою];\\ 
2. поточні процедури захисту носіїв інформації [Призначення: з визначеною організацією частотою]. \\ \hline
115 & ДОСТУП ДО НОСІЇВ ІНФОРМАЦІЇ (MP-2) & Параметр: Доступ до типів цифрових носіїв інформації обмежено для персоналу або ролей\newline Тип: list\newline Значення: \textbf{admin, security\_\allowbreak{}officer}\newline\newline Параметр: Доступ до типів нецифрових носіїв інформації обмежено для персоналу або ролей\newline Тип: list\newline Значення: \textbf{admin, security\_\allowbreak{}officer}\newline\newline Параметр: Визначено персонал або ролі, уповноважені на доступ до цифрових носіїв інформації\newline Тип: list\newline Значення: \textbf{admin, security\_\allowbreak{}officer}\newline\newline Параметр: Визначено персонал або ролі, уповноважені на доступ до нецифрових носіїв інформації\newline Тип: list\newline Значення: \textbf{admin, security\_\allowbreak{}officer} & MP-2 & Обмежити доступ до [Призначення: визначених організацією типів цифрових та/або нецифрових носіїв інформації] [Призначення: визначеним організацією персоналом або ролями]. \\ \hline
116 & МАРКУВАННЯ НОСІЇВ ІНФОРМАЦІЇ (MP-3) & Параметр: Визначено типи носіїв інформації, які звільняються від маркування під час перебування на контрольованих зонах\newline Тип: integer\newline Значення: \textbf{30} & MP-3 & a. Наносити на носії інформації маркування, що вказують на обмеження поширення, обробки, а також застереження та відповідні мітки безпеки (якщо такі є) інформації.\\ 
b. Звільнити [Призначення: визначені організацією типи носіїв системи] від маркування, якщо носії залишаються в межах [Призначення: визначених організацією контрольованих зон]. \\ \hline
117 & ЗБЕРІГАННЯ НОСІЇВ ІНФОРМАЦІЇ (MP-4) &  & MP-4 & a. Фізично контролювати та безпечно зберігати [Призначення: визначені організацією типи цифрових та/або нецифрових носіїв інформації] в межах [Призначення: визначених організацією контрольованих зон].\\ 
b. Захищати системні носії, які визначені в МР-4 до того часу, як носії знищуються або очищаються, з використанням затвердженого обладнання, методів та процедур. \\ \hline
118 & Транспортування носіїв інформації (MP-5) & Параметр: Типи носіїв інформації системи захищаються під час транспортування за межі контрольованих зон за допомогою заходів безпеки\newline Тип: integer\newline Значення: \textbf{30}\newline\newline Параметр: Типи носіїв інформації системи контролюються під час транспортування за межі контрольованих зон за допомогою заходів безпеки\newline Тип: integer\newline Значення: \textbf{30}\newline\newline Параметр: Під час транспортування за межі контрольованих зон ведеться облік носіїв інформації системи\newline Тип: integer\newline Значення: \textbf{30}\newline\newline Параметр: Визначено персонал, уповноважений здійснювати діяльність з транспортування носіїв інформації\newline Тип: list\newline Значення: \textbf{admin, security\_\allowbreak{}officer}\newline\newline Параметр: Діяльність, пов'язана з транспортуванням носіїв інформації системи, обмежується визначеним уповноваженим персоналом\newline Тип: list\newline Значення: \textbf{admin, security\_\allowbreak{}officer}\newline\newline Параметр: Визначено типи носіїв інформації системи для захисту та контролю під час транспортування за межі контрольованих зон\newline Тип: integer\newline Значення: \textbf{30} & MP-5 & a. Захищати та контролювати [Призначення: визначені організацією типи носіїв системи] під час транспортування за межами контрольованих зон, використовуючи [Призначення: визначені організацією заходи безпеки].\\ 
b. Вести облік носіїв системи інформації під час транспортування за межами контрольованих зон.\\ 
c. Документувати дії, пов'язані з транспортуванням носіїв системи.\\ 
d. Обмежити діяльність уповноваженого персоналу, пов'язану з транспортуванням носіїв системи. \\ \hline
119 & ЗНИЩЕННЯ ІНФОРМАЦІЇ НА НОСІЯХ ІНФОРМАЦІЇ (MP-6) & Параметр: Застосовуються механізми очищення, надійність і цілісність яких відповідає категорії безпеки або рівню секретності інформації\newline Тип: string\newline Значення: \textbf{автоматизований засіб моніторингу} & MP-6 & a. Очищувати [Призначення: визначені організацією системні носії] перед утилізацією, випуском за межі організаційного контролю, або перед повторним використанням [Призначення: методами та процедурами очищення, визначеними організацією].\\ 
b. Використовувати механізми очищення зі стійкістю та цілісністю, що відповідає категорії безпеки або рівню секретності інформації. \\ \hline
120 & ВИКОРИСТАННЯ НОСІЇВ ІНФОРМАЦІЇ (MP-7) &  & MP-7 & a. [Вибір: обмежити; заборонити] використання [Призначення: визначених організацією типів носіїв системи] на [Призначення: визначені організацією системи або компоненти системи], використовуючи [Призначення: визначені організацією заходи безпеки].\\ 
b. Заборонити використання портативних пристроїв зберігання даних в системах організації, якщо такі пристрої не мають визначеного власника. \\ \hline
121 & ЗНИЖЕННЯ КАТЕГОРІЇ БЕЗПЕКИ НОСІЇВ ІНФОРМАЦІЇ - ТАЄМНА ІНФОРМАЦІЯ (MP-8(4)) & Параметр: Носії інформації, що містять інформацію з обмеженим доступом, знижуються в класі перед передачею особам, які не мають необхідних дозволів на доступ\newline Тип: list\newline Значення: \textbf{admin, security\_\allowbreak{}officer} & MP-8(4) &  \\ \hline
\multicolumn{5}{|l|}{\textbf{11. ФІЗИЧНИЙ ЗАХИСТ ТА ЗАХИСТ НАВКОЛИШНЬОГО СЕРЕДОВИЩА (PE) (PE)}} \\ \hline
122 & РЕ-1 фізичного захисту та захисту робочого середовища Авторизація фізичного (PE-1) & Параметр: Розроблено та задокументовано політику фізичного захисту та захисту робочого середовища\newline Тип: list\newline Значення: \textbf{default\_\allowbreak{}deny\_\allowbreak{}rule, abac\_\allowbreak{}rule\_\allowbreak{}1}\newline\newline Параметр: Політика фізичного захисту та захисту робочого середовища розповсюджується серед персоналу або ролей\newline Тип: list\newline Значення: \textbf{admin, security\_\allowbreak{}officer}\newline\newline Параметр: Посадова особа призначається для управління розробкою, документуванням та розповсюдженням політики та процедур фізичного захисту та захисту робочого середовища; PE-01(c)[01][01] переглядається та оновлюється поточна політика фізичного захисту та захисту робочого середовища частота; PE-01(c)[01][02] поточна політика фізичного захисту та захисту робочого середовища переглядається та оновлюється після подій; PE-01(c)[02][01] переглядаються та оновлюються поточні процедури фізичного захисту та захисту робочого середовища частота; поточні процедури фізичного та екологічного захисту переглядаються та оновлюються після подій. PE-01(c)[02][02]\newline Тип: list\newline Значення: \textbf{admin, security\_\allowbreak{}officer}\newline\newline Параметр: Визначено персонал або ролі, до яких має бути доведена політика фізичного захисту та захисту робочого середовища\newline Тип: list\newline Значення: \textbf{admin, security\_\allowbreak{}officer}\newline\newline Параметр: Визначено персонал або ролі, на які поширюються процедури фізичного захисту та захисту робочого середовища\newline Тип: list\newline Значення: \textbf{admin, security\_\allowbreak{}officer}\newline\newline Параметр: Визначено посадову особу, яка керуватиме політикою та процедурами фізичного захисту та захисту робочого середовища\newline Тип: list\newline Значення: \textbf{admin, security\_\allowbreak{}officer}\newline\newline Параметр: Визначено частоту, з якою переглядається та оновлюється поточна політика фізичного захисту та захисту робочого середовища\newline Тип: list\newline Значення: \textbf{default\_\allowbreak{}deny\_\allowbreak{}rule, abac\_\allowbreak{}rule\_\allowbreak{}1}\newline\newline Параметр: Визначено події, які потребують перегляду та оновлення поточної політики фізичного захисту та захисту робочого середовища\newline Тип: list\newline Значення: \textbf{default\_\allowbreak{}deny\_\allowbreak{}rule, abac\_\allowbreak{}rule\_\allowbreak{}1}\newline\newline Параметр: Визначено частоту, з якою переглядаються та оновлюються поточні процедури фізичного захисту та захисту робочого середовища\newline Тип: integer\newline Значення: \textbf{30}\newline\newline Параметр: Визначено події, які потребують перегляду та оновлення процедур фізичного захисту та захисту робочого середовища\newline Тип: list\newline Значення: \textbf{login, logout, failed\_\allowbreak{}attempt} & PE-1 &  \\ \hline
123 & РЕ-2 доступу (PE-2) & Параметр: Розроблено перелік осіб, які мають право авторизованого доступу до об'єкта, де перебуває система\newline Тип: list\newline Значення: \textbf{}\newline\newline Параметр: Затверджено перелік осіб, які мають право авторизованого доступу до об'єкта, де перебуває система\newline Тип: list\newline Значення: \textbf{}\newline\newline Параметр: Ведеться перелік осіб, які мають право авторизованого доступу до об'єкта, де перебуває система\newline Тип: list\newline Значення: \textbf{}\newline\newline Параметр: Переглядається список доступу, , у якому закріплений перелік персоналу або ролей, яким дозволений санкціонований доступ до об'єкта частота\newline Тип: list\newline Значення: \textbf{admin, security\_\allowbreak{}officer}\newline\newline Параметр: Особи видаляються зі списку доступу до об'єкта, коли доступ більше не потрібен\newline Тип: list\newline Значення: \textbf{admin, security\_\allowbreak{}officer}\newline\newline Параметр: Визначено періодичність перегляду списку доступу, у якому закріплений перелік персоналу або ролей, яким дозволений санкціонований доступ до об'єкта\newline Тип: list\newline Значення: \textbf{admin, security\_\allowbreak{}officer} & PE-2 &  \\ \hline
124 & ФІЗИЧНИЙ ДОСТУП ДО СИСТЕМИ & Параметр: Активність відвідувачів контролюється умови\newline Тип: list\newline Значення: \textbf{}\newline\newline Параметр: Пристрої фізичного доступу інвентаризуються частота\newline Тип: string\newline Значення: \textbf{щорічно}\newline\newline Параметр: Коди доступу змінюється частота, коли код скомпрометовано, або коли особи, які володіють кодом, переводяться або звільняються; <PE- PE-03(g)[02] ключі змінюються частота, коли ключі втрачено, або коли особи, що володіють ключами, переводяться або звільняються\newline Тип: list\newline Значення: \textbf{admin, security\_\allowbreak{}officer}\newline\newline Параметр: Визначено умови, що вимагають супроводу відвідувачів та моніторингу активності відвідувачів\newline Тип: list\newline Значення: \textbf{}\newline\newline Параметр: Визначено частоту проведення інвентаризації пристроїв фізичного доступу\newline Тип: integer\newline Значення: \textbf{30}\newline\newline Параметр: Визначено частоту, з якою потрібно змінювати коди доступу\newline Тип: integer\newline Значення: \textbf{30}\newline\newline Параметр: Визначено частоту, з якою потрібно змінювати ключі\newline Тип: integer\newline Значення: \textbf{30} & PE-3 &  \\ \hline
125 & РЕ-4 ліній електроживлення (PE-4) &  & PE-4 &  \\ \hline
126 & КОНТРОЛЬ ДОСТУПУ В ПРИМІЩЕННЯ ДЛЯ ВІДОБРАЖЕННЯ ІНФОРМАЦІЇ &  & PE-5 &  \\ \hline
127 & МОНІТОРИНГ ФІЗИЧНОГО ДОСТУПУ & Параметр: Переглядаються журнали фізичного доступу частота\newline Тип: string\newline Значення: \textbf{щорічно}\newline\newline Параметр: Визначено частоту перегляду журналів фізичного доступу\newline Тип: integer\newline Значення: \textbf{30}\newline\newline Параметр: Визначено події або потенційні ознаки подій, що вимагають перегляду журналів фізичного доступу\newline Тип: list\newline Значення: \textbf{login, logout, failed\_\allowbreak{}attempt} & PE-6 &  \\ \hline
128 & АЛЬТЕРНАТИВНЕ РОБОЧЕ МІСЦЕ (PE-17) & Параметр: Визначаються заходи захисту, які будуть застосовуватися на альтернативних робочих місцях; РЕ-17(a) альтернативні робочі місця визначені та задокументовані; РЕ-17(b) заходи захисту впроваджені на альтернативних робочих місцях; РЕ-17(c) оцінюється ефективність заходів захисту на альтернативних робочих місцях; РЕ-17(d) працівникам надаються засоби комунікації з персоналом служби інформаційної безпеки на випадок інцидентів\newline Тип: list\newline Значення: \textbf{admin, security\_\allowbreak{}officer} & PE-17 &  \\ \hline
129 & МАРКУВАННЯ КОМПОНЕНТІВ (PE-22) &  & PE-22 &  \\ \hline
\multicolumn{5}{|l|}{\textbf{12. ПЛАНУВАННЯ (PL) (PL)}} \\ \hline
130 & ПОЛІТИКИ ТА ПРОЦЕДУРИ ПЛАНУВАННЯ БЕЗПЕКИ (PL-1) & Параметр: Визначено персонал або ролі, до яких має бути доведена політика планування безпеки\newline Тип: list\newline Значення: \textbf{admin, security\_\allowbreak{}officer}\newline\newline Параметр: Визначено персонал або ролі, на які поширюватимуться процедури планування безпеки\newline Тип: list\newline Значення: \textbf{admin, security\_\allowbreak{}officer}\newline\newline Параметр: Визначено посадову особу, яка керуватиме політикою та процедурами планування безпеки\newline Тип: list\newline Значення: \textbf{admin, security\_\allowbreak{}officer}\newline\newline Параметр: Визначено періодичність перегляду та оновлення поточної політики планування безпеки\newline Тип: list\newline Значення: \textbf{default\_\allowbreak{}deny\_\allowbreak{}rule, abac\_\allowbreak{}rule\_\allowbreak{}1}\newline\newline Параметр: Є події, які потребують перегляду та оновлення поточної політики планування безпеки\newline Тип: list\newline Значення: \textbf{default\_\allowbreak{}deny\_\allowbreak{}rule, abac\_\allowbreak{}rule\_\allowbreak{}1}\newline\newline Параметр: Визначена частота, з якою переглядаються та оновлюються поточні процедури планування безпеки\newline Тип: string\newline Значення: \textbf{щорічно} & PL-1 & a. Розробити, задокументувати та поширити серед [Призначення: визначеного організацією персоналу або ролей]:\\ 
1. 2. [вибір (один або декілька): рівень організації; рівень місії/бізнес процесу; рівень системи] політику планування, яка:\\ 
(a) містить мету, сферу застосування, ролі, обов'язки, відповідальність керівництва, координацію між організаційними підрозділами та системою контролю (complaince);\\ 
(b) відповідає чинному законодавству, виконавчим наказам, директивам, нормам, політикам, стандартам та керівним принципам; процедури, що полегшують здійснення планування політики безпеки та приватності й пов'язані з ними заходи.\\ 
b. Призначити [Призначення: визначену організацією посадову особу] для управління політикою та процедурами планування політики безпеки та приватності.\\ 
c. Переглядати та оновлювати поточне планування:\\ 
1. політики планування безпеки та приватності [Призначення:з визначеною організацією частотою] та наступні [Призначення: визначені організацією події];\\ 
2. поточні процедури планування політики [Призначення: з визначеною організацією [Призначення: визначені організацією події]. безпеки та приватності частотою] та наступні \\ \hline
131 & ПЛАНИ ЗАХИСТУ ІНФОРМАЦІЇ ТА ПЕРСОНАЛЬНИХ ДАНИХ (PL-2) & Параметр: Призначені особи або групи, з якими пов'язана діяльність з безпекою та конфіденційністю, що впливає на систему, яка потребує планування та координації\newline Тип: list\newline Значення: \textbf{admin, security\_\allowbreak{}officer}\newline\newline Параметр: Призначено персонал або ролі для отримання розповсюджуваних копій планів захисту інформації та конфіденційності системи\newline Тип: list\newline Значення: \textbf{admin, security\_\allowbreak{}officer}\newline\newline Параметр: Визначено періодичність перегляду інформації та конфіденційності системи; планів захисту PL-02a.01[01] розроблено план захисту інформації, архітектурі підприємства організації; який відповідає PL-02a.01[02] розроблено план конфіденційності для системи, відповідає архітектурі підприємства організації; який PL-02a.02[01] розроблено план захисту інформації, який чітко визначає складові компоненти системи; PL-02a.02[02] розроблено для системи план забезпечення конфіденційності, який чітко визначає складові компоненти системи; PL-02a.03[01] розроблено план захисту інформації системи, який описує операційний контекст системи з точки зору місії та бізнеспроцесів; PL-02a.03[02] розроблено для системи план забезпечення конфіденційності, який описує операційний контекст системи з точки зору місії та бізнес-процесів; PL-02a.04[01] розроблено план захисту інформації, який визначає осіб, що виконують системні ролі та обов'язки; PL-02a.04[02] розроблено для системи план забезпечення конфіденційності, який визначає осіб, що виконують ролі та обов'язки в системі; PL-02a.05[01] розроблено план захисту інформації, який визначає типи інформації, що обробляється, зберігається та передається системою; PL-02a.05[02] розроблено план конфіденційності для системи, який визначає типи інформації, що обробляється, зберігається та передається системою; PL-02a.06[01] розроблено план захисту інформації, який передбачає категоризацію безпеки системи, включаючи відповідне обґрунтування; PL-02a.06[02] розроблено для системи план забезпечення конфіденційності, який передбачає категоризацію системи за рівнем безпеки, включаючи обґрунтування; PL-02a.07[01] розроблено план захисту інформації, який описує будь-які конкретні загрози для системи, що викликають потенційні ризики для рганізації; PL-02a.07[02] розроблено план конфіденційності для системи, який описує будь-які конкретні загрози для системи, що викликають потенційні ризики для організації; PL-02a.08[01] розроблено план захисту інформації, який містить результати оцінки ризиків конфіденційності для систем, що обробляють інформацію, яка ідентифікує особу; PL-02a.08[02] розроблено для системи план забезпечення конфіденційності, який містить результати оцінки ризиків конфіденційності для систем, що обробляють інформацію, яка ідентифікує особу; PL-02a.09[01] розроблено план захисту інформації, який описує операційне середовище системи та будь-які залежності або зв'язки з іншими системами чи компонентами системи; PL-02a.09[02] розроблено для системи план забезпечення конфіденційності, який описує робоче середовище системи та будь-які залежності або зв'язки з іншими системами чи компонентами системи; PL-02a.10[01] розроблено план захисту інформації, який містить огляд вимог до безпеки системи; PL-02a.10[02] розроблено для системи план забезпечення конфіденційності, який містить огляд вимог до конфіденційності системи; PL-02a.11[01] розроблено план захисту інформації, який визначає будь-які відповідні базові рівні контролю або обмеження, якщо такі є; PL-02a.11[02] розроблено для системи план забезпечення конфіденційності, який визначає будь-які відповідні базові рівні контролю або обмеження, якщо такі є; PL-02a.12[01] розроблено план захисту інформації, який описує наявні або заплановані засоби контролю для виконання вимог безпеки, включаючи обґрунтування будь-яких рішень щодо адаптації; PL-02a.12[02] розроблено для системи план забезпечення конфіденційності, який описує наявні або заплановані засоби контролю для виконання вимог щодо конфіденційності, включаючи обґрунтування будь-яких рішень, пов'язаних з адаптацією; PL-02a.13[01] розроблено план захисту інформації, який включає визначення ризиків для архітектури безпеки та проектних рішень; PL-02a.13[02] розроблено план забезпечення конфіденційності для системи, який включає визначення ризиків для архітектури конфіденційності та проектних рішень; PL-02a.14[01] розроблено захисту інформації, який включає діяльність, пов'язану з безпекою, що впливає на систему і потребує планування та координації з окремими особами або групами; PL-02a.14[02] розроблено для системи план забезпечення конфіденційності, який включає діяльність, пов'язану з конфіденційністю, що впливає на систему і потребує планування та координації з окремими особами або групами; PL-02a.15[01] розроблено план захисту інформації, який розглядається та затверджується уповноваженою посадовою особою або призначеним представником до початку реалізації плану; PL-02a.15[02] розроблено план забезпечення конфіденційності для системи, який перевіряється та затверджується уповноваженою посадовою особою або призначеним представником перед впровадженням плану. PL-02b.[01] розповсюджуються копії персонал або ролі>; PL-02b.[02] повідомляються наступні зміни до планів персонал або ролі; PL-02c. переглядаються плани частота; планів серед <PL-02\_\allowbreak{}ODP[02] PL-02d.[01] оновлюються плани відповідно до змін у системі та середовищі діяльності; PL-02d.[02] оновлюються плани для вирішення проблем, виявлених під час реалізації плану; PL-02d.[03] оновлюються плани для вирішення проблем, виявлених під час контрольних оцінок; PL-02e.[01] захищені плани від несанкціонованого розголошення; PL-02e.[02] захищені плани від несанкціонованої модифікації\newline Тип: list\newline Значення: \textbf{admin, security\_\allowbreak{}officer} & PL-2 & a. Розробити план захисту інформації та персональних даних для інформаційної системи, який:\\ 
1. узгоджується з архітектурою підприємства організації;\\ 
2. чітко визначає складові компоненти системи;\\ 
3. описує оперативний контекст інформаційної системи з точки зору завдань та процесів;\\ 
4. визначає осіб, які виконують системні ролі та обов'язки;\\ 
5. визначає тип інформації, яка обробляється, зберігається та передається системою;\\ 
6. надає огляд вимог безпеки та приватності інформаційної системи;\\ 
7. описує будь які конкретні загрози системи, які викликають стурбованість організації;\\ 
8. надає результати оцінки ризику конфіденційності для систем, в яких обробляються персональні дані;\\ 
9. описує робоче середовище інформаційної системи та будь які залежності від систем або компонентів систем або підключень до таких систем та їх компонентів;\\ 
10. надає огляд вимог безпеки та конфіденційності системи;\\ 
11. визначає будь які відповідні контрольні базові рівні або накладання, якщо вони застосовуються;\\ 
12. описує чинні або заплановані заходи щодо забезпечення безпеки та приватності, включно з обґрунтуванням рішень щодо налаштування\\ 
13. включає виявлення ризиків для архітектури безпеки і приватності, а також проєктних рішень;\\ 
14. включає дії, пов'язані з безпекою та конфіденційністю, які впливають на систему, виконання яких вимагає планування та координацію з [Призначення: визначені організацією окремі особи або групи];\\ 
15. розглядається та затверджується уповноваженою посадовою особою або призначеним представником до початку реалізації плану.\\ 
b. Поширити копії планів захисту інформації та персональних даних і повідомляти про подальші зміни планів серед [Призначення:визначеного організацією персоналу або ролей].\\ 
c. Переглядати плани захисту інформації та персональних даних [Призначення: з визначеною організацією частотою].\\ 
d. Оновлювати плани захисту інформації та персональних даних для врахування змін в інформаційній системі й робочому середовищі або проблем, виявлених у ході реалізації або оцінювання заходів безпеки та приватності.\\ 
e. забезпечити захист планів захисту інформації та персональних даних від несанкціонованого розголошення та змін. \\ \hline
132 & ПРАВИЛА ПОВЕДІНКИ (PL-4) & Параметр: Визначено періодичність перегляду та оновлення правил поведінки\newline Тип: string\newline Значення: \textbf{щорічно}\newline\newline Параметр: Вибрано одне або декілька з наступних ЗНАЧЕНЬ ПАРАМЕТРА:\{частота; коли правила переглядаються або оновлюються\}\newline Тип: list\newline Значення: \textbf{default\_\allowbreak{}deny\_\allowbreak{}rule, abac\_\allowbreak{}rule\_\allowbreak{}1}\newline\newline Параметр: Визначена періодичність перегляду та повторного підтвердження правил поведінки (якщо вибрано); PL-04a.[01] встановлені правила, які описують обов'язки та очікувану поведінку щодо використання інформації та системи, безпеки та конфіденційності для осіб, яким потрібен доступ до системи; PL-04a.[02] надаються правила, які описують обов'язки та очікувану поведінку щодо використання інформації та системи, безпеки та конфіденційності, особам, які отримують доступ до системи; PL-04b. отримано перед наданням доступу до інформації та системи задокументоване підтвердження від таких осіб про те, що вони прочитали, зрозуміли та згодні дотримуватися правил поведінки; PL-04c. переглядаються та оновлюються правила поведінки частота; PL-04d. потрібно особам, які визнали попередню версію правил поведінки, прочитати та повторно визнати ЗНАЧЕННЯ ВИБРАНОГО ПАРАМЕТРА(ів)\newline Тип: list\newline Значення: \textbf{admin, security\_\allowbreak{}officer} & PL-4 & a. Створити та надати особам, що потребують доступу до інформаційної системи, правила, які описують їхні обов'язки й очікувану поведінку щодо інформації та використання інформаційної системи, безпеки та приватності.\\ 
b. Отримати документальне підтвердження від таких осіб про те, що вони прочитали, зрозуміли та погодилися дотримуватися правил поведінки, перш ніж дозволяти доступ до інформації та інформаційної системи.\\ 
c. Переглядати й оновлювати правила поведінки [Призначення: з визначеною організацією частотою].\\ 
d. Вимагати від осіб, які підписали попередню версію правил поведінки, перечитати та повторно підписати правила [Вибір (один або декілька): [Призначення: з визначеною організацією частотою]; коли правила переглядаються чи оновлюються]. \\ \hline
\multicolumn{5}{|l|}{\textbf{13. БЕЗПЕКА ПЕРСОНАЛУ (PS) (PS)}} \\ \hline
133 & Політика та процедури кадрової безпеки (PS-1) & Параметр: Визначено персонал або ролі, на які поширюється політика кадрової безпеки\newline Тип: list\newline Значення: \textbf{admin, security\_\allowbreak{}officer}\newline\newline Параметр: Визначено персонал або ролі, на які поширюються процедури кадрової безпеки\newline Тип: list\newline Значення: \textbf{admin, security\_\allowbreak{}officer}\newline\newline Параметр: Визначено посадову особу, яка керуватиме політикою та процедурами кадрової безпеки\newline Тип: list\newline Значення: \textbf{admin, security\_\allowbreak{}officer}\newline\newline Параметр: Визначена періодичність перегляду та оновлення поточної політики кадрової безпеки\newline Тип: list\newline Значення: \textbf{default\_\allowbreak{}deny\_\allowbreak{}rule, abac\_\allowbreak{}rule\_\allowbreak{}1}\newline\newline Параметр: Є події, які вимагають перегляду та оновлення поточної політики кадрової безпеки\newline Тип: list\newline Значення: \textbf{default\_\allowbreak{}deny\_\allowbreak{}rule, abac\_\allowbreak{}rule\_\allowbreak{}1}\newline\newline Параметр: Визначено періодичність перегляду та оновлення поточних процедур кадрової безпеки\newline Тип: string\newline Значення: \textbf{щорічно} & PS-1 &  \\ \hline
134 & Перевірка персоналу (PS-3) & Параметр: Визначені умови, що вимагають повторної перевірки осіб\newline Тип: list\newline Значення: \textbf{}\newline\newline Параметр: Визначена частота повторної перевірки осіб, для яких це показано; PS-03a. проходять особи перевірку перед тим, як надати їм доступ до системи; PS-03b.[01] проходять особи повторну перевірку відповідно до умови, що вимагають повторної перевірки; PS-03b.[02] проводиться повторна перевірка у випадках, коли це зазначено, частота\newline Тип: list\newline Значення: \textbf{admin, security\_\allowbreak{}officer} & PS-3 &  \\ \hline
135 & Звільнення персоналу (PS-4) & Параметр: Визначено період часу, протягом якого забороняється доступ до системи\newline Тип: integer\newline Значення: \textbf{30}\newline\newline Параметр: Визначені теми інформаційної безпеки для обговорення під час проведення співбесід; PS-04a. при звільненні працівника доступ до системи вимикається протягом часового періоду; PS-04b. припиняють дію або анулюють будь-які автентифікатори та облікові дані після припинення трудових відносин з окремими особами; PS-04c. проводяться при звільненні окремих працівників співбесіди, які включають обговорення питань інформаційної безпеки; PS-04d. отримується після звільнення особи все майно, пов'язане з безпекою організаційної системи; PS-04e. зберігається доступ до організаційної інформації та систем, які раніше перебували під контролем особи, що звільняється, після припинення нею трудових відносин\newline Тип: list\newline Значення: \textbf{admin, security\_\allowbreak{}officer} & PS-4 &  \\ \hline
136 & Переведення персоналу (PS-5) & Параметр: Визначені дії, які мають бути ініційовані після переведення або перепризначення\newline Тип: list\newline Значення: \textbf{login, logout, failed\_\allowbreak{}attempt}\newline\newline Параметр: Визначено період часу, протягом якого мають бути здійснені дії з переведення або перепризначення після переведення або перепризначення\newline Тип: list\newline Значення: \textbf{login, logout, failed\_\allowbreak{}attempt}\newline\newline Параметр: Визначено персонал або ролі, про які необхідно повідомляти, коли осіб призначають на інші посади або переводять на інші посади в організації\newline Тип: list\newline Значення: \textbf{admin, security\_\allowbreak{}officer}\newline\newline Параметр: Визначено період часу, протягом якого необхідно повідомляти визначений організацією персонал або ролі, коли осіб перепризначають або переводять на інші посади в межах організації; PS-05a. переглядаються та підтверджуються поточні потреби в логічних та фізичних дозволах на доступ до систем та об'єктів при перепризначенні або переведенні осіб на інші посади в організації; PS-05b. були ініційовані дії з переведення або перепризначення протягом періоду часу після формальної дії з переведення; PS-05c. змінюється авторизація доступу за необхідності, щоб відповідати будь-яким змінам в оперативних потребах у зв'язку з перепризначенням або переведенням; PS-05d. було повідомлено персонал або ролі протягом часового періоду\newline Тип: list\newline Значення: \textbf{admin, security\_\allowbreak{}officer} & PS-5 &  \\ \hline
\multicolumn{5}{|l|}{\textbf{14. ОЦІНКА РИЗИКІВ (RA) (RA)}} \\ \hline
137 & ПОЛІТИКА ТА ПРОЦЕДУРИ ОЦІНЮВАННЯ РИЗИКУ (RA-1) & Параметр: Визначено персонал або ролі, на які поширюється політика оцінювання ризику\newline Тип: list\newline Значення: \textbf{admin, security\_\allowbreak{}officer}\newline\newline Параметр: Визначено персонал або ролі, на які поширюються процедури, що сприяють здійсненню політики оцінювання ризику\newline Тип: list\newline Значення: \textbf{admin, security\_\allowbreak{}officer}\newline\newline Параметр: Визначена посадова особа, відповідальна за управління політикою та процедурами оцінювання ризику\newline Тип: list\newline Значення: \textbf{admin, security\_\allowbreak{}officer}\newline\newline Параметр: Визначена періодичність перегляду та оновлення поточної політики оцінювання ризику\newline Тип: list\newline Значення: \textbf{default\_\allowbreak{}deny\_\allowbreak{}rule, abac\_\allowbreak{}rule\_\allowbreak{}1}\newline\newline Параметр: Є події, які вимагають перегляду та оновлення поточної політики оцінювання ризику\newline Тип: list\newline Значення: \textbf{default\_\allowbreak{}deny\_\allowbreak{}rule, abac\_\allowbreak{}rule\_\allowbreak{}1}\newline\newline Параметр: Визначено періодичність перегляду поточних процедур оцінювання ризику\newline Тип: string\newline Значення: \textbf{щорічно} & RA-1 & a. Розробити, задокументувати та поширити серед [Призначення: визначеного організацією персоналу або посадових осіб]:\\ 
1. 2. Політику оцінювання ризику, яка:\\ 
(a) містить мету, сферу застосування, ролі, обов'язки, відповідальність керівництва, координацію між організаційними підрозділами та систему контролю відповідності (complaince);\\ 
(b) відповідає чинному законодавству, виконавчим наказам, директивам, нормам, політикам, стандартам і керівним принципам. Процедури, що сприяють здійсненню політики оцінювання ризику та пов'язаних з ними заходів оцінювання ризику.\\ 
b. Призначити [Призначення: визначена організацією посадову особу] для управління політикою та процедурами оцінювання ризику.\\ 
c. Переглядати й оновлювати:\\ 
1. Поточну політику оцінювання організацією частотою]. ризику [Призначення: з визначеною\\ 
2. Поточні процедури оцінювання ризику [Призначення: з визначеною організацією частотою]. \\ \hline
138 & ОЦІНЮВАННЯ РИЗИКУ (RA-3) & Параметр: Визначено періодичність оцінювання ризику\newline Тип: string\newline Значення: \textbf{щорічно}\newline\newline Параметр: Визначено персонал або ролі, до яких мають бути доведені результати оцінювання ризику\newline Тип: list\newline Значення: \textbf{admin, security\_\allowbreak{}officer}\newline\newline Параметр: Визначена періодичність оновлення оцінювання ризику; перегляду результатів RA-03a.01 проводиться оцінювання ризику для виявлення загроз і вразливостей у системі; RA-03a.02 проводиться оцінювання ризику для визначення ймовірності та розміру шкоди від несанкціонованого доступу, використання, розкриття, порушення, модифікації або знищення системи; інформації, яку вона обробляє, зберігає або передає; а також будь-якої пов'язаної з нею інформації; RA-03a.03 проводиться оцінювання ризику для визначення ймовірності та впливу несприятливих наслідків для фізичних осіб, що виникають у зв'язку з обробкою інформації, яка ідентифікує особу; RA-03b. інтегровані результати оцінювання ризику та рішення з управління ризиками з точки зору організації та місії або бізнес-процесів з оцінкою ризиків на системному рівні; RA-03c. результати оцінювання ризику задокументовані в ЗНАЧЕННЯ ВИБРАНОГО ПАРАМЕТРА; RA-03d. переглядаються результати 03\_\allowbreak{}ODP[03] частота>; RA-03e. поширюються результати оцінювання ризику серед персоналу або ролей; RA-03f. оновлюється оцінювання ризику частота або коли відбуваються значні зміни в системі, середовищі її функціонування або інших умовах, які можуть вплинути на стан безпеки або приватності системи оцінювання ризику <RA-\newline Тип: list\newline Значення: \textbf{admin, security\_\allowbreak{}officer} & RA-3 &  \\ \hline
139 & ОЦІНЮВАННЯ ПОСТАЧАННЯ (RA-3(1)) &  & RA-3(1) &  \\ \hline
140 & СКАНУВАННЯ ВРАЗЛИВОСТЕЙ (RA-5) & Параметр: Визначена періодичність перевірки систем та розміщених на них застосунків на наявність вразливостей\newline Тип: string\newline Значення: \textbf{щорічно}\newline\newline Параметр: Визначено час реагування на усунення законних вразливостей відповідно до організаційної оцінення ризику\newline Тип: integer\newline Значення: \textbf{30}\newline\newline Параметр: Потрібно ділитися інформацією, отриманою в процесі сканування вразливостей та оцінок контролю, з персоналом або ролями, з якими потрібно ділитися; RA-05a.[01] здійснюється моніторинг систем та розміщених застосунків на наявність вразливостей частота та/або випадковість відповідно до визначеного організацією процесу, а також коли виявляються та повідомляються нові вразливості, що потенційно можуть вплинути на систему; RA-05a.[02] перевіряються системи та розміщені застосунки на наявність вразливостей частота та/або випадковим чином відповідно до визначеного організацією процесу, а також коли виявляються та повідомляються нові вразливості, що потенційно можуть вплинути на систему; RA-05b. застосовуються інструменти та методи моніторингу вразливостей для забезпечення сумісності між інструментами; RA-05b.01 застосовуються інструменти та методи моніторингу вразливостей для автоматизації частини процесу управління вразливостями, використовуючи стандарти для переліку платформ, недоліків програмного забезпечення та неправильних конфігурацій; RA-05b.02 застосовуються інструменти та методи моніторингу вразливостей для полегшення взаємодії між інструментами та автоматизації частини процесу управління вразливостями шляхом використання стандартів для формування контрольних списків та процедур тестування; RA-05b.03 застосовуються інструменти та методи моніторингу вразливостей для полегшення взаємодії між інструментами та автоматизації частин процесу управління вразливостями шляхом використання стандартів для вимірювання впливу вразливостей; RA-05c. аналізуються звіти про сканування вразливостей та результати моніторингу вразливостей; RA-05d. усуваються легітимні вразливості час реагування відповідно до організаційної оцінки ризиків; RA-05e. надається інформація, отримана в процесі моніторингу вразливостей та оцінки контролю, персонал або ролі, щоб допомогти усунути подібні вразливості в інших системах; RA-05f. використовуються інструменти моніторингу вразливостей, які передбачають можливість швидкого оновлення вразливостей, що підлягають скануванню\newline Тип: list\newline Значення: \textbf{admin, security\_\allowbreak{}officer} & RA-5 & Використовувати заходи ПДТР за [Призначення: визначені організацією місця] [Вибір (один або кілька): [Призначення: з визначеною організацією частотою]; [Призначення: за визначеними організацією подіями або показниками]]. \\ \hline
141 & СКАНУВАННЯ ВРАЗЛИВОСТЕЙ - ОНОВЛЕННЯ ЗА ЧАСТОТОЮ, ПЕРЕД НОВИМ СКАНУВАННЯМ АБО ПРИ ІДЕНТИФІКАЦІЇ (RA-5(2)) & Параметр: Визначено час реагування на усунення законних вразливостей відповідно до організаційної оцінення ризику\newline Тип: integer\newline Значення: \textbf{30} & RA-5(2) &  \\ \hline
142 & РЕАГУВАННЯ НА РИЗИК (RA-7) &  & RA-7 & Реагувати на результати оцінювання, моніторингу й аудиту безпеки та приватності. \\ \hline
\multicolumn{5}{|l|}{\textbf{15. ПРИДБАННЯ СИСТЕМ ТА ПОСЛУГ (SA) (SA)}} \\ \hline
143 & ПОЛІТИКИ ТА ПРОЦЕДУРИ ПРИДБАННЯ СИСТЕМ ТА ПОСЛУГ (SA-1) & Параметр: Визначено персонал або ролі, на які поширюватиметься політика придбання систем і послуг\newline Тип: list\newline Значення: \textbf{admin, security\_\allowbreak{}officer}\newline\newline Параметр: Визначено персонал або ролі, на які поширюються процедури придбання систем і послуг\newline Тип: list\newline Значення: \textbf{admin, security\_\allowbreak{}officer}\newline\newline Параметр: Визначено посадову особу, яка керуватиме політикою та процедурами придбання систем і послуг\newline Тип: list\newline Значення: \textbf{admin, security\_\allowbreak{}officer}\newline\newline Параметр: Визначено періодичність перегляду та оновлення поточної політики придбання систем і послуг\newline Тип: list\newline Значення: \textbf{default\_\allowbreak{}deny\_\allowbreak{}rule, abac\_\allowbreak{}rule\_\allowbreak{}1}\newline\newline Параметр: Є події, які вимагають перегляду та оновлення поточної політики придбання систем і послуг\newline Тип: list\newline Значення: \textbf{default\_\allowbreak{}deny\_\allowbreak{}rule, abac\_\allowbreak{}rule\_\allowbreak{}1}\newline\newline Параметр: Визначено частоту, з якою переглядаються та оновлюються поточні процедури придбання систем і послуг\newline Тип: integer\newline Значення: \textbf{30} & SA-1 & a. Розробити, задокументувати та поширити серед [Призначення: визначеного організацією персоналу або ролей]:\\ 
1. 2.\\ 
b. [Вибір (один або декілька): рівень організації; рівень місії/бізнес-процесу; рівень системи] політики придбання систем і послуг, яка:\\ 
(a) містить мету, сферу застосування, ролі, обов'язки, відповідальність керівництва, координацію між організаційними підрозділами та систему контролю (complaince);\\ 
(b) відповідає чинному законодавству, виконавчим наказам, директивам, нормам, політикам, стандартам і рекомендаціям. Процедури, що полегшують впровадження політики та заходів придбання систем і послуг. Призначити [Призначення: визначену організацією посадову особу] для управління політикою та процедурами придбання системи та послуг.\\ 
c. Переглядати й оновлювати поточні політику та процедури придбання систем та послуг:\\ 
1. Поточну політику придбання системи та послуг [Призначення: з визначеною організацією частотою].\\ 
2. Поточні процедури придбання системи та послуг [Призначення: з визначеною організацією частотою] та наступні [Призначення: події, визначені організацією]. \\ \hline
144 & БЕЗПЕКА ТА ПРИВАТНІСТЬ ПРИНЦИПІВ ІНЖИНІРИНГУ (SA-8) &  & SA-8 & Застосовувати [Призначення: визначені організацією принципи інжинірингу безпеки та конфіденційності системи] в специфікації, проєктуванні, розробці, впровадженні та зміні системи й компонентів системи. \\ \hline
145 & ЗОВНІШНІ ПОСЛУГИ ДЛЯ СИСТЕМИ (SA-9) & Параметр: Визначені засоби контролю, які будуть застосовуватися зовнішніми постачальниками системних послуг\newline Тип: string\newline Значення: \textbf{автоматизований засіб моніторингу}\newline\newline Параметр: Визначені процеси, методи та техніки, що застосовуються для моніторингу дотримання вимог контролю зовнішніми постачальниками послуг; SA-09a.[01] дотримуються постачальники зовнішніх системних послуг вимог організаційної безпеки; SA-09a.[02] дотримуються постачальники зовнішніх системних послуг вимог приватності організації; SA-09a.[03] використовують постачальники зовнішніх системних послуг елементи керування; SA-09b.[01] визначено та задокументовано організаційний нагляд за зовнішніми системними послугами; SA-09b.[02] визначені та задокументовані ролі та обов'язки користувачів щодо зовнішніх системних сервісів; SA-09c. визначені та задокументовані ролі та обов'язки користувачів щодо зовнішніх системних послуг\newline Тип: list\newline Значення: \textbf{admin, security\_\allowbreak{}officer} & SA-9 & a. Вимагати, щоб постачальники зовнішніх послуг для системи відповідали вимогам безпеки та приватності в організації та застосовували такі заходи захисту [Призначення: встановлені організацією заходи безпеки та приватності].\\ 
b. Визначити та задокументувати нагляд організаціїповноваження та обов'язки користувачів щодо зовнішніх послуг для системи.\\ 
c. Використовувати наступні процеси, методи та техніки для постійного моніторингу дотримання контролю зовнішніми постачальниками послуг: [Призначення: визначені організацією процеси, методи та техніки]. \\ \hline
146 & КОМПОНЕНТИ СИСТЕМИ, ЩО НЕ ПІДТРИМУЮТЬСЯ (SA-22) &  & SA-22 & a. Замінювати компоненти системи, якщо підтримка компонентів більше не доступна розробнику, постачальнику або виробнику.\\ 
b. Надавати такі варіанти альтернативних джерел для подальшої підтримки непідтримуваних компонентів [Вибір (один або більше): внутрішня підтримка; [Призначення: підтримка, визначена організацією від зовнішніх постачальників]]. \\ \hline
\multicolumn{5}{|l|}{\textbf{16. ЗАХИСТ СИСТЕМ ТА КОМУНІКАЦІЙ (SC) (SC)}} \\ \hline
147 & ПОЛІТИКА ТА ПРОЦЕДУРИ ЗАХИСТУ СИСТЕМИ ТА КОМУНІКАЦІЙ (SC-1) & Параметр: Визначено персонал або ролі, до яких має бути доведена політика захисту системи та комунікацій\newline Тип: list\newline Значення: \textbf{admin, security\_\allowbreak{}officer}\newline\newline Параметр: Визначено персонал або ролі, на які поширюються процедури захисту системи та комунікацій\newline Тип: list\newline Значення: \textbf{admin, security\_\allowbreak{}officer}\newline\newline Параметр: Визначено посадову особу, яка керуватиме політикою та процедурами захисту системи та комунікацій\newline Тип: list\newline Значення: \textbf{admin, security\_\allowbreak{}officer}\newline\newline Параметр: Визначена періодичність перегляду та оновлення поточної політики захисту системи та комунікацій\newline Тип: list\newline Значення: \textbf{default\_\allowbreak{}deny\_\allowbreak{}rule, abac\_\allowbreak{}rule\_\allowbreak{}1}\newline\newline Параметр: Є події, які вимагають перегляду та оновлення поточної політики захисту системи та комунікацій\newline Тип: list\newline Значення: \textbf{default\_\allowbreak{}deny\_\allowbreak{}rule, abac\_\allowbreak{}rule\_\allowbreak{}1}\newline\newline Параметр: Визначена періодичність перегляду та оновлення поточних процедур захисту системи та засобів зв'язку\newline Тип: string\newline Значення: \textbf{щорічно} & SC-1 & a. Розробити, задокументувати та поширити серед [Призначення: визначеного організацією персоналу або посадових осіб]:\\ 
1. 2. Політику захисту системи та комунікацій, яка:\\ 
(a) містить мету, сферу застосування, ролі, обов'язки, відповідальність керівництва, координацію між організаційними підрозділами та систему контролю відповідності (complaince);\\ 
(b) відповідає чинному законодавству, виконавчим наказам, директивам, нормам, політикам, стандартам та керівним принципам. Процедури для сприяння впровадженню політики в області захисту систем і комунікацій, а також пов'язаних з ними систем і засобів захисту зв'язку.\\ 
b. Призначити [Призначення: визначена організацією посадову особу] для управління політикою та процедурами захисту системи та комунікацій.\\ 
c. Переглядати та оновлювати:\\ 
1. поточну політику захисту системи та комунікацій [Призначення: визначена організацією частота];\\ 
2. поточні процедури захисту системи та комунікацій [Призначення: визначена організацією частота]. \\ \hline
148 & ІНФОРМАЦІЯ В ЗАГАЛЬНИХ СИСТЕМНИХ РЕСУРСАХ (SC-4) &  & SC-4 & Запобігати несанкціонованій та ненавмисній передачі інформації через спільні системні ресурси. \\ \hline
149 & ДОСТУПНІСТЬ РЕСУРСІВ (SC-7) &  & SC-7 & a. Контролювати та управляти зв'язком на зовнішньому периметрі системи та на ключових внутрішніх периметрах всередині системи.\\ 
b. Реалізувати підмережі для загальнодоступних компонентів системи, які є [Вибір: фізично; логічно] відділені від внутрішніх мереж організації.\\ 
c. Підключатися до зовнішніх мереж або систем тільки через керовані інтерфейси, що складаються з пристроїв захисту периметру, і розташованих відповідно до архітектури безпеки та приватності організації. \\ \hline
150 & ЗАХИСТ ПЕРИМЕТРА - ВІДМОВА ЗА ЗАМОВЧУВАННЯМ - ДОЗВІЛ ЗА ВИНЯТКОМ (SC-7(5)) &  & SC-7(5) &  \\ \hline
151 & КОНФІДЕНЦІЙНІСТЬ ТА ЦІЛІСНІСТЬ ПЕРЕДАЧІ (SC-8) &  & SC-8 & Забезпечити [Вибір (один або кілька): конфіденційність; цілісність] інформації, що передається. \\ \hline
152 & КОНФІДЕНЦІЙНІСТЬ ТА КРИПТОГРАФІЧНИЙ ЗАХИСТ (SC-8(1)) &  & SC-8(1) &  \\ \hline
153 & ВІДКЛЮЧЕННЯ МЕРЕЖІ (SC-10) & Параметр: Визначено період бездіяльності, після якого система розриває мережеве з'єднання, пов'язане з сеансом зв'язку; SC08(05)\_\allowbreak{}ODP[02] мережеве з'єднання, пов'язане з сеансом зв'язку, розірвано в кінці сеансу або після період часу бездіяльності\newline Тип: integer\newline Значення: \textbf{30} & SC-10 & Завершити з'єднання з мережею, яке пов'язане із сеансом зв'язку в кінці сеансу або після [Призначення: визначений організацією період часу] бездіяльності. \\ \hline
154 & ВСТАНОВЛЕННЯ КЛЮЧАМИ (SC-12) & Параметр: Встановлюються криптографічні ключі, коли в системі використовується криптографія відповідно до < SC-12\_\allowbreak{}ODP вимог >\newline Тип: string\newline Значення: \textbf{AES-256-GCM}\newline\newline Параметр: Здійснюється управління криптографічними ключами, коли в системі використовується криптографія, відповідно до вимог \newline Тип: string\newline Значення: \textbf{AES-256-GCM} & SC-12 & Встановити та управляти криптографічними ключами для криптографічних засобів, які використовуються в системі відповідно до [Призначення: визначені організацією вимоги до генерації, поширення, зберігання, доступу та знищення ключів]. \\ \hline
155 & КРИПТОГРАФІЧНИЙ ЗАХИСТ (SC-13) & Параметр: КРИПТОГРАФІЧНИЙ ЗАХИСТ МЕТА ОЦІНКИ: Визначити, чи:\newline Тип: string\newline Значення: \textbf{AES-256-GCM}\newline\newline Параметр: Визначено використання криптографічних засобів\newline Тип: string\newline Значення: \textbf{AES-256-GCM}\newline\newline Параметр: Визначено типи криптографії для кожного вказаного криптографічного використання; SC-13a. ідентифіковано використання>; SC-13b. реалізовано типи криптографії для кожного вказаного криптографічного використання (визначеного в SC-13\_\allowbreak{}ODP[01]). криптографічне <SC-13\_\allowbreak{}ODP[01]\newline Тип: string\newline Значення: \textbf{AES-256-GCM} & SC-13 & a. Визначити [Призначення: використання криптографічних засобів, визначених організацією];\\ 
b. Впровадити [Завдання: визначені організацією види криптографії для кожного визначеного криптографічного використання]. \\ \hline
156 & СПІЛЬНІ ОБЧИСЛЮВАЛЬНІ ПРИСТРОЇ ТА ЗАСТОСУНКИ (SC-15) &  & SC-15 & a. Заборонити віддалену активацію спільних обчислювальних пристроїв (хмар) та застосунків з такими виключеннями: [Призначення: визначені організацією виключення, у яких дозволена віддалена активація].\\ 
b. Надати явну вказівку щодо використання користувачами фізично присутніми пристроями. \\ \hline
157 & СЕРТИФІКАТИ ІНФРАСТРУКТУРИ ВІДКРИТИХ КЛЮЧІВ (SC-17) &  & SC-17 & a. Випускати сертифікати відкритого ключа відповідно до [Призначення: визначеної організацією політики сертифікації];\\ 
b. Отримувати сертифікати відкритого ключа від затвердженого постачальника послуг. \\ \hline
158 & МОБІЛЬНИЙ КОД (SC-18) &  & SC-18 & a. Визначати прийнятні та неприйнятні мобільні коди та технології мобільних кодів.\\ 
b. Проводити авторизацію, відстежувати та контролювати використання мобільного коду всередині системи. \\ \hline
159 & АВТЕНТИФІКАЦІЯ СЕСІЇ (SC-23) &  & SC-23 & Забезпечити автентифікацію сеансів зв'язку. \\ \hline
160 & ЗАХИСТ ІНФОРМАЦІЇ В СТАНІ СПОКОЮ (SC-28) &  & SC-28 & Забезпечити [Вибір (один або кілька): конфіденційність; цілісність] [Призначення: визначена організацією інформація] в стані спокою. \\ \hline
161 & ЗАХИСТ ІНФОРМАЦІЇ В СТАНІ СПОКОЮ | КРИПТОГРАФІЧНИЙ ЗАХИСТ & Параметр: Реалізовані криптографічні механізми для запобігання несанкціонованому розкриттю інформації, що знаходиться в стані спокою на системних компонентах або носіях\newline Тип: string\newline Значення: \textbf{AES-256-GCM}\newline\newline Параметр: Реалізовані криптографічні механізми для запобігання несанкціонованій модифікації інформації, що знаходиться в стані спокою на системних компонентах або носіях\newline Тип: string\newline Значення: \textbf{AES-256-GCM} & SC-28(1) &  \\ \hline
162 & ЕКРАНОВАНІ КАМЕРИ (SC-44) &  & SC-44 & Впровадити екрановані камери в [Призначення: визначену організацією систему, компонент системи або місце розташування]. \\ \hline
163 & ПРИМУСОВЕ АПАРАТНЕ ЗАБЕЗПЕЧЕННЯ ВИКОНАННЯ (SC-49) & Параметр: Впроваджено механізми апаратного розділення та застосування політик між доменами безпеки\newline Тип: list\newline Значення: \textbf{default\_\allowbreak{}deny\_\allowbreak{}rule, abac\_\allowbreak{}rule\_\allowbreak{}1}\newline\newline Параметр: Визначені домени безпеки, які потребують апаратного розділення та механізмів забезпечення дотримання політики\newline Тип: list\newline Значення: \textbf{default\_\allowbreak{}deny\_\allowbreak{}rule, abac\_\allowbreak{}rule\_\allowbreak{}1} & SC-49 & Впровадити механізми апаратного поділу та застосування політики між [Призначення: домени безпеки, визначені організацією]. \\ \hline
164 & АПАРАТНИЙ ЗАХИСТ (SC-51) & Параметр: Визначені уповноважені особи, які повинні виконувати процедури вимкнення та повторного ввімкнення апаратного захисту від запису; SC-51a. використовується апаратний захист від запису для компонентів мікропрограми системи; SC-51b.[01] впроваджено спеціальні процедури для уповноважених осіб для ручного вимкнення апаратного захисту від запису для модифікацій мікропрограми; SC-51b.[02] реалізовано спеціальні процедури для уповноважених осіб для повторного увімкнення захисту від запису перед поверненням до робочого режиму\newline Тип: list\newline Значення: \textbf{admin, security\_\allowbreak{}officer} & SC-51 & a. Перевіряти правильність роботи [Призначення: визначені організацією функції безпеки та приватності].\\ 
b. Виконувати перевірку [Вибір (один або кілька): [Призначення: визначені організацією системні перехідні стани]; за командою користувача з відповідними повноваженнями; [Призначення: визначена організацією частота]].\\ 
c. Повідомляти [Призначення: визначені організацією персонал або посадові особи] про невдалі перевірки безпеки та приватності.\\ 
d. [Вибір (один або кілька): Вимкнути систему; Перезапустити систему; [Призначення: визначені організацією альтернативні дії]], коли виявляються аномалії. \\ \hline
\multicolumn{5}{|l|}{\textbf{17. ЦІЛІСНІСТЬ СИСТЕМИ ТА ІНФОРМАЦІЇ (SI) (SI)}} \\ \hline
165 & ПОЛІТИКА І ПРОЦЕДУРИ ЦІЛІСНОСТІ ІНФОРМАЦІЇ (SI-1) & Параметр: Визначено персонал або ролі, до яких має бути доведена політика цілісності системи та інформації\newline Тип: list\newline Значення: \textbf{admin, security\_\allowbreak{}officer}\newline\newline Параметр: Визначено персонал або ролі, на які поширюються процедури цілісності системи та інформації \newline Тип: list\newline Значення: \textbf{admin, security\_\allowbreak{}officer}\newline\newline Параметр: Визначено посадову особу, відповідальну за управління системою та політикою і процедурами цілісності інформації\newline Тип: list\newline Значення: \textbf{admin, security\_\allowbreak{}officer}\newline\newline Параметр: Визначено періодичність перегляду та оновлення поточної політики цілісності системи та інформації\newline Тип: list\newline Значення: \textbf{default\_\allowbreak{}deny\_\allowbreak{}rule, abac\_\allowbreak{}rule\_\allowbreak{}1}\newline\newline Параметр: Є події, які вимагають перегляду та оновлення поточної політики цілісності системи та інформації\newline Тип: list\newline Значення: \textbf{default\_\allowbreak{}deny\_\allowbreak{}rule, abac\_\allowbreak{}rule\_\allowbreak{}1}\newline\newline Параметр: Визначено частоту, з якою переглядаються та оновлюються поточні цілісності системи та інформації\newline Тип: integer\newline Значення: \textbf{30} & SI-1 &  \\ \hline
166 & ВИПРАВЛЕННЯ ДЕФЕКТІВ (SI-2) & Параметр: Визначено період часу, протягом якого необхідно встановити оновлення програмного забезпечення, пов'язані з безпекою, після виходу оновлень; SI-02a.[01] виявлено недоліки системи; SI-02a.[02] повідомляється про недоліки системи; SI-02a.[03] виправлені недоліки системи; SI-02b.[01] перевіряються оновлення програмного забезпечення, пов'язані з усуненням недоліків, на ефективність перед встановленням; SI-02b.[02] перевіряються оновлення програмного забезпечення, пов'язані з виправленням дефектів, на наявність потенційних побічних ефектів перед встановленням; SI-02b.[03] перевіряються оновлення прошивки, пов'язані з усуненням недоліків, на ефективність перед встановленням; SI-02b.[04] перевіряються оновлення прошивки, пов'язані з усуненням недоліків, на наявність потенційних побічних ефектів перед встановленням; SI-02c.[01] встановлено оновлення програмного забезпечення, що стосуються безпеки, протягом часовий проміжок з моменту випуску оновлень; SI-02c.[02] встановлено оновлення мікропрограми, що стосуються безпеки, протягом часового періоду з моменту випуску оновлень; SI-02d. включено відновлення порушених прав у процес управління організаційною конфігурацією\newline Тип: integer\newline Значення: \textbf{30} & SI-2 &  \\ \hline
167 & ЗАХИСТ ВІД ШКІДЛИВОГО КОДУ (SI-3) & Параметр: Визначено частоту, з якою механізми шкідливого коду виконують сканування\newline Тип: integer\newline Значення: \textbf{30}\newline\newline Параметр: Визначено дії, яких слід вжити у відповідь на виявлення шкідливого коду (якщо вибрано)\newline Тип: list\newline Значення: \textbf{login, logout, failed\_\allowbreak{}attempt} & SI-3 &  \\ \hline
168 & МОНІТОРИНГ СИСТЕМИ (SI-4) & Параметр: Визначені методи та способи, що використовуються для виявлення несанкціонованого використання системи\newline Тип: list\newline Значення: \textbf{admin, security\_\allowbreak{}officer}\newline\newline Параметр: Визначена інформація про моніторинг системи, яка повинна надаватися персоналу або ролям\newline Тип: list\newline Значення: \textbf{admin, security\_\allowbreak{}officer}\newline\newline Параметр: Визначено персонал або ролі, яким має надаватися інформація про моніторинг системи\newline Тип: list\newline Значення: \textbf{admin, security\_\allowbreak{}officer}\newline\newline Параметр: Вибрано одне або декілька з наступних ЗНАЧЕНЬ ПАРАМЕТРІВ: \{за потребою; частота\}\newline Тип: string\newline Значення: \textbf{щорічно} & SI-4 &  \\ \hline
169 & МОНІТОРИНГ СИСТЕМИ КОМУНІКАЦІЙ (SI-4(4)) & Параметр: Визначені критерії незвичної або несанкціонованої діяльності або умови для вхідного трафіку зв'язку\newline Тип: list\newline Значення: \textbf{}\newline\newline Параметр: Визначені критерії незвичайної або несанкціонованої діяльності або умови для вихідного трафіку зв'язку\newline Тип: list\newline Значення: \textbf{}\newline\newline Параметр: Здійснюється моніторинг вхідного комунікаційного трафіку частота на предмет незвичних або несанкціонованих дій або умов\newline Тип: string\newline Значення: \textbf{щорічно}\newline\newline Параметр: Контролюється вихідний трафік зв'язку частота на предмет незвичних або несанкціонованих дій або умов. вихідного виявлення\newline Тип: string\newline Значення: \textbf{щорічно} & SI-4(4) &  \\ \hline
170 & ПОПЕРЕДЖЕННЯ, РЕКОМЕНДАЦІЇ ТА ДИРЕКТИВИ З БЕЗПЕКИ (SI-5) & Параметр: Вибрано одне або декілька з наступних ЗНАЧЕНЬ ПАРАМЕТРІВ: \{персонал або ролі; елементи; зовнішні організації\}\newline Тип: list\newline Значення: \textbf{admin, security\_\allowbreak{}officer}\newline\newline Параметр: Визначено персонал або ролі, до яких мають бути доведені попередження, поради та директиви з безпеки (якщо визначено)\newline Тип: list\newline Значення: \textbf{admin, security\_\allowbreak{}officer} & SI-5 &  \\ \hline
171 & ЦІЛІСНІСТЬ ПРОГРАМНОГО ЗАБЕЗПЕЧЕННЯ, ПРОГРАМНОГО ЗАБЕЗПЕЧЕННЯ ТА ІНФОРМАЦІЇ (SI-7) & Параметр: Визначені дії, яких слід вжити при виявленні несанкціонованих змін у програмному забезпеченні\newline Тип: list\newline Значення: \textbf{login, logout, failed\_\allowbreak{}attempt}\newline\newline Параметр: Визначені дії, яких слід вжити несанкціонованих змін у прошивці; при виявленні SI-07\_\allowbreak{}ODP[06] визначені дії, яких слід вжити несанкціонованих змін до інформації; при виявленні SI-07a.[01] використовуються засоби перевірки цілісності для виявлення несанкціонованих змін у програмному забезпеченні; SI-07a.[02] використовуються засоби перевірки цілісності для виявлення несанкціонованих змін у мікропрограмі; SI-07a.[03] використовуються засоби перевірки цілісності для виявлення несанкціонованих змін до інформації; SI-07b.[01] виконуються дії при виявленні несанкціонованих змін у програмному забезпеченні; SI-07b.[02] виконуються дії несанкціонованих змін у прошивці; при виявленні SI-07b.[03] виконуються дії несанкціонованих змін в інформації. при виявленні\newline Тип: list\newline Значення: \textbf{login, logout, failed\_\allowbreak{}attempt} & SI-7 & a. Впровадити механізми захисту від спаму в точках входу та виходу системи, щоб виявляти та протидіяти небажаним повідомленням.\\ 
b. Оновлювати механізми захисту від спаму, коли доступні нові механізми відповідно до організаційної політики та процедур управління конфігурацією. \\ \hline
172 & УПРАВЛІННЯ ТА ЗБЕРЕЖЕННЯ ІНФОРМАЦІЇ (SI-12) & Параметр: Здійснюється управління інформацією в системі відповідно до чинних законів, наказів, директив, положень, політик, стандартів, інструкцій та експлуатаційних вимог\newline Тип: list\newline Значення: \textbf{default\_\allowbreak{}deny\_\allowbreak{}rule, abac\_\allowbreak{}rule\_\allowbreak{}1}\newline\newline Параметр: Зберігається інформація в системі відповідно до чинних законів, указів Президента, директив, положень, політик, стандартів, інструкцій та експлуатаційних вимог\newline Тип: list\newline Значення: \textbf{default\_\allowbreak{}deny\_\allowbreak{}rule, abac\_\allowbreak{}rule\_\allowbreak{}1}\newline\newline Параметр: Управління інформацією, що виводиться з системи, здійснюється відповідно до чинних законів, указів Президента, директив, положень, політик, стандартів, інструкцій та операційних вимог\newline Тип: list\newline Значення: \textbf{default\_\allowbreak{}deny\_\allowbreak{}rule, abac\_\allowbreak{}rule\_\allowbreak{}1}\newline\newline Параметр: Зберігається інформація, що виводиться з системи, відповідно до чинних законів, наказів, директив, положень, політик, стандартів, інструкцій та експлуатаційних вимог\newline Тип: list\newline Значення: \textbf{default\_\allowbreak{}deny\_\allowbreak{}rule, abac\_\allowbreak{}rule\_\allowbreak{}1} & SI-12 & Управляти та зберігати інформацію всередині системи та виводити інформацію із системи відповідно до чинного законодавства, виконавчих наказів, директив, правил, політик, стандартів, керівних принципів та експлуатаційних вимог. \\ \hline
\multicolumn{5}{|l|}{\textbf{18. УПРАВЛІННЯ РИЗИКАМИ В ЛАНЦЮГУ ПОСТАЧАННЯ (SR) (SR)}} \\ \hline
173 & ПОЛІТИКА ТА ПРОЦЕДУРИ УПРАВЛІННЯ РИЗИКАМИ ЛАНЦЮГА ПОСТАЧАННЯ (SR-1) & Параметр: Визначено персонал або ролі, на які поширюється політика управління ризиками ланцюга постачання\newline Тип: list\newline Значення: \textbf{admin, security\_\allowbreak{}officer}\newline\newline Параметр: Визначено персонал або ролі, на які поширюються процедури управління ризиками ланцюга постачання\newline Тип: list\newline Значення: \textbf{admin, security\_\allowbreak{}officer}\newline\newline Параметр: Визначена посадова особа, відповідальна за розробку, документування та розповсюдження політики та процедур управління ризиками ланцюга постачання\newline Тип: list\newline Значення: \textbf{admin, security\_\allowbreak{}officer}\newline\newline Параметр: Визначена періодичність перегляду та оновлення поточної політики управління ризиками ланцюга постачання\newline Тип: list\newline Значення: \textbf{default\_\allowbreak{}deny\_\allowbreak{}rule, abac\_\allowbreak{}rule\_\allowbreak{}1}\newline\newline Параметр: Є події, які вимагають перегляду та оновлення поточної політики управління ризиками ланцюга постачання\newline Тип: list\newline Значення: \textbf{default\_\allowbreak{}deny\_\allowbreak{}rule, abac\_\allowbreak{}rule\_\allowbreak{}1}\newline\newline Параметр: Визначена періодичність перегляду та оновлення поточної процедури управління ризиками ланцюга постачання\newline Тип: string\newline Значення: \textbf{щорічно} & SR-1 & a. Розробіть, задокументуйте та поширте [Призначення: персонал або ролі, визначені організацією]:\\ 
1. [Вибір (один або декілька): Рівень організації; Рівень місії/бізнес-процесу; рівень системи] політика управління ризиками ланцюга постачання, яка: a) Розглядає мету, сферу діяльності, ролі, відповідальність, зобов'язання керівництва, координацію між організаційними підрозділами та відповідність; b) Відповідає чинним законам, виконавчим наказам, директивам, положенням, політикам, стандартам і вказівкам;\\ 
2. Процедури для сприяння впровадженню політики управління ризиками ланцюга постачання та відповідних засобів контролю управління ризиками ланцюга постачання;\\ 
b. Призначити [Призначення: посадова особа, визначена організацією] для управління розробкою, документуванням і розповсюдженням політики та процедур управління ризиками ланцюга постачання;\\ 
c. Перегляньте та оновіть поточне управління ризиками ланцюга постачання:\\ 
1. Політика [Призначення: частота, визначена організацією] та наступне [Призначення: події, визначені організацією];\\ 
2. Процедури [Призначення: частота, визначена організацією] та наступні [Призначення: події, визначені організацією]. \\ \hline
174 & ПЛАН УПРАВЛІННЯ РИЗИКАМИ ЛАНЦЮГА ПОСТАЧАННЯ (SR-2) &  & SR-2 & a. Розробіть план управління ризиками ланцюга постачання, пов'язаними з дослідженнями та розробкою, проектуванням, виробництвом, придбанням, доставкою, інтеграцією, експлуатацією та обслуговуванням, а також утилізацією таких систем, компонентів системи або послуг для системи: [Призначення: системи, визначені організацією, системні компоненти або системні служби];\\ 
b. Перегляньте та оновіть план управління ризиками ланцюга постачання [Призначення: частота, визначена організацією] або за потреби для усунення загроз;\\ 
c. Захистіть план управління ризиками ланцюга постачання від несанкціонованого розголошення та модифікації. \\ \hline
175 & КОНТРОЛЬ ЛАНЦЮГА ПОСТАЧАННЯ І ПРОЦЕСІВ (SR-3) & Параметр: Визначено персонал ланцюга поставок, з яким необхідно координувати процес або процеси виявлення та усунення слабких місць або недоліків в елементах і процесах ланцюга постачання\newline Тип: list\newline Значення: \textbf{admin, security\_\allowbreak{}officer}\newline\newline Параметр: Визначені засоби контролю ланцюга постачання, що застосовуються для захисту від ризиків ланцюга постачання для системи, системного компонента або системної послуги, а також для обмеження шкоди або наслідків від подій, пов'язаних з ланцюгом постачання\newline Тип: string\newline Значення: \textbf{автоматизований засіб моніторингу}\newline\newline Параметр: Вибрано одне або декілька з наступних ЗНАЧЕНЬ ПАРАМЕТРІВ: \{плани безпеки та конфіденційності; план управління ризиками ланцюга постачання; документ\}; SR-03\_\allowbreak{}ODP[05] визначено документ, що ідентифікує обрані та впроваджені процеси та засоби контролю ланцюга постачання (якщо обрано); SR-03a.[01] запроваджено процес або процеси для виявлення та усунення слабких місць або недоліків в елементах та процесах ланцюга постачання; SR-03a.[02] процес або процеси виявлення та усунення слабких місць або недоліків в елементах та процесах ланцюга постачання системи або компонента системи координується/координуються з персоналом ланцюга постачання; SR-03b. застосовуються засоби контролю ланцюга постачання для захисту від ризиків ланцюга постачання для системи, системного компонента або системної послуги, а також для обмеження шкоди або наслідків від подій, пов'язаних з ланцюгом постачання; SR-03c. задокументовані обрані та впроваджені процеси та засоби контролю ланцюга постачання в ЗНАЧЕННЯ ВИБІРКОВОГО ПАРАМЕТРА(ів)\newline Тип: list\newline Значення: \textbf{admin, security\_\allowbreak{}officer}\newline\newline Параметр: Визначено документ, що ідентифікує обрані та впроваджені процеси та засоби контролю ланцюга постачання (якщо обрано); SR-03a.[01] запроваджено процес або процеси для виявлення та усунення слабких місць або недоліків в елементах та процесах ланцюга постачання; SR-03a.[02] процес або процеси виявлення та усунення слабких місць або недоліків в елементах та процесах ланцюга постачання системи або компонента системи координується/координуються з персоналом ланцюга постачання; SR-03b. застосовуються засоби контролю ланцюга постачання для захисту від ризиків ланцюга постачання для системи, системного компонента або системної послуги, а також для обмеження шкоди або наслідків від подій, пов'язаних з ланцюгом постачання; SR-03c. задокументовані обрані та впроваджені процеси та засоби контролю ланцюга постачання в ЗНАЧЕННЯ ВИБІРКОВОГО ПАРАМЕТРА(ів)\newline Тип: list\newline Значення: \textbf{admin, security\_\allowbreak{}officer} & SR-3 & a. Встановлення процесу або процесів для виявлення та усунення слабких місць або недоліків в елементах і процесах ланцюга постачання [Призначення: визначена організацією система або компонент системи] у координації з [Завдання: персонал ланцюга постачання, визначений організацією];\\ 
b. Використовуйте такі заходи захисту, щоб захистити систему, компонент системи або системну службу від ризиків ланцюга постачання та обмежити шкоду чи наслідки від подій, пов'язаних із ланцюгом постачання: [Призначення: заходи захисту ланцюга постачання, визначені організацією];\\ 
c. Задокументуйте обрані та впроваджені процеси та заходи захисту ланцюгом постачання у [Вибір: плани безпеки та приватності; план управління ризиками ланцюга постачання; [Призначення: документ, визначений організацією]]. \\ \hline
176 & СТРАТЕГІЇ ПРИДБАННЯ, ІНСТРУМЕНТИ І МЕТОДИ (SR-5) &  & SR-5 & Використовуйте наступні стратегії придбання, контрактні інструменти та методи закупівель, щоб захистити від ризиків ланцюга постачання, визначити та пом'якшити їх: [Призначення: визначені організацією стратегії придбання, контрактні інструменти та методи закупівель]. \\ \hline
177 & ОЦІНКА ПОСТАЧАЛЬНИКІВ (SR-6) & Параметр: Оцінюються та аналізуються ризики, пов'язані з ланцюгом постачання, які стосуються постачальників або підрядників та систем, компонентів системи або системних послуг, які вони надають < SR-06\_\allowbreak{}ODP частота >\newline Тип: string\newline Значення: \textbf{щорічно}\newline\newline Параметр: Визначена періодичність оцінки та аналізу ризиків, пов'язаних з ланцюгом постачання, що стосуються постачальників або підрядників, а також систем, компонентів системи або системних послуг, які вони надають\newline Тип: string\newline Значення: \textbf{щорічно} & SR-6 & Оцініть і перегляньте ризики ланцюга постачання, пов'язані з постачальниками або підрядниками, системою, системним компонентом або системною послугою, яку вони надають [Призначення: частота, визначена організацією]. \\ \hline


\end{longtable}

\end{document}
